% Complete reviewed Unit 13; editorial master: chapters/ch13/student.md
% Markdown SHA-256: 23b35a247e66b72cf78bfe01c9ce7882c58af054ecffc4c42d5187efa43cb635
% One-use initial assembly followed by complete content and page review.
\newcommand{\runningtitle}{Unit 13. Clause order and reference}
% Layout inherited from the reviewed previous course; set a running title per unit.
% This is a small hand-maintained TeX layout, not a conversion pipeline.
\RequirePackage{fix-cm}
\documentclass[11pt,twoside]{article}
\usepackage{fontspec}
\setmainfont{Cambria}[
  Renderer=HarfBuzz,
  BoldFont=cambriab.ttf,
  ItalicFont=cambriai.ttf,
  BoldItalicFont=cambriaz.ttf
]
\setmonofont{lmmono10-regular.otf}[BoldFont=lmmonolt10-bold.otf]
\usepackage[paperwidth=165mm,paperheight=241mm,inner=21mm,outer=19mm,
  top=22mm,bottom=20mm,headheight=14pt,headsep=7mm,footskip=10mm]{geometry}
\usepackage{xcolor,array,booktabs,longtable,tabularx}
\usepackage{enumitem,fancyhdr}
\usepackage[unicode,hidelinks]{hyperref}
\definecolor{ink}{HTML}{000000}
\definecolor{muted}{HTML}{4C5558}
\setlength{\parindent}{0pt}
\setlength{\parskip}{3.5pt}
\AtBeginDocument{\fontsize{11}{14.3}\selectfont}
\setlength{\tabcolsep}{4pt}
\renewcommand{\arraystretch}{1.10}
\setlist[itemize]{leftmargin=1.3em,label=\textbullet,itemsep=2pt,topsep=3pt,parsep=0pt}
\setlist[enumerate]{leftmargin=1.8em,itemsep=3pt,topsep=4pt,parsep=0pt}
\widowpenalty=10000
\clubpenalty=10000
\emergencystretch=1em
\raggedbottom
\pagestyle{fancy}
\fancyhf{}
\providecommand{\runningtitle}{Old English}
\fancyhead[LE,RO]{\small\itshape \runningtitle}
\fancyhead[LO,RE]{\small Old English}
\fancyfoot[LE,RO]{\small\thepage}
\renewcommand{\headrulewidth}{0.3pt}
\newcommand{\booktitle}[2]{
  {\fontsize{9}{11}\selectfont\color{muted}#1\par}
  \vspace{3pt}
  {\fontsize{20}{24}\selectfont\color{ink}#2\par}
  \vspace{6pt}
}
\newcommand{\source}[1]{\par{\footnotesize\color{muted}#1\par}}
\newcommand{\readingtext}[1]{
  \par\addvspace{5pt}
  \begingroup\leftskip=1.5em
    \fontsize{12}{16}\selectfont #1\par
  \endgroup\addvspace{5pt}
}
\makeatletter
\renewcommand{\section}{\@startsection{section}{1}{\z@}
  {-2.4ex \@plus -0.6ex \@minus -0.2ex}{0.8ex \@plus 0.2ex}
  {\normalfont\fontsize{12}{15}\selectfont\bfseries\color{ink}}}
\renewcommand{\subsection}{\@startsection{subsection}{2}{\z@}
  {-1.8ex \@plus -0.5ex \@minus -0.2ex}{0.5ex \@plus 0.2ex}
  {\normalfont\fontsize{11}{14}\selectfont\bfseries\color{ink}}}
\makeatother

\hyphenpenalty=10000
\exhyphenpenalty=10000
\emergencystretch=3em
\newcommand{\itempara}[2]{\par\addvspace{3pt}\noindent\hangindent=1.8em\hangafter=1\makebox[1.8em][l]{#1.}#2\par}
\newcommand{\reserve}[1]{\par\penalty-100\begingroup\dimen0=\pagegoal\advance\dimen0 by-\pagetotal\ifdim\dimen0<#1\newpage\fi\endgroup}
\begin{document}
\addfontfeatures{Ligatures=TeXOff}
\raggedright
\booktitle{UNIT 13}{Clause order and reference}

% mdparagraph:P001:begin
A sentence may begin with its subject, an object, an adverbial, or a dependent clause. Its opening words do not by themselves identify the participants. This unit brings clause structure, agreement and reference together, then applies them to a connected narrative.
% mdparagraph:P001:end

% mdparagraph:P002:begin
Sections 6–8 supply the vocabulary, forms and constructions required by the practice. The four Cook sets are retained in full: 15, 13, 12 and 12 items.
% mdparagraph:P002:end

% mdparagraph:P003:begin
Examples 1–12, Exercises F, A, C, R and S, and The bishop and the monk are new teaching material. The final reading preserves the selected Old English of Cook’s The Passing of Chad; its English orientation and local helps are separate additions. The original notes remain at the end, with their exact attachment points.
% mdparagraph:P003:end

% mdparagraph:P004:begin
\textbf{Scope.} “Production” means use in the senses and constructions taught here or earlier. “Reference” forms are for recognition; their display does not add a new meaning or construction. “Supplied help” stays visible while a task is done and is not a test of unaided recall. In tables, 1, 2 and 3 mean first, second and third person; sg and pl mean singular and plural. Nominative, genitive, dative and accusative are abbreviated N, G, D and A only in diagnostic headings. Preserve the printed distinction between e and ę and the marked vowel lengths.
% mdparagraph:P004:end

\reserve{7\baselineskip}
% mdparagraph:P005:begin
\section*{1. Clause roles and the first constituent}
% mdparagraph:P005:end

% teaching: NEW-U13-GRAM-U13-CLAUSE-ROLE-PARSING
% teaching: NEW-U13-GRAM-U13-MAIN-FRONTING-ORDER
% mdparagraph:P006:begin
Find the finite verb and its subject before deciding how to translate the order. The nominative marks the subject; an accusative may be a direct object, but its role must also fit the verb. Agreement and the verb’s construction help when endings coincide. A constituent is a whole phrase functioning as a unit: se munuc and ðone sang each contain two words but each is one constituent.
% mdparagraph:P006:end

\reserve{8\baselineskip}
% example: NEW-U13-EXAMPLE01
\reserve{5\baselineskip}
% mdparagraph:P007:begin
\textbf{Example 1.}
% mdparagraph:P007:end

\reserve{5\baselineskip}
% mdparagraph:P008:begin
\textbf{Se munuc hīerde ðone sang.}
% mdparagraph:P008:end

% mdparagraph:P009:begin
The monk heard the song.
% mdparagraph:P009:end

% mdparagraph:P010:begin
Se munuc is nominative; ðone sang is accusative. The finite verb hīerde agrees with the singular subject.
% mdparagraph:P010:end

\reserve{8\baselineskip}
% example: NEW-U13-EXAMPLE02
\reserve{5\baselineskip}
% mdparagraph:P011:begin
\textbf{Example 2.}
% mdparagraph:P011:end

\reserve{5\baselineskip}
% mdparagraph:P012:begin
\textbf{Ðone sang hīerde se munuc.}
% mdparagraph:P012:end

% mdparagraph:P013:begin
The monk heard the song.
% mdparagraph:P013:end

% mdparagraph:P014:begin
The object phrase comes first. Its accusative ending still makes it the object. The finite verb follows that complete phrase, not the first written word.
% mdparagraph:P014:end

\reserve{8\baselineskip}
% example: NEW-U13-EXAMPLE03
\reserve{5\baselineskip}
% mdparagraph:P015:begin
\textbf{Example 3.}
% mdparagraph:P015:end

\reserve{5\baselineskip}
% mdparagraph:P016:begin
\textbf{Ðā hīerde se munuc ðone sang.}
% mdparagraph:P016:end

% mdparagraph:P017:begin
Then the monk heard the song.
% mdparagraph:P017:end

% mdparagraph:P018:begin
The initial adverb ðā means “then.” Here the finite verb precedes the noun subject. This is a common main-clause order after an initial constituent.
% mdparagraph:P018:end

% mdparagraph:P019:begin
Finite-second order is a tendency, not a requirement. It concerns constituents, not word counts. A coordinating and, ac or oððe does not itself count as the first constituent of the following clause. Fronting gives a constituent a prominent position, but the exact effect depends on context; no change of case follows merely from a change of order.
% mdparagraph:P019:end

\reserve{7\baselineskip}
% mdparagraph:P020:begin
\section*{2. Dependent clauses and limits on order rules}
% mdparagraph:P020:end

% teaching: NEW-U13-DEPENDENT-ORDER-HELP
% teaching: NEW-U13-GRAM-U13-DEPENDENT-ORDER-TENDENCIES
% teaching: NEW-U13-GRAM-U13-CORRELATIVE-CLAUSE-REFERENCE
% teaching: NEW-U13-GRAM-U13-PRONOUN-POSITION-LIMITS
% mdparagraph:P021:begin
A dependent clause has its own finite verb and subject. Its finite verb often comes late, frequently last. The conjunction or relative marker begins the clause; it is not its subject unless the construction makes it so. Identify the boundary of each clause before comparing orders.
% mdparagraph:P021:end

\reserve{8\baselineskip}
% example: NEW-U13-EXAMPLE04
\reserve{5\baselineskip}
% mdparagraph:P022:begin
\textbf{Example 4.}
% mdparagraph:P022:end

\reserve{5\baselineskip}
% mdparagraph:P023:begin
\textbf{Ðā hē ðone sang hīerde, ðā ēode hē tō ðǣm bisceope.}
% mdparagraph:P023:end

% mdparagraph:P024:begin
When he heard the song, he went to the bishop.
% mdparagraph:P024:end

% mdparagraph:P025:begin
The first ðā means “when” and introduces the dependent clause ending in hīerde. The second means “then” and belongs to the main clause. Both occurrences of hē refer to the person already in view; the order alone does not supply his identity.
% mdparagraph:P025:end

\reserve{8\baselineskip}
% example: NEW-U13-EXAMPLE09
\reserve{5\baselineskip}
% mdparagraph:P026:begin
\textbf{Example 9.}
% mdparagraph:P026:end

\reserve{5\baselineskip}
% mdparagraph:P027:begin
\textbf{Se munuc, ðe se bisceop grētte, ēode tō ðǣre cirican.}
% mdparagraph:P027:end

% mdparagraph:P028:begin
The monk whom the bishop greeted went to the church.
% mdparagraph:P028:end

% mdparagraph:P029:begin
Invariant ðe introduces the relative clause and refers back to se munuc. The monk is the subject of ēode, but is the person greeted inside the relative; se bisceop is the relative clause’s subject. Keep the two roles separate.
% mdparagraph:P029:end

\reserve{8\baselineskip}
% example: NEW-U13-EXAMPLE10
\reserve{5\baselineskip}
% mdparagraph:P030:begin
\textbf{Example 10.}
% mdparagraph:P030:end

\reserve{5\baselineskip}
% mdparagraph:P031:begin
\textbf{Ðā se bisceop cōm, ðā sungon ðā brōðor.}
% mdparagraph:P031:end

% mdparagraph:P032:begin
When the bishop came, the brothers sang.
% mdparagraph:P032:end

% mdparagraph:P033:begin
The first ðā introduces “when the bishop came”; the second is the main-clause adverb “then.” The last ðā is the plural demonstrative with brōðor. The plural subject agrees with sungon.
% mdparagraph:P033:end

% mdparagraph:P034:begin
After an initial adverbial a subject pronoun may precede the finite verb. The following short source fragment gives an actual instance. It belongs to the continuation of the Chad passage, outside the continuous extract assigned below.
% mdparagraph:P034:end

% source-example: NEW-U13-SOURCE-ORDER-FRAGMENT
\reserve{5\baselineskip}
% mdparagraph:P035:begin
\textbf{Æfter ðon hē underðēodde}
% mdparagraph:P035:end

% mdparagraph:P036:begin
After that he added
% mdparagraph:P036:end

% teaching: NEW-U13-SOURCE-ORDER-HELP
% support: NEW-U13-SOURCE-ORDER-HELP
% mdparagraph:P037:begin
\textbf{Order help.} This complete fragment is supplied as corroborative evidence outside the assigned continuous reading. Æfter ðon means after that; hē is Chad; underðēodde means added to his speech (Cook note13: Latin subjunxit). Identify its order only; no independent translation or lexical recall is required. The fronted temporal phrase is followed by subject pronoun, then finite verb.
% mdparagraph:P037:end

% teaching: NEW-U13-VOC-SE
% teaching: NEW-U13-VOC-AEFTER
% teaching: NEW-U13-VOC-HE
% teaching: NEW-U13-VOC-UNDERTHEODAN
% teaching: NEW-U13-PAR-SE
% teaching: NEW-U13-PAR-AEFTER
% teaching: NEW-U13-PAR-HE
% teaching: NEW-U13-PAR-UNDERTHEODAN
% teaching: NEW-U13-GRAM-u13-source-after-that
% mdparagraph:P038:begin
The initial phrase Æfter ðon is followed by hē and then underðēodde. Thus even counting constituents does not turn main-clause finite-second order into a universal rule.
% mdparagraph:P038:end

% teaching: NEW-U13-SOURCE-DEPENDENT-ORDER
% mdparagraph:P039:begin
The assigned reading also contains a dependent clause whose finite verb is followed by another constituent:
% mdparagraph:P039:end

% source-example: NEW-U13-SOURCE-DEPENDENT-ORDER
\reserve{5\baselineskip}
% mdparagraph:P040:begin
\textbf{ðæt hē ǣrest gehīerde fram ēastsūðdǣle heofones}
% mdparagraph:P040:end

% support: NEW-U13-DEPENDENT-ORDER-HELP
% mdparagraph:P041:begin
\textbf{Order help.} This is a fragment of the assigned Chad reading: “that he first heard [the sound and song] from the south-eastern part of the sky.” Hē is Owini; gehīerde is the finite preterite “heard”; fram ēastsūðdǣle heofones names the origin of the sound and follows that verb inside the content clause. The objects have been fronted earlier in the full sentence, as Cook note 152.8 explains; they are supplied in brackets only in the English help. Identify the order with this help; no independent recall of the supplied forms is required.
% mdparagraph:P041:end

% teaching: NEW-U13-VOC-EASTSUTHDAEL
% teaching: NEW-U13-VOC-AEREST
% teaching: NEW-U13-SOURCE-HELPS
% teaching: NEW-U13-PAR-AEREST
% mdparagraph:P042:begin
The phrase beginning fram follows gehīerde inside the ðæt clause. “Often late” therefore does not mean “always last.” Cook 25.7 nevertheless asks for a final verb in its specified relative clause. A stated exercise target can select one order without declaring every other order ungrammatical.
% mdparagraph:P042:end

\reserve{7\baselineskip}
% mdparagraph:P043:begin
\section*{3. Pronouns and sustained reference}
% mdparagraph:P043:end

% teaching: NEW-U13-GRAM-U13-REFERENCE-RESOLUTION-AMBIGUITY
% mdparagraph:P044:begin
A pronoun agrees with the grammatical gender and number of its antecedent. Its case expresses its role in its own clause. Masculine hē or hine can therefore refer to a masculine thing; English may require “it.” Agreement restricts the possible antecedents but need not select one uniquely.
% mdparagraph:P044:end

\reserve{8\baselineskip}
% example: NEW-U13-EXAMPLE05
\reserve{5\baselineskip}
% mdparagraph:P045:begin
\textbf{Example 5.}
% mdparagraph:P045:end

\reserve{5\baselineskip}
% mdparagraph:P046:begin
\textbf{Se sang wæs swēte. Hē wæs ēac wynsum.}
% mdparagraph:P046:end

% mdparagraph:P047:begin
The song was sweet. It was also pleasant.
% mdparagraph:P047:end

% mdparagraph:P048:begin
Sang is masculine. Hē refers to the song, so English uses “it.” The pronoun’s nominative case belongs to its own clause.
% mdparagraph:P048:end

\reserve{8\baselineskip}
% example: NEW-U13-EXAMPLE06
\reserve{5\baselineskip}
% mdparagraph:P049:begin
\textbf{Example 6.}
% mdparagraph:P049:end

\reserve{5\baselineskip}
% mdparagraph:P050:begin
\textbf{Se bisceop geseah ðone munuc and grētte hine.}
% mdparagraph:P050:end

% mdparagraph:P051:begin
The bishop saw the monk and greeted him.
% mdparagraph:P051:end

% mdparagraph:P052:begin
The second finite verb shares the bishop as its subject. Hine is accusative and refers to the monk, who receives the greeting. The preceding direct object does not become the next subject automatically.
% mdparagraph:P052:end

\reserve{8\baselineskip}
% example: NEW-U13-EXAMPLE07
\reserve{5\baselineskip}
% mdparagraph:P053:begin
\textbf{Example 7.}
% mdparagraph:P053:end

\reserve{5\baselineskip}
% mdparagraph:P054:begin
\textbf{Se bisceop geseah ðone munuc. Hē wæs on ðǣm hūse.}
% mdparagraph:P054:end

% mdparagraph:P055:begin
The bishop saw the monk. He was in the house.
% mdparagraph:P055:end

% mdparagraph:P056:begin
Both men are masculine singular. Hē can refer to either without a fuller context. Neither the nearest noun nor the previous subject is an infallible rule. Repeating the intended noun phrase removes the ambiguity.
% mdparagraph:P056:end

% mdparagraph:P057:begin
In a longer passage, keep a separate record for each participant and for any thing later named by a pronoun. Recheck the reference when a new subject, apposition or quotation appears. A statement such as hē, se bisceop makes the reference explicit at that point; it need not give every earlier hē the same referent.
% mdparagraph:P057:end

\reserve{7\baselineskip}
% mdparagraph:P058:begin
\section*{4. Coordination and omitted material}
% mdparagraph:P058:end

% teaching: NEW-U13-GRAM-U13-COORDINATION-PREDICATE-RECOVERY
% mdparagraph:P059:begin
Two coordinated predicates can share an expressed subject, as in Example 6. Two coordinated clauses may instead have different subjects and omit a repeated predicate. Recover only what the syntax and context support.
% mdparagraph:P059:end

\reserve{8\baselineskip}
% example: NEW-U13-EXAMPLE11
\reserve{5\baselineskip}
% mdparagraph:P060:begin
\textbf{Example 11.}
% mdparagraph:P060:end

\reserve{5\baselineskip}
% mdparagraph:P061:begin
\textbf{Se bisceop wæs on ðǣm hūse, and se munuc on ðǣre cirican.}
% mdparagraph:P061:end

% mdparagraph:P062:begin
The bishop was in the house, and the monk in the church.
% mdparagraph:P062:end

% mdparagraph:P063:begin
The second clause has its own subject, se munuc, and omits the repeated copula. Both subjects here are singular. If a restored verb belongs to a plural subject, it must have plural agreement; recovering the lexical predicate does not license copying a singular ending.
% mdparagraph:P063:end

% mdparagraph:P064:begin
Coordinate clauses remain separate for agreement and reference. The same caution applies when one clause contains several verbs: a shared subject is recovered from the construction, not from proximity to the last noun.
% mdparagraph:P064:end

\reserve{7\baselineskip}
% mdparagraph:P065:begin
\section*{5. Speech and long sentences}
% mdparagraph:P065:end

% teaching: NEW-U13-GRAM-U13-QUOTED-REFERENCE
% teaching: NEW-U13-GRAM-U13-LONG-SENTENCE-ANALYSIS
% teaching: NEW-U13-GRAM-U13-NEGATION-QUESTION-ORDER
% teaching: NEW-U13-GRAM-U13-SUPPORTED-LONG-FRONTING
\reserve{8\baselineskip}
% example: NEW-U13-EXAMPLE08
\reserve{5\baselineskip}
% mdparagraph:P066:begin
\textbf{Example 8.}
% mdparagraph:P066:end

\reserve{5\baselineskip}
% mdparagraph:P067:begin
\textbf{Se bisceop cwæð: “Ic wille mid ēow sprecan.”}
% mdparagraph:P067:end

% mdparagraph:P068:begin
The bishop said, “I want to speak with you.”
% mdparagraph:P068:end

% mdparagraph:P069:begin
Within the quotation, Ic is the bishop and ēow addresses several people. Sprecan with mid and the dative means “speak with”; the plain infinitive follows wille.
% mdparagraph:P069:end

\reserve{8\baselineskip}
% example: NEW-U13-EXAMPLE12
\reserve{5\baselineskip}
% mdparagraph:P070:begin
\textbf{Example 12.}
% mdparagraph:P070:end

\reserve{5\baselineskip}
% mdparagraph:P071:begin
\textbf{Hē sægde ðæt se sang swēte wǣre.}
% mdparagraph:P071:end

% mdparagraph:P072:begin
He said that the song was sweet.
% mdparagraph:P072:end

% mdparagraph:P073:begin
Ðæt introduces reported content. Se sang is its subject and wǣre its finite verb. The report subjunctive does not decide who an otherwise ambiguous pronoun refers to.
% mdparagraph:P073:end

% mdparagraph:P074:begin
Direct speech has its own first and second persons. In an indirect report, a third-person pronoun must be resolved from the reporting context. Questions and negatives keep their previously taught constructions: a yes/\allowbreak{}no question may begin with the finite verb; an interrogative word has a role within its question; ne belongs with the finite verb, and negative concord remains possible. No new negative-copula contraction is required here.
% mdparagraph:P074:end

% mdparagraph:P075:begin
For a long sentence, first mark each finite verb. Then bracket conjunction and relative clauses, identify each subject, and attach the remaining complements to the verb that governs them. Resolve pronouns only after these relations are clear. Finally restore the temporal sequence and translate in connected English. Coordination, parenthesis and direct speech each need their own boundaries.
% mdparagraph:P075:end

% mdparagraph:P076:begin
One supported construction in the Chad reading has objects placed before a reporting clause even though their verb is inside the following content clause. Cook’s note identifies that dependence. Use the printed help to analyse it; the unit does not ask you to create new sentences by freely moving material out of embedded clauses.
% mdparagraph:P076:end

\reserve{7\baselineskip}
% mdparagraph:P077:begin
\section*{6. Nouns and adjectives}
% mdparagraph:P077:end

% mdparagraph:P078:begin
The following vocabulary gives the ordinary meanings needed in this unit. The genitive connects one noun with another; prepositions and verbs determine their own governed cases. A complete declension supplies forms, not a permission to substitute any case in any construction.
% mdparagraph:P078:end

% teaching: NEW-U13-VOC-IDELNES
% teaching: NEW-U13-NOM-IDELNES
% teaching: NEW-U13-VOC-MODOR
% teaching: NEW-U13-NOM-MODOR
% teaching: NEW-U13-VOC-SYNN
% teaching: NEW-U13-NOM-SYNN
% teaching: NEW-U13-VOC-DOHTOR
% teaching: NEW-U13-NOM-DOHTOR
% teaching: NEW-U13-VOC-BISCEOP
% teaching: NEW-U13-NOM-BISCEOP
% teaching: NEW-U13-VOC-MYNSTER
% teaching: NEW-U13-NOM-MYNSTER
% teaching: NEW-U13-VOC-LEORNUNG
% teaching: NEW-U13-NOM-LEORNUNG
% teaching: NEW-U13-VOC-CIRICE
% teaching: NEW-U13-NOM-CIRICE
% teaching: NEW-U13-VOC-UNRIHTWISNES
% teaching: NEW-U13-NOM-UNRIHTWISNES
% teaching: NEW-U13-VOC-SANG
% teaching: NEW-U13-NOM-SANG
% teaching: NEW-U13-VOC-LICHAMA
% teaching: NEW-U13-NOM-LICHAMA
% teaching: NEW-U13-VOC-EAGTHYREL
% teaching: NEW-U13-NOM-EAGTHYREL
% teaching: NEW-U13-VOC-DAEL
% teaching: NEW-U13-NOM-DAEL
% teaching: NEW-U13-VOC-ROD
% teaching: NEW-U13-NOM-ROD
% teaching: NEW-U13-VOC-AEX
% teaching: NEW-U13-NOM-AEX
% teaching: NEW-U13-VOC-ADESA
% teaching: NEW-U13-NOM-ADESA
% teaching: NEW-U13-VOC-WYRTTRUMA
% teaching: NEW-U13-NOM-WYRTTRUMA
% teaching: NEW-U13-VOC-UNTRYMNES
% teaching: NEW-U13-NOM-UNTRYMNES
% teaching: NEW-U13-VOC-GEFERSCIPE
% teaching: NEW-U13-NOM-GEFERSCIPE
% teaching: NEW-U13-VOC-GELEAFA
% teaching: NEW-U13-NOM-GELEAFA
% teaching: NEW-U13-VOC-WUNDOR
% teaching: NEW-U13-NOM-WUNDOR
% teaching: NEW-U13-VOC-HELL
% teaching: NEW-U13-NOM-HELL
% teaching: NEW-U13-VOC-HROF
% teaching: NEW-U13-NOM-HROF
% teaching: NEW-U13-VOC-MED
% teaching: NEW-U13-NOM-MED
% teaching: NEW-U13-VOC-AECER
% teaching: NEW-U13-NOM-AECER
% teaching: NEW-U13-VOC-OFEN
% teaching: NEW-U13-NOM-OFEN
% teaching: NEW-U13-VOC-BROTHOR
% teaching: NEW-U13-NOM-BROTHOR
% teaching: NEW-U13-VOC-FORTHFOR
% teaching: NEW-U13-NOM-FORTHFOR
% teaching: NEW-U13-VOC-WEG
% teaching: NEW-U13-NOM-WEG
% teaching: NEW-U13-VOC-FREOND
% teaching: NEW-U13-NOM-FREOND
% teaching: NEW-U13-VOC-BURG
% teaching: NEW-U13-NOM-BURG
% teaching: NEW-U13-VOC-WORD
% teaching: NEW-U13-NOM-WORD
% mdtable:T01:begin
\par\begingroup
\fontsize{10.5}{13.2}\selectfont
\setlength{\LTleft}{0pt}\setlength{\LTright}{\fill}
\begin{longtable}{@{}>{\raggedright\arraybackslash}p{38.39mm}>{\raggedright\arraybackslash}p{31.99mm}>{\raggedright\arraybackslash}p{47.99mm}@{}}
\toprule
\textbf{Noun} & \textbf{Gender} & \textbf{Meaning} \\
\midrule
\endfirsthead
\toprule
\textbf{Noun} & \textbf{Gender} & \textbf{Meaning} \\
\midrule
\endhead
\midrule
\multicolumn{3}{r}{\footnotesize Continued on the following page.}\\
\endfoot
\bottomrule
\endlastfoot
īdelnes & feminine & idleness \\
mōdor & feminine & mother \\
synn & feminine & sin \\
dohtor & feminine & daughter \\
bisceop & masculine & bishop \\
mynster & neuter & monastery \\
leornung & feminine & study; learning \\
cirice & feminine & church \\
unrihtwīsnes & feminine & unrighteousness \\
sang & masculine & song \\
līchama & masculine & body \\
ēagðȳrel & neuter & window \\
dǣl & masculine & part \\
rōd & feminine & cross \\
æx & feminine & ax \\
adesa & masculine & hatchet \\
wyrttruma & masculine & root \\
untrymnes & feminine & infirmity \\
gefērscipe & masculine & companionship \\
gelēafa & masculine & faith \\
wundor & neuter & wonder \\
hęll & feminine & hell \\
hrōf & masculine & roof \\
mēd & feminine & reward \\
æcer & masculine & field \\
ofen & masculine & oven \\
brōðor & masculine & brother \\
forðfōr & feminine & departure \\
weg & masculine & way; path \\
frēond & masculine & friend \\
burg & feminine & town; city \\
word & neuter & word \\
\end{longtable}
\par\endgroup

% mdtable:T01:end
% mdparagraph:P079:begin
The plural forms of īdelnes, leornung, unrihtwīsnes, untrymnes, gefērscipe, gelēafa and hęll are reference forms only. Their assigned abstract senses are singular. All other noun systems below give productive N, G, D and A forms in both numbers, within the meanings just listed. The source spelling ciricean and its sense “oratory” are supplied separately with the reading.
% mdparagraph:P079:end

\reserve{7\baselineskip}
% mdparagraph:P080:begin
\subsection*{6.1. Complete noun models}
% mdparagraph:P080:end

% teaching: NEW-U13-GRAM-U13-NOM-MASC
% teaching: NEW-U13-GRAM-U13-NOM-MASC-E
% teaching: NEW-U13-GRAM-U13-NOM-NEUTER
% teaching: NEW-U13-GRAM-U13-NOM-SYNCOPE
% teaching: NEW-U13-GRAM-U13-NOM-FEMININE
% teaching: NEW-U13-GRAM-U13-NOM-NES-UNG
% teaching: NEW-U13-GRAM-U13-NOM-WEAK
% teaching: NEW-U13-GRAM-U13-NOM-KINSHIP
% teaching: NEW-U13-GRAM-U13-NOM-BURG
% teaching: NEW-U13-GRAM-U13-NOM-FREOND
% mdparagraph:P081:begin
Strong masculines of the bisceop type have singular N/\allowbreak{}A with no ending, G -es and D -e; plural N/\allowbreak{}A -as, G -a and D -um. Retain the vowels in bisceop, æcer and ofen. In gefērscipe, remove citation -e before adding another ending; N/\allowbreak{}A and D singular retain -e.
% mdparagraph:P081:end

% mdparagraph:P082:begin
Strong neuters have G singular -es, D singular -e, G plural -a and D plural -um. Their N/\allowbreak{}A plural is lexical: word and wundor have no added ending; ēagðȳrel and mynster select -u with the shortened stem. Before vowel endings, mynster loses its second e, ēagðȳrel its second e, and wundor its o: mynstr-, ēagðȳrl-, wundr-. Keep the full citation form wherever the tables show it.
% mdparagraph:P082:end

% mdparagraph:P083:begin
Long feminine nouns have singular N with no added ending and G/\allowbreak{}D/\allowbreak{}A -e; selected plural N/\allowbreak{}A/\allowbreak{}G -a and D -um. Plural N/\allowbreak{}A -e is also found. Æx retains æ throughout: no masculine æ/\allowbreak{}a alternation applies. Nouns in -nes double s before an ending. In leornung, the selected oblique singular ending is -e; -a is a documented alternative.
% mdparagraph:P083:end

% mdparagraph:P084:begin
Weak masculine nouns have N singular -a, all other singular cases and N/\allowbreak{}A plural -an, G plural -ena and D plural -um. Weak feminine cirice has N singular -e and follows the same selected endings otherwise. Its productive stem is ciric-.
% mdparagraph:P084:end

% mdparagraph:P085:begin
Kinship nouns require their full tables. Mōdor, dohtor and brōðor have unchanged N/\allowbreak{}G/\allowbreak{}A singular, but D singular mēder, dehter and brēðer. Dohtor has short o. Burg has byrig in selected G/\allowbreak{}D singular and N/\allowbreak{}A plural. Frēond follows the earlier fēond pattern, with frēond- and frīend-; its selected G plural is frēonda.
% mdparagraph:P085:end

\reserve{7\baselineskip}
% mdparagraph:P086:begin
\subsection*{Table 1. īdelnes: feminine noun}
% mdparagraph:P086:end

% table: NEW-U13-NOM-IDELNES
% mdtable:T02:begin
\par\begingroup
\fontsize{10.5}{13.2}\selectfont
\setlength{\LTleft}{0pt}\setlength{\LTright}{\fill}
\noindent
\begin{tabular}{@{}>{\raggedright\arraybackslash}p{39.46mm}>{\raggedright\arraybackslash}p{39.46mm}>{\raggedright\arraybackslash}p{39.46mm}@{}}
\toprule
\textbf{Case} & \textbf{Singular} & \textbf{Plural} \\
\midrule
Nominative & īdelnes & īdelnessa \\
Genitive & īdelnesse & īdelnessa \\
Dative & īdelnesse & īdelnessum \\
Accusative & īdelnesse & īdelnessa \\
\bottomrule
\end{tabular}
\par\endgroup

% mdtable:T02:end
% mdparagraph:P087:begin
Plural: reference only.
% mdparagraph:P087:end

\reserve{7\baselineskip}
% mdparagraph:P088:begin
\subsection*{Table 2. mōdor: feminine noun}
% mdparagraph:P088:end

% table: NEW-U13-NOM-MODOR
% mdtable:T03:begin
\par\begingroup
\fontsize{10.5}{13.2}\selectfont
\setlength{\LTleft}{0pt}\setlength{\LTright}{\fill}
\noindent
\begin{tabular}{@{}>{\raggedright\arraybackslash}p{39.46mm}>{\raggedright\arraybackslash}p{39.46mm}>{\raggedright\arraybackslash}p{39.46mm}@{}}
\toprule
\textbf{Case} & \textbf{Singular} & \textbf{Plural} \\
\midrule
Nominative & mōdor & mōdru \\
Genitive & mōdor & mōdra \\
Dative & mēder & mōdrum \\
Accusative & mōdor & mōdru \\
\bottomrule
\end{tabular}
\par\endgroup

% mdtable:T03:end
\reserve{7\baselineskip}
% mdparagraph:P089:begin
\subsection*{Table 3. dohtor: feminine noun}
% mdparagraph:P089:end

% table: NEW-U13-NOM-DOHTOR
% mdtable:T04:begin
\par\begingroup
\fontsize{10.5}{13.2}\selectfont
\setlength{\LTleft}{0pt}\setlength{\LTright}{\fill}
\noindent
\begin{tabular}{@{}>{\raggedright\arraybackslash}p{39.46mm}>{\raggedright\arraybackslash}p{39.46mm}>{\raggedright\arraybackslash}p{39.46mm}@{}}
\toprule
\textbf{Case} & \textbf{Singular} & \textbf{Plural} \\
\midrule
Nominative & dohtor & dohtor \\
Genitive & dohtor & dohtra \\
Dative & dehter & dohtrum \\
Accusative & dohtor & dohtor \\
\bottomrule
\end{tabular}
\par\endgroup

% mdtable:T04:end
\reserve{7\baselineskip}
% mdparagraph:P090:begin
\subsection*{Table 4. bisceop: masculine noun}
% mdparagraph:P090:end

% table: NEW-U13-NOM-BISCEOP
% mdtable:T05:begin
\par\begingroup
\fontsize{10.5}{13.2}\selectfont
\setlength{\LTleft}{0pt}\setlength{\LTright}{\fill}
\noindent
\begin{tabular}{@{}>{\raggedright\arraybackslash}p{39.46mm}>{\raggedright\arraybackslash}p{39.46mm}>{\raggedright\arraybackslash}p{39.46mm}@{}}
\toprule
\textbf{Case} & \textbf{Singular} & \textbf{Plural} \\
\midrule
Nominative & bisceop & bisceopas \\
Genitive & bisceopes & bisceopa \\
Dative & bisceope & bisceopum \\
Accusative & bisceop & bisceopas \\
\bottomrule
\end{tabular}
\par\endgroup

% mdtable:T05:end
\reserve{7\baselineskip}
% mdparagraph:P091:begin
\subsection*{Table 5. leornung: feminine noun}
% mdparagraph:P091:end

% table: NEW-U13-NOM-LEORNUNG
% mdtable:T06:begin
\par\begingroup
\fontsize{10.5}{13.2}\selectfont
\setlength{\LTleft}{0pt}\setlength{\LTright}{\fill}
\noindent
\begin{tabular}{@{}>{\raggedright\arraybackslash}p{39.46mm}>{\raggedright\arraybackslash}p{39.46mm}>{\raggedright\arraybackslash}p{39.46mm}@{}}
\toprule
\textbf{Case} & \textbf{Singular} & \textbf{Plural} \\
\midrule
Nominative & leornung & leornunga \\
Genitive & leornunge & leornunga \\
Dative & leornunge & leornungum \\
Accusative & leornunge & leornunga \\
\bottomrule
\end{tabular}
\par\endgroup

% mdtable:T06:end
% mdparagraph:P092:begin
Plural: reference only.
% mdparagraph:P092:end

\reserve{7\baselineskip}
% mdparagraph:P093:begin
\subsection*{Table 6. cirice: feminine noun}
% mdparagraph:P093:end

% table: NEW-U13-NOM-CIRICE
% mdtable:T07:begin
\par\begingroup
\fontsize{10.5}{13.2}\selectfont
\setlength{\LTleft}{0pt}\setlength{\LTright}{\fill}
\noindent
\begin{tabular}{@{}>{\raggedright\arraybackslash}p{39.46mm}>{\raggedright\arraybackslash}p{39.46mm}>{\raggedright\arraybackslash}p{39.46mm}@{}}
\toprule
\textbf{Case} & \textbf{Singular} & \textbf{Plural} \\
\midrule
Nominative & cirice & cirican \\
Genitive & cirican & ciricena \\
Dative & cirican & ciricum \\
Accusative & cirican & cirican \\
\bottomrule
\end{tabular}
\par\endgroup

% mdtable:T07:end
\reserve{7\baselineskip}
% mdparagraph:P094:begin
\subsection*{Table 7. līchama: masculine noun}
% mdparagraph:P094:end

% table: NEW-U13-NOM-LICHAMA
% mdtable:T08:begin
\par\begingroup
\fontsize{10.5}{13.2}\selectfont
\setlength{\LTleft}{0pt}\setlength{\LTright}{\fill}
\noindent
\begin{tabular}{@{}>{\raggedright\arraybackslash}p{39.46mm}>{\raggedright\arraybackslash}p{39.46mm}>{\raggedright\arraybackslash}p{39.46mm}@{}}
\toprule
\textbf{Case} & \textbf{Singular} & \textbf{Plural} \\
\midrule
Nominative & līchama & līchaman \\
Genitive & līchaman & līchamena \\
Dative & līchaman & līchamum \\
Accusative & līchaman & līchaman \\
\bottomrule
\end{tabular}
\par\endgroup

% mdtable:T08:end
\reserve{7\baselineskip}
% mdparagraph:P095:begin
\subsection*{Table 8. ēagðȳrel: neuter noun}
% mdparagraph:P095:end

% table: NEW-U13-NOM-EAGTHYREL
% mdtable:T09:begin
\par\begingroup
\fontsize{10.5}{13.2}\selectfont
\setlength{\LTleft}{0pt}\setlength{\LTright}{\fill}
\noindent
\begin{tabular}{@{}>{\raggedright\arraybackslash}p{39.46mm}>{\raggedright\arraybackslash}p{39.46mm}>{\raggedright\arraybackslash}p{39.46mm}@{}}
\toprule
\textbf{Case} & \textbf{Singular} & \textbf{Plural} \\
\midrule
Nominative & ēagðȳrel & ēagðȳrlu \\
Genitive & ēagðȳrles & ēagðȳrla \\
Dative & ēagðȳrle & ēagðȳrlum \\
Accusative & ēagðȳrel & ēagðȳrlu \\
\bottomrule
\end{tabular}
\par\endgroup

% mdtable:T09:end
\reserve{7\baselineskip}
% mdparagraph:P096:begin
\subsection*{Table 9. æx: feminine noun}
% mdparagraph:P096:end

% table: NEW-U13-NOM-AEX
% mdtable:T10:begin
\par\begingroup
\fontsize{10.5}{13.2}\selectfont
\setlength{\LTleft}{0pt}\setlength{\LTright}{\fill}
\noindent
\begin{tabular}{@{}>{\raggedright\arraybackslash}p{39.46mm}>{\raggedright\arraybackslash}p{39.46mm}>{\raggedright\arraybackslash}p{39.46mm}@{}}
\toprule
\textbf{Case} & \textbf{Singular} & \textbf{Plural} \\
\midrule
Nominative & æx & æxa \\
Genitive & æxe & æxa \\
Dative & æxe & æxum \\
Accusative & æxe & æxa \\
\bottomrule
\end{tabular}
\par\endgroup

% mdtable:T10:end
\reserve{7\baselineskip}
% mdparagraph:P097:begin
\subsection*{Table 10. gefērscipe: masculine noun}
% mdparagraph:P097:end

% table: NEW-U13-NOM-GEFERSCIPE
% mdtable:T11:begin
\par\begingroup
\fontsize{10.5}{13.2}\selectfont
\setlength{\LTleft}{0pt}\setlength{\LTright}{\fill}
\noindent
\begin{tabular}{@{}>{\raggedright\arraybackslash}p{36.99mm}>{\raggedright\arraybackslash}p{40.69mm}>{\raggedright\arraybackslash}p{40.69mm}@{}}
\toprule
\textbf{Case} & \textbf{Singular} & \textbf{Plural} \\
\midrule
Nominative & gefērscipe & gefērscipas \\
Genitive & gefērscipes & gefērscipa \\
Dative & gefērscipe & gefērscipum \\
Accusative & gefērscipe & gefērscipas \\
\bottomrule
\end{tabular}
\par\endgroup

% mdtable:T11:end
% mdparagraph:P098:begin
Plural: reference only.
% mdparagraph:P098:end

\reserve{7\baselineskip}
% mdparagraph:P099:begin
\subsection*{Table 11. wundor: neuter noun}
% mdparagraph:P099:end

% table: NEW-U13-NOM-WUNDOR
% mdtable:T12:begin
\par\begingroup
\fontsize{10.5}{13.2}\selectfont
\setlength{\LTleft}{0pt}\setlength{\LTright}{\fill}
\noindent
\begin{tabular}{@{}>{\raggedright\arraybackslash}p{39.46mm}>{\raggedright\arraybackslash}p{39.46mm}>{\raggedright\arraybackslash}p{39.46mm}@{}}
\toprule
\textbf{Case} & \textbf{Singular} & \textbf{Plural} \\
\midrule
Nominative & wundor & wundor \\
Genitive & wundres & wundra \\
Dative & wundre & wundrum \\
Accusative & wundor & wundor \\
\bottomrule
\end{tabular}
\par\endgroup

% mdtable:T12:end
\reserve{7\baselineskip}
% mdparagraph:P100:begin
\subsection*{Table 12. brōðor: masculine noun}
% mdparagraph:P100:end

% table: NEW-U13-NOM-BROTHOR
% mdtable:T13:begin
\par\begingroup
\fontsize{10.5}{13.2}\selectfont
\setlength{\LTleft}{0pt}\setlength{\LTright}{\fill}
\noindent
\begin{tabular}{@{}>{\raggedright\arraybackslash}p{39.46mm}>{\raggedright\arraybackslash}p{39.46mm}>{\raggedright\arraybackslash}p{39.46mm}@{}}
\toprule
\textbf{Case} & \textbf{Singular} & \textbf{Plural} \\
\midrule
Nominative & brōðor & brōðor \\
Genitive & brōðor & brōðra \\
Dative & brēðer & brōðrum \\
Accusative & brōðor & brōðor \\
\bottomrule
\end{tabular}
\par\endgroup

% mdtable:T13:end
\reserve{7\baselineskip}
% mdparagraph:P101:begin
\subsection*{Table 13. frēond: masculine noun}
% mdparagraph:P101:end

% table: NEW-U13-NOM-FREOND
% mdtable:T14:begin
\par\begingroup
\fontsize{10.5}{13.2}\selectfont
\setlength{\LTleft}{0pt}\setlength{\LTright}{\fill}
\noindent
\begin{tabular}{@{}>{\raggedright\arraybackslash}p{39.46mm}>{\raggedright\arraybackslash}p{39.46mm}>{\raggedright\arraybackslash}p{39.46mm}@{}}
\toprule
\textbf{Case} & \textbf{Singular} & \textbf{Plural} \\
\midrule
Nominative & frēond & frīend \\
Genitive & frēondes & frēonda \\
Dative & frēonde & frēondum \\
Accusative & frēond & frīend \\
\bottomrule
\end{tabular}
\par\endgroup

% mdtable:T14:end
\reserve{7\baselineskip}
% mdparagraph:P102:begin
\subsection*{Table 14. burg: feminine noun}
% mdparagraph:P102:end

% table: NEW-U13-NOM-BURG
% mdtable:T15:begin
\par\begingroup
\fontsize{10.5}{13.2}\selectfont
\setlength{\LTleft}{0pt}\setlength{\LTright}{\fill}
\noindent
\begin{tabular}{@{}>{\raggedright\arraybackslash}p{39.46mm}>{\raggedright\arraybackslash}p{39.46mm}>{\raggedright\arraybackslash}p{39.46mm}@{}}
\toprule
\textbf{Case} & \textbf{Singular} & \textbf{Plural} \\
\midrule
Nominative & burg & byrig \\
Genitive & byrig & burga \\
Dative & byrig & burgum \\
Accusative & burg & byrig \\
\bottomrule
\end{tabular}
\par\endgroup

% mdtable:T15:end
\reserve{7\baselineskip}
% mdparagraph:P103:begin
\subsection*{Table 15. word: neuter noun}
% mdparagraph:P103:end

% table: NEW-U13-NOM-WORD
% mdtable:T16:begin
\par\begingroup
\fontsize{10.5}{13.2}\selectfont
\setlength{\LTleft}{0pt}\setlength{\LTright}{\fill}
\noindent
\begin{tabular}{@{}>{\raggedright\arraybackslash}p{39.46mm}>{\raggedright\arraybackslash}p{39.46mm}>{\raggedright\arraybackslash}p{39.46mm}@{}}
\toprule
\textbf{Case} & \textbf{Singular} & \textbf{Plural} \\
\midrule
Nominative & word & word \\
Genitive & wordes & worda \\
Dative & worde & wordum \\
Accusative & word & word \\
\bottomrule
\end{tabular}
\par\endgroup

% mdtable:T16:end
\reserve{7\baselineskip}
% mdparagraph:P104:begin
\subsection*{6.2. Diagnostic stems and shared models}
% mdparagraph:P104:end

% mdparagraph:P105:begin
The following rows give every noun’s full diagnostic stem and its G singular, D singular and N plural. For a shared model, use the indicated table in the same case and number, replace its initial stem by the stated stem, and retain the complete ending. A zero ending remains zero. If the recipe specifies the citation form for uninflected N/\allowbreak{}A cells, use the full headword there. This is a complete route through both numbers, not a guess from one oblique form.
% mdparagraph:P105:end

% mdtable:T17:begin
% table-reflow: nominal-diagnostic-panels
\par\begingroup
\fontsize{10.5}{13.2}\selectfont
\setlength{\LTleft}{0pt}\setlength{\LTright}{\fill}
\begin{longtable}{@{}>{\raggedright\arraybackslash}p{48mm}>{\raggedright\arraybackslash}p{48mm}>{\raggedright\arraybackslash}p{22.37mm}@{}}
\toprule
\textbf{Noun} & \textbf{Stem} & \textbf{Model} \\
\midrule
\endfirsthead
\toprule
\textbf{Noun} & \textbf{Stem} & \textbf{Model} \\
\midrule
\endhead
\midrule
\multicolumn{3}{r}{\footnotesize Continued on the following page.}\\
\endfoot
\bottomrule
\endlastfoot
īdelnes & īdelness & Table 1 \\
mōdor & mōdr & Table 2 \\
synn & synn & Table 9 \\
dohtor & dohtr & Table 3 \\
bisceop & bisceop & Table 4 \\
mynster & mynstr & Table 8 \\
leornung & leornung & Table 5 \\
cirice & ciric & Table 6 \\
unrihtwīsnes & unrihtwīsness & Table 1 \\
sang & sang & Table 4 \\
līchama & līcham & Table 7 \\
ēagðȳrel & ēagðȳrl & Table 8 \\
dǣl & dǣl & Table 4 \\
rōd & rōd & Table 9 \\
æx & æx & Table 9 \\
adesa & ades & Table 7 \\
wyrttruma & wyrttrum & Table 7 \\
untrymnes & untrymness & Table 1 \\
gefērscipe & gefērscip & Table 10 \\
gelēafa & gelēaf & Table 7 \\
wundor & wundr & Table 11 \\
hęll & hęll & Table 9 \\
hrōf & hrōf & Table 4 \\
mēd & mēd & Table 9 \\
æcer & æcer & Table 4 \\
ofen & ofen & Table 4 \\
brōðor & brōðr & Table 12 \\
forðfōr & forðfōr & Table 9 \\
weg & weg & Table 4 \\
frēond & frēond /\allowbreak{} frīend & Table 13 \\
burg & burg /\allowbreak{} byrig & Table 14 \\
word & word & Table 15 \\
\end{longtable}
\par\addvspace{5pt}
\setlength{\LTleft}{0pt}\setlength{\LTright}{\fill}
\begin{longtable}{@{}>{\raggedright\arraybackslash}p{31mm}>{\raggedright\arraybackslash}p{29mm}>{\raggedright\arraybackslash}p{29mm}>{\raggedright\arraybackslash}p{26.56mm}@{}}
\toprule
\textbf{Noun} & \textbf{G singular} & \textbf{D singular} & \textbf{N plural} \\
\midrule
\endfirsthead
\toprule
\textbf{Noun} & \textbf{G singular} & \textbf{D singular} & \textbf{N plural} \\
\midrule
\endhead
\midrule
\multicolumn{4}{r}{\footnotesize Continued on the following page.}\\
\endfoot
\bottomrule
\endlastfoot
īdelnes & īdelnesse & īdelnesse & īdelnessa \\
mōdor & mōdor & mēder & mōdru \\
synn & synne & synne & synna \\
dohtor & dohtor & dehter & dohtor \\
bisceop & bisceopes & bisceope & bisceopas \\
mynster & mynstres & mynstre & mynstru \\
leornung & leornunge & leornunge & leornunga \\
cirice & cirican & cirican & cirican \\
unrihtwīsnes & unrihtwīsnesse & unrihtwīsnesse & unrihtwīsnessa \\
sang & sanges & sange & sangas \\
līchama & līchaman & līchaman & līchaman \\
ēagðȳrel & ēagðȳrles & ēagðȳrle & ēagðȳrlu \\
dǣl & dǣles & dǣle & dǣlas \\
rōd & rōde & rōde & rōda \\
æx & æxe & æxe & æxa \\
adesa & adesan & adesan & adesan \\
wyrttruma & wyrttruman & wyrttruman & wyrttruman \\
untrymnes & untrymnesse & untrymnesse & untrymnessa \\
gefērscipe & gefērscipes & gefērscipe & gefērscipas \\
gelēafa & gelēafan & gelēafan & gelēafan \\
wundor & wundres & wundre & wundor \\
hęll & hęlle & hęlle & hęlla \\
hrōf & hrōfes & hrōfe & hrōfas \\
mēd & mēde & mēde & mēda \\
æcer & æceres & æcere & æceras \\
ofen & ofenes & ofene & ofenas \\
brōðor & brōðor & brēðer & brōðor \\
forðfōr & forðfōre & forðfōre & forðfōra \\
weg & weges & wege & wegas \\
frēond & frēondes & frēonde & frīend \\
burg & byrig & byrig & byrig \\
word & wordes & worde & word \\
\end{longtable}
\par\addvspace{5pt}
\par\endgroup
% mdtable:T17:end

% mdtable:T18:begin
\par\begingroup
\fontsize{10.5}{13.2}\selectfont
\setlength{\LTleft}{0pt}\setlength{\LTright}{\fill}
\begin{longtable}{@{}>{\raggedright\arraybackslash}p{25mm}>{\raggedright\arraybackslash}p{19mm}>{\raggedright\arraybackslash}p{74.37mm}@{}}
\toprule
\textbf{Noun} & \textbf{Complete model} & \textbf{Operation in the corresponding cell} \\
\midrule
\endfirsthead
\toprule
\textbf{Noun} & \textbf{Complete model} & \textbf{Operation in the corresponding cell} \\
\midrule
\endhead
\midrule
\multicolumn{3}{r}{\footnotesize Continued on the following page.}\\
\endfoot
\bottomrule
\endlastfoot
synn & Table 9 & Replace æx- by synn-; retain the ending. \\
mynster & Table 8 & In N singular, A singular, use the full citation mynster. Replace ēagðȳrl- by mynstr-; retain the ending. \\
unrihtwīsnes & Table 1 & In N singular, use the full citation unrihtwīsnes. Replace īdelness- by unrihtwīsness-; retain the ending. \\
sang & Table 4 & Replace bisceop- by sang-; retain the ending. \\
dǣl & Table 4 & Replace bisceop- by dǣl-; retain the ending. \\
rōd & Table 9 & Replace æx- by rōd-; retain the ending. \\
adesa & Table 7 & Replace līcham- by ades-; retain the ending. \\
wyrttruma & Table 7 & Replace līcham- by wyrttrum-; retain the ending. \\
untrymnes & Table 1 & In N singular, use the full citation untrymnes. Replace īdelness- by untrymness-; retain the ending. \\
gelēafa & Table 7 & Replace līcham- by gelēaf-; retain the ending. \\
hęll & Table 9 & Replace æx- by hęll-; retain the ending. \\
hrōf & Table 4 & Replace bisceop- by hrōf-; retain the ending. \\
mēd & Table 9 & Replace æx- by mēd-; retain the ending. \\
æcer & Table 4 & Replace bisceop- by æcer-; retain the ending. \\
ofen & Table 4 & Replace bisceop- by ofen-; retain the ending. \\
forðfōr & Table 9 & Replace æx- by forðfōr-; retain the ending. \\
weg & Table 4 & Replace bisceop- by weg-; retain the ending. \\
\end{longtable}
\par\endgroup

% mdtable:T18:end
% mdparagraph:P106:begin
Recognize these documented alternatives without replacing the selected course forms: mōdra for N/\allowbreak{}A plural mōdru; dohtru or dohtra for N/\allowbreak{}A plural dohtor; brōðru for N/\allowbreak{}A plural brōðor; late G singular mēder and dehter; G singular burge beside byrig; D singular frīend and N/\allowbreak{}A plural frēond or frēondas beside the selected frēonde and frīend. Cook’s citation syn(n) includes syn beside selected synn. These alternatives do not change the listed meanings or case relations.
% mdparagraph:P106:end

\reserve{7\baselineskip}
% mdparagraph:P107:begin
\subsection*{6.3. Adjectives and comparison}
% mdparagraph:P107:end

% teaching: NEW-U13-GRAM-U13-ADJ-STRONG
% teaching: NEW-U13-GRAM-U13-ADJ-WEAK
% teaching: NEW-U13-GRAM-U13-ADJ-HALIG
% teaching: NEW-U13-GRAM-U13-ADJ-COMPARISON
% teaching: NEW-U13-VOC-HALIG
% teaching: NEW-U13-ADJ-HALIG
% teaching: NEW-U13-VOC-SWETE
% teaching: NEW-U13-ADJ-SWETE
% teaching: NEW-U13-VOC-LICHAMLIC
% teaching: NEW-U13-ADJ-LICHAMLIC
% teaching: NEW-U13-VOC-HEFIG
% teaching: NEW-U13-ADJ-HEFIG
% teaching: NEW-U13-VOC-WYNSUM
% teaching: NEW-U13-ADJ-WYNSUM
% teaching: NEW-U13-VOC-HWIT
% teaching: NEW-U13-ADJ-HWIT
% teaching: NEW-U13-VOC-WIERTHE
% teaching: NEW-U13-ADJ-WIERTHE
% teaching: NEW-U13-VOC-HEOFONLIC
% teaching: NEW-U13-ADJ-HEOFONLIC
% teaching: NEW-U13-VOC-SWETE-COMP
% teaching: NEW-U13-ADJ-SWETE-COMP
% teaching: NEW-U13-VOC-SWETE-SUPERL
% teaching: NEW-U13-ADJ-SWETE-SUPERL
% mdtable:T19:begin
\par\begingroup
\fontsize{10.5}{13.2}\selectfont
\setlength{\LTleft}{0pt}\setlength{\LTright}{\fill}
\noindent
\begin{tabular}{@{}>{\raggedright\arraybackslash}p{53.62mm}>{\raggedright\arraybackslash}p{67.56mm}@{}}
\toprule
\textbf{Adjective} & \textbf{Meaning} \\
\midrule
hālig & holy \\
swēte & sweet \\
līchamlic & bodily \\
hęfig & grievous \\
wynsum & pleasant; winsome \\
hwīt & white \\
wierðe & worthy \\
heofonlic & heavenly \\
swētra & sweeter \\
swētest & sweetest \\
\bottomrule
\end{tabular}
\par\endgroup

% mdtable:T19:end
% mdparagraph:P108:begin
Use strong and weak adjective forms in the environments already taught. The following complete tables cover N, G, D and A in all three genders and both numbers. Masculine and neuter singular instrumental forms are printed for reference; their display does not add independent productive instrumental syntax. Comparatives have weak inflection only. Superlatives have both strong and weak forms.
% mdparagraph:P108:end

% mdparagraph:P109:begin
An adjective in -e loses that citation ending before another ending: swēte gives swēt-, wierðe gives wierð-. Their strong feminine N singular and neuter N/\allowbreak{}A plural have -u. Long monosyllables and swētest omit that -u. Līchamlic, hęfig, wynsum and heofonlic retain it in those cells.
% mdparagraph:P109:end

% mdparagraph:P110:begin
Hālig needs its own full pattern: feminine N singular and neuter N/\allowbreak{}A plural are hālig. Before vowel endings use hālg-; before -ne, -re and -ra retain hālig-. Hęfig retains its medial vowel throughout the selected pattern. The weak G plural ending is selected -ra; -ena is also documented.
% mdparagraph:P110:end

% mdparagraph:P111:begin
Swēte forms comparative swētra and superlative swētest. In the comparative G plural, swētr- plus -ra simplifies the doubled r after t, giving swētra. Do not mechanically write swētrra. The shared model for hwīt below is morphological only: “white” does not acquire the meaning “sweetest.”
% mdparagraph:P111:end

\reserve{7\baselineskip}
% mdparagraph:P112:begin
\subsection*{Table 16. hālig: strong adjective}
% mdparagraph:P112:end

% table: NEW-U13-ADJ-HALIG-STRONG
% mdtable:T20:begin
% course-table-style: 2026-10
% course-panel:T20-P01:begin
\par\begingroup
\fontsize{10.5}{13.2}\selectfont
\setlength{\LTleft}{0pt}\setlength{\LTright}{\fill}
\noindent
\begin{tabular}{@{}>{\raggedright\arraybackslash}p{22.0mm}>{\raggedright\arraybackslash}p{31.188mm}>{\raggedright\arraybackslash}p{31.188mm}>{\raggedright\arraybackslash}p{31.188mm}@{}}
\toprule
\textbf{Case} & \textbf{Masculine singular} & \textbf{Neuter singular} & \textbf{Feminine singular} \\
\midrule
Nominative & hālig & hālig & hālig \\
Genitive & hālges & hālges & hāligre \\
Dative & hālgum & hālgum & hāligre \\
Accusative & hāligne & hālig & hālge \\
Instrumental & hālge & hālge & as dative \\
\bottomrule
\end{tabular}
\par\endgroup
% course-panel:T20-P01:end
\par\addvspace{5pt}
% course-panel:T20-P02:begin
\par\begingroup
\fontsize{10.5}{13.2}\selectfont
\setlength{\LTleft}{0pt}\setlength{\LTright}{\fill}
\noindent
\begin{tabular}{@{}>{\raggedright\arraybackslash}p{22.0mm}>{\raggedright\arraybackslash}p{31.188mm}>{\raggedright\arraybackslash}p{31.188mm}>{\raggedright\arraybackslash}p{31.188mm}@{}}
\toprule
\textbf{Case} & \textbf{Masculine plural} & \textbf{Neuter plural} & \textbf{Feminine plural} \\
\midrule
Nominative & hālge & hālig & hālga \\
Genitive & hāligra & hāligra & hāligra \\
Dative & hālgum & hālgum & hālgum \\
Accusative & hālge & hālig & hālga \\
Instrumental & as dative & as dative & as dative \\
\bottomrule
\end{tabular}
\par\endgroup
% course-panel:T20-P02:end
% mdtable:T20:end

\reserve{7\baselineskip}
% mdparagraph:P113:begin
\subsection*{Table 17. hālig: weak adjective}
% mdparagraph:P113:end

% table: NEW-U13-ADJ-HALIG-WEAK
% mdtable:T21:begin
% course-table-style: 2026-10
% course-panel:T21-P01:begin
\par\begingroup
\fontsize{10.5}{13.2}\selectfont
\setlength{\LTleft}{0pt}\setlength{\LTright}{\fill}
\noindent
\begin{tabular}{@{}>{\raggedright\arraybackslash}p{22.0mm}>{\raggedright\arraybackslash}p{31.188mm}>{\raggedright\arraybackslash}p{31.188mm}>{\raggedright\arraybackslash}p{31.188mm}@{}}
\toprule
\textbf{Case} & \textbf{Masculine singular} & \textbf{Neuter singular} & \textbf{Feminine singular} \\
\midrule
Nominative & hālga & hālge & hālge \\
Genitive & hālgan & hālgan & hālgan \\
Dative & hālgan & hālgan & hālgan \\
Accusative & hālgan & hālge & hālgan \\
\bottomrule
\end{tabular}
\par\endgroup
% course-panel:T21-P01:end
\par\addvspace{5pt}
% course-panel:T21-P02:begin
\par\begingroup
\fontsize{10.5}{13.2}\selectfont
\setlength{\LTleft}{0pt}\setlength{\LTright}{\fill}
\noindent
\begin{tabular}{@{}>{\raggedright\arraybackslash}p{22.0mm}>{\raggedright\arraybackslash}p{31.188mm}>{\raggedright\arraybackslash}p{31.188mm}>{\raggedright\arraybackslash}p{31.188mm}@{}}
\toprule
\textbf{Case} & \textbf{Masculine plural} & \textbf{Neuter plural} & \textbf{Feminine plural} \\
\midrule
Nominative & hālgan & hālgan & hālgan \\
Genitive & hāligra & hāligra & hāligra \\
Dative & hālgum & hālgum & hālgum \\
Accusative & hālgan & hālgan & hālgan \\
\bottomrule
\end{tabular}
\par\endgroup
% course-panel:T21-P02:end
% mdtable:T21:end

\reserve{7\baselineskip}
% mdparagraph:P114:begin
\subsection*{Table 18. swēte: strong adjective}
% mdparagraph:P114:end

% table: NEW-U13-ADJ-SWETE-STRONG
% mdtable:T22:begin
% course-table-style: 2026-10
% course-panel:T22-P01:begin
\par\begingroup
\fontsize{10.5}{13.2}\selectfont
\setlength{\LTleft}{0pt}\setlength{\LTright}{\fill}
\noindent
\begin{tabular}{@{}>{\raggedright\arraybackslash}p{22.0mm}>{\raggedright\arraybackslash}p{31.188mm}>{\raggedright\arraybackslash}p{31.188mm}>{\raggedright\arraybackslash}p{31.188mm}@{}}
\toprule
\textbf{Case} & \textbf{Masculine singular} & \textbf{Neuter singular} & \textbf{Feminine singular} \\
\midrule
Nominative & swēte & swēte & swētu \\
Genitive & swētes & swētes & swētre \\
Dative & swētum & swētum & swētre \\
Accusative & swētne & swēte & swēte \\
Instrumental & swēte & swēte & as dative \\
\bottomrule
\end{tabular}
\par\endgroup
% course-panel:T22-P01:end
\par\addvspace{5pt}
% course-panel:T22-P02:begin
\par\begingroup
\fontsize{10.5}{13.2}\selectfont
\setlength{\LTleft}{0pt}\setlength{\LTright}{\fill}
\noindent
\begin{tabular}{@{}>{\raggedright\arraybackslash}p{22.0mm}>{\raggedright\arraybackslash}p{31.188mm}>{\raggedright\arraybackslash}p{31.188mm}>{\raggedright\arraybackslash}p{31.188mm}@{}}
\toprule
\textbf{Case} & \textbf{Masculine plural} & \textbf{Neuter plural} & \textbf{Feminine plural} \\
\midrule
Nominative & swēte & swētu & swēta \\
Genitive & swētra & swētra & swētra \\
Dative & swētum & swētum & swētum \\
Accusative & swēte & swētu & swēta \\
Instrumental & as dative & as dative & as dative \\
\bottomrule
\end{tabular}
\par\endgroup
% course-panel:T22-P02:end
% mdtable:T22:end

\reserve{7\baselineskip}
% mdparagraph:P115:begin
\subsection*{Table 19. swēte: weak adjective}
% mdparagraph:P115:end

% table: NEW-U13-ADJ-SWETE-WEAK
% mdtable:T23:begin
% course-table-style: 2026-10
% course-panel:T23-P01:begin
\par\begingroup
\fontsize{10.5}{13.2}\selectfont
\setlength{\LTleft}{0pt}\setlength{\LTright}{\fill}
\noindent
\begin{tabular}{@{}>{\raggedright\arraybackslash}p{22.0mm}>{\raggedright\arraybackslash}p{31.188mm}>{\raggedright\arraybackslash}p{31.188mm}>{\raggedright\arraybackslash}p{31.188mm}@{}}
\toprule
\textbf{Case} & \textbf{Masculine singular} & \textbf{Neuter singular} & \textbf{Feminine singular} \\
\midrule
Nominative & swēta & swēte & swēte \\
Genitive & swētan & swētan & swētan \\
Dative & swētan & swētan & swētan \\
Accusative & swētan & swēte & swētan \\
\bottomrule
\end{tabular}
\par\endgroup
% course-panel:T23-P01:end
\par\addvspace{5pt}
% course-panel:T23-P02:begin
\par\begingroup
\fontsize{10.5}{13.2}\selectfont
\setlength{\LTleft}{0pt}\setlength{\LTright}{\fill}
\noindent
\begin{tabular}{@{}>{\raggedright\arraybackslash}p{22.0mm}>{\raggedright\arraybackslash}p{31.188mm}>{\raggedright\arraybackslash}p{31.188mm}>{\raggedright\arraybackslash}p{31.188mm}@{}}
\toprule
\textbf{Case} & \textbf{Masculine plural} & \textbf{Neuter plural} & \textbf{Feminine plural} \\
\midrule
Nominative & swētan & swētan & swētan \\
Genitive & swētra & swētra & swētra \\
Dative & swētum & swētum & swētum \\
Accusative & swētan & swētan & swētan \\
\bottomrule
\end{tabular}
\par\endgroup
% course-panel:T23-P02:end
% mdtable:T23:end

\reserve{7\baselineskip}
% mdparagraph:P116:begin
\subsection*{Table 20. līchamlic: strong adjective}
% mdparagraph:P116:end

% table: NEW-U13-ADJ-LICHAMLIC-STRONG
% mdtable:T24:begin
% course-table-style: 2026-10
% course-panel:T24-P01:begin
\par\begingroup
\fontsize{10.5}{13.2}\selectfont
\setlength{\LTleft}{0pt}\setlength{\LTright}{\fill}
\noindent
\begin{tabular}{@{}>{\raggedright\arraybackslash}p{22.0mm}>{\raggedright\arraybackslash}p{31.188mm}>{\raggedright\arraybackslash}p{31.188mm}>{\raggedright\arraybackslash}p{31.188mm}@{}}
\toprule
\textbf{Case} & \textbf{Masculine singular} & \textbf{Neuter singular} & \textbf{Feminine singular} \\
\midrule
Nominative & līchamlic & līchamlic & līchamlicu \\
Genitive & līchamlices & līchamlices & līchamlicre \\
Dative & līchamlicum & līchamlicum & līchamlicre \\
Accusative & līchamlicne & līchamlic & līchamlice \\
Instrumental & līchamlice & līchamlice & as dative \\
\bottomrule
\end{tabular}
\par\endgroup
% course-panel:T24-P01:end
\par\addvspace{5pt}
% course-panel:T24-P02:begin
\par\begingroup
\fontsize{10.5}{13.2}\selectfont
\setlength{\LTleft}{0pt}\setlength{\LTright}{\fill}
\noindent
\begin{tabular}{@{}>{\raggedright\arraybackslash}p{22.0mm}>{\raggedright\arraybackslash}p{31.188mm}>{\raggedright\arraybackslash}p{31.188mm}>{\raggedright\arraybackslash}p{31.188mm}@{}}
\toprule
\textbf{Case} & \textbf{Masculine plural} & \textbf{Neuter plural} & \textbf{Feminine plural} \\
\midrule
Nominative & līchamlice & līchamlicu & līchamlica \\
Genitive & līchamlicra & līchamlicra & līchamlicra \\
Dative & līchamlicum & līchamlicum & līchamlicum \\
Accusative & līchamlice & līchamlicu & līchamlica \\
Instrumental & as dative & as dative & as dative \\
\bottomrule
\end{tabular}
\par\endgroup
% course-panel:T24-P02:end
% mdtable:T24:end

\reserve{7\baselineskip}
% mdparagraph:P117:begin
\subsection*{Table 21. līchamlic: weak adjective}
% mdparagraph:P117:end

% table: NEW-U13-ADJ-LICHAMLIC-WEAK
% mdtable:T25:begin
% course-table-style: 2026-10
% course-panel:T25-P01:begin
\par\begingroup
\fontsize{10.5}{13.2}\selectfont
\setlength{\LTleft}{0pt}\setlength{\LTright}{\fill}
\noindent
\begin{tabular}{@{}>{\raggedright\arraybackslash}p{22.0mm}>{\raggedright\arraybackslash}p{31.188mm}>{\raggedright\arraybackslash}p{31.188mm}>{\raggedright\arraybackslash}p{31.188mm}@{}}
\toprule
\textbf{Case} & \textbf{Masculine singular} & \textbf{Neuter singular} & \textbf{Feminine singular} \\
\midrule
Nominative & līchamlica & līchamlice & līchamlice \\
Genitive & līchamlican & līchamlican & līchamlican \\
Dative & līchamlican & līchamlican & līchamlican \\
Accusative & līchamlican & līchamlice & līchamlican \\
\bottomrule
\end{tabular}
\par\endgroup
% course-panel:T25-P01:end
\par\addvspace{5pt}
% course-panel:T25-P02:begin
\par\begingroup
\fontsize{10.5}{13.2}\selectfont
\setlength{\LTleft}{0pt}\setlength{\LTright}{\fill}
\noindent
\begin{tabular}{@{}>{\raggedright\arraybackslash}p{22.0mm}>{\raggedright\arraybackslash}p{31.188mm}>{\raggedright\arraybackslash}p{31.188mm}>{\raggedright\arraybackslash}p{31.188mm}@{}}
\toprule
\textbf{Case} & \textbf{Masculine plural} & \textbf{Neuter plural} & \textbf{Feminine plural} \\
\midrule
Nominative & līchamlican & līchamlican & līchamlican \\
Genitive & līchamlicra & līchamlicra & līchamlicra \\
Dative & līchamlicum & līchamlicum & līchamlicum \\
Accusative & līchamlican & līchamlican & līchamlican \\
\bottomrule
\end{tabular}
\par\endgroup
% course-panel:T25-P02:end
% mdtable:T25:end

\reserve{7\baselineskip}
% mdparagraph:P118:begin
\subsection*{Table 22. swētra: weak adjective}
% mdparagraph:P118:end

% table: NEW-U13-ADJ-SWETE-COMP-WEAK
% mdtable:T26:begin
% course-table-style: 2026-10
% course-panel:T26-P01:begin
\par\begingroup
\fontsize{10.5}{13.2}\selectfont
\setlength{\LTleft}{0pt}\setlength{\LTright}{\fill}
\noindent
\begin{tabular}{@{}>{\raggedright\arraybackslash}p{22.0mm}>{\raggedright\arraybackslash}p{31.188mm}>{\raggedright\arraybackslash}p{31.188mm}>{\raggedright\arraybackslash}p{31.188mm}@{}}
\toprule
\textbf{Case} & \textbf{Masculine singular} & \textbf{Neuter singular} & \textbf{Feminine singular} \\
\midrule
Nominative & swētra & swētre & swētre \\
Genitive & swētran & swētran & swētran \\
Dative & swētran & swētran & swētran \\
Accusative & swētran & swētre & swētran \\
\bottomrule
\end{tabular}
\par\endgroup
% course-panel:T26-P01:end
\par\addvspace{5pt}
% course-panel:T26-P02:begin
\par\begingroup
\fontsize{10.5}{13.2}\selectfont
\setlength{\LTleft}{0pt}\setlength{\LTright}{\fill}
\noindent
\begin{tabular}{@{}>{\raggedright\arraybackslash}p{22.0mm}>{\raggedright\arraybackslash}p{31.188mm}>{\raggedright\arraybackslash}p{31.188mm}>{\raggedright\arraybackslash}p{31.188mm}@{}}
\toprule
\textbf{Case} & \textbf{Masculine plural} & \textbf{Neuter plural} & \textbf{Feminine plural} \\
\midrule
Nominative & swētran & swētran & swētran \\
Genitive & swētra & swētra & swētra \\
Dative & swētrum & swētrum & swētrum \\
Accusative & swētran & swētran & swētran \\
\bottomrule
\end{tabular}
\par\endgroup
% course-panel:T26-P02:end
% mdtable:T26:end

\reserve{7\baselineskip}
% mdparagraph:P119:begin
\subsection*{Table 23. swētest: strong adjective}
% mdparagraph:P119:end

% table: NEW-U13-ADJ-SWETE-SUPERL-STRONG
% mdtable:T27:begin
% course-table-style: 2026-10
% course-panel:T27-P01:begin
\par\begingroup
\fontsize{10.5}{13.2}\selectfont
\setlength{\LTleft}{0pt}\setlength{\LTright}{\fill}
\noindent
\begin{tabular}{@{}>{\raggedright\arraybackslash}p{22.0mm}>{\raggedright\arraybackslash}p{31.188mm}>{\raggedright\arraybackslash}p{31.188mm}>{\raggedright\arraybackslash}p{31.188mm}@{}}
\toprule
\textbf{Case} & \textbf{Masculine singular} & \textbf{Neuter singular} & \textbf{Feminine singular} \\
\midrule
Nominative & swētest & swētest & swētest \\
Genitive & swētestes & swētestes & swētestre \\
Dative & swētestum & swētestum & swētestre \\
Accusative & swētestne & swētest & swēteste \\
Instrumental & swēteste & swēteste & as dative \\
\bottomrule
\end{tabular}
\par\endgroup
% course-panel:T27-P01:end
\par\addvspace{5pt}
% course-panel:T27-P02:begin
\par\begingroup
\fontsize{10.5}{13.2}\selectfont
\setlength{\LTleft}{0pt}\setlength{\LTright}{\fill}
\noindent
\begin{tabular}{@{}>{\raggedright\arraybackslash}p{22.0mm}>{\raggedright\arraybackslash}p{31.188mm}>{\raggedright\arraybackslash}p{31.188mm}>{\raggedright\arraybackslash}p{31.188mm}@{}}
\toprule
\textbf{Case} & \textbf{Masculine plural} & \textbf{Neuter plural} & \textbf{Feminine plural} \\
\midrule
Nominative & swēteste & swētest & swētesta \\
Genitive & swētestra & swētestra & swētestra \\
Dative & swētestum & swētestum & swētestum \\
Accusative & swēteste & swētest & swētesta \\
Instrumental & as dative & as dative & as dative \\
\bottomrule
\end{tabular}
\par\endgroup
% course-panel:T27-P02:end
% mdtable:T27:end

\reserve{7\baselineskip}
% mdparagraph:P120:begin
\subsection*{Table 24. swētest: weak adjective}
% mdparagraph:P120:end

% table: NEW-U13-ADJ-SWETE-SUPERL-WEAK
% mdtable:T28:begin
% course-table-style: 2026-10
% course-panel:T28-P01:begin
\par\begingroup
\fontsize{10.5}{13.2}\selectfont
\setlength{\LTleft}{0pt}\setlength{\LTright}{\fill}
\noindent
\begin{tabular}{@{}>{\raggedright\arraybackslash}p{22.0mm}>{\raggedright\arraybackslash}p{31.188mm}>{\raggedright\arraybackslash}p{31.188mm}>{\raggedright\arraybackslash}p{31.188mm}@{}}
\toprule
\textbf{Case} & \textbf{Masculine singular} & \textbf{Neuter singular} & \textbf{Feminine singular} \\
\midrule
Nominative & swētesta & swēteste & swēteste \\
Genitive & swētestan & swētestan & swētestan \\
Dative & swētestan & swētestan & swētestan \\
Accusative & swētestan & swēteste & swētestan \\
\bottomrule
\end{tabular}
\par\endgroup
% course-panel:T28-P01:end
\par\addvspace{5pt}
% course-panel:T28-P02:begin
\par\begingroup
\fontsize{10.5}{13.2}\selectfont
\setlength{\LTleft}{0pt}\setlength{\LTright}{\fill}
\noindent
\begin{tabular}{@{}>{\raggedright\arraybackslash}p{22.0mm}>{\raggedright\arraybackslash}p{31.188mm}>{\raggedright\arraybackslash}p{31.188mm}>{\raggedright\arraybackslash}p{31.188mm}@{}}
\toprule
\textbf{Case} & \textbf{Masculine plural} & \textbf{Neuter plural} & \textbf{Feminine plural} \\
\midrule
Nominative & swētestan & swētestan & swētestan \\
Genitive & swētestra & swētestra & swētestra \\
Dative & swētestum & swētestum & swētestum \\
Accusative & swētestan & swētestan & swētestan \\
\bottomrule
\end{tabular}
\par\endgroup
% course-panel:T28-P02:end
% mdtable:T28:end

% mdtable:T29:begin
\par\begingroup
\fontsize{10.5}{13.2}\selectfont
\setlength{\LTleft}{0pt}\setlength{\LTright}{\fill}
\noindent
\begin{tabular}{@{}>{\raggedright\arraybackslash}p{16.83mm}>{\raggedright\arraybackslash}p{16.83mm}>{\raggedright\arraybackslash}p{29.87mm}>{\raggedright\arraybackslash}p{28.19mm}>{\raggedright\arraybackslash}p{21.04mm}@{}}
\toprule
\textbf{Adjective} & \textbf{Stem} & \textbf{Strong feminine N singular} & \textbf{Weak masculine N singular} & \textbf{Complete model(s)} \\
\midrule
hālig & hālig & hālig & hālga & Table 16, Table 17 \\
swēte & swēt & swētu & swēta & Table 18, Table 19 \\
līchamlic & līchamlic & līchamlicu & līchamlica & Table 20, Table 21 \\
hęfig & hęfig & hęfigu & hęfiga & Table 20, Table 21 \\
wynsum & wynsum & wynsumu & wynsuma & Table 20, Table 21 \\
hwīt & hwīt & hwīt & hwīta & Table 23, Table 24 \\
wierðe & wierð & wierðu & wierða & Table 18, Table 19 \\
heofonlic & heofonlic & heofonlicu & heofonlica & Table 20, Table 21 \\
swētra & swētr & No strong comparative & swētra & Table 22 \\
swētest & swētest & swētest & swētesta & Table 23, Table 24 \\
\bottomrule
\end{tabular}
\par\endgroup

% mdtable:T29:end
% mdtable:T30:begin
\par\begingroup
\fontsize{10.5}{13.2}\selectfont
\setlength{\LTleft}{0pt}\setlength{\LTright}{\fill}
\noindent
\begin{tabular}{@{}>{\raggedright\arraybackslash}p{22.31mm}>{\raggedright\arraybackslash}p{27.89mm}>{\raggedright\arraybackslash}p{68.17mm}@{}}
\toprule
\textbf{Adjective} & \textbf{Complete model(s)} & \textbf{Operation in the corresponding cell} \\
\midrule
hęfig & Table 20, Table 21 & Replace līchamlic- by hęfig-; retain the ending. \\
wynsum & Table 20, Table 21 & Replace līchamlic- by wynsum-; retain the ending. \\
hwīt & Table 23, Table 24 & Replace swētest- by hwīt-; retain the ending. \\
wierðe & Table 18, Table 19 & Replace swēt- by wierð-; retain the ending. \\
heofonlic & Table 20, Table 21 & Replace līchamlic- by heofonlic-; retain the ending. \\
\bottomrule
\end{tabular}
\par\endgroup

% mdtable:T30:end
\reserve{7\baselineskip}
% mdparagraph:P121:begin
\section*{7. Verbal systems}
% mdparagraph:P121:end

% mdparagraph:P122:begin
The finite tables give all three persons in both numbers. Present and preterite indicative are separate from the subjunctive, imperative, infinitive and participle systems. Learn each listed system as a whole, but use it only with its taught meaning and construction. A citation infinitive used to name a verb does not by itself assign that infinitive for sentence production.
% mdparagraph:P122:end

% teaching: NEW-U13-VOC-FORDILGIAN
% teaching: NEW-U13-PAR-FORDILGIAN
% teaching: NEW-U13-VOC-GEWUNIAN
% teaching: NEW-U13-PAR-GEWUNIAN
% teaching: NEW-U13-VOC-ASCIAN
% teaching: NEW-U13-PAR-ASCIAN
% teaching: NEW-U13-VOC-HALSIAN
% teaching: NEW-U13-PAR-HALSIAN
% teaching: NEW-U13-VOC-WUNIAN
% teaching: NEW-U13-PAR-WUNIAN
% teaching: NEW-U13-VOC-HATIAN
% teaching: NEW-U13-PAR-HATIAN
% teaching: NEW-U13-VOC-ASTRECCEAN
% teaching: NEW-U13-PAR-ASTRECCEAN
% teaching: NEW-U13-VOC-ONLIESAN
% teaching: NEW-U13-PAR-ONLIESAN
% teaching: NEW-U13-VOC-NEALAECAN
% teaching: NEW-U13-PAR-NEALAECAN
% teaching: NEW-U13-VOC-FORCEORFAN
% teaching: NEW-U13-PAR-FORCEORFAN
% teaching: NEW-U13-VOC-ASETTAN
% teaching: NEW-U13-PAR-ASETTAN
% teaching: NEW-U13-VOC-OFERWINNAN
% teaching: NEW-U13-PAR-OFERWINNAN
% teaching: NEW-U13-VOC-GESCRYDAN
% teaching: NEW-U13-PAR-GESCRYDAN
% teaching: NEW-U13-VOC-ONDRAEDAN
% teaching: NEW-U13-PAR-ONDRAEDAN
% teaching: NEW-U13-VOC-FORDON
% teaching: NEW-U13-PAR-FORDON
% teaching: NEW-U13-VOC-INGAN
% teaching: NEW-U13-PAR-INGAN
% teaching: NEW-U13-VOC-ASENDAN
% teaching: NEW-U13-PAR-ASENDAN
% teaching: NEW-U13-VOC-BEHEALDAN
% teaching: NEW-U13-PAR-BEHEALDAN
% teaching: NEW-U13-VOC-BEBEODAN
% teaching: NEW-U13-PAR-BEBEODAN
% teaching: NEW-U13-VOC-EFTHWEORFAN
% teaching: NEW-U13-PAR-EFTHWEORFAN
% teaching: NEW-U13-VOC-LAEDAN
% teaching: NEW-U13-PAR-LAEDAN
% teaching: NEW-U13-VOC-SPRECAN
% teaching: NEW-U13-PAR-SPRECAN
% teaching: NEW-U13-VOC-BEGAN
% teaching: NEW-U13-VOC-ASTIGAN
% teaching: NEW-U13-VOC-SELLAN
% teaching: NEW-U13-VOC-WYRCEAN
% mdtable:T31:begin
\par\begingroup
\fontsize{10.5}{13.2}\selectfont
\setlength{\LTleft}{0pt}\setlength{\LTright}{\fill}
\begin{longtable}{@{}>{\raggedright\arraybackslash}p{44.96mm}>{\raggedright\arraybackslash}p{76.23mm}@{}}
\toprule
\textbf{Verb} & \textbf{Meaning} \\
\midrule
\endfirsthead
\toprule
\textbf{Verb} & \textbf{Meaning} \\
\midrule
\endhead
\midrule
\multicolumn{2}{r}{\footnotesize Continued on the following page.}\\
\endfoot
\bottomrule
\endlastfoot
fordilgian & destroy; annihilate \\
gewunian & be accustomed; be wont \\
āscian & ask a person about something \\
hālsian & entreat; implore \\
wunian & dwell; remain \\
hatian & hate \\
āstręccean & prostrate (oneself) \\
onlīesan & release \\
nēalǣcan & approach; draw near \\
forceorfan & cut down; cut off \\
āsęttan & set; place \\
oferwinnan & overcome; conquer \\
gescrȳdan & clothe \\
ondrǣdan & fear \\
fordōn & destroy \\
ingān & enter; go in \\
āsęndan & send \\
behealdan & behold; look at \\
bebēodan & command; instruct \\
ęfthweorfan & return; turn back \\
lǣdan & lead; bring; take a person \\
sprecan & speak with someone \\
begān & pursue; practise study \\
āstīgan & descend \\
sęllan & sell \\
wyrcean & make; work; perform \\
\end{longtable}
\par\endgroup

% mdtable:T31:end
\reserve{7\baselineskip}
% mdparagraph:P123:begin
\subsection*{7.1. Weak verbs}
% mdparagraph:P123:end

% mdparagraph:P124:begin
The selected weak II pattern has present -ie, -ast, -að, plural -iað; preterite -ode, -odest, -ode, plural -odon. The finite model below supplies every indicative cell for the six verbs in the diagnostic table. Alternative -ige and -ade formations are outside the selected pattern.
% mdparagraph:P124:end

% teaching: NEW-U13-PAR-FORDILGIAN-IND
% teaching: NEW-U13-GRAM-U13-FORDILGIAN-INDICATIVE
\reserve{7\baselineskip}
% mdparagraph:P125:begin
\subsection*{Table 25. fordilgian: present and preterite indicative}
% mdparagraph:P125:end

% table: NEW-U13-PAR-FORDILGIAN-IND
% mdtable:T32:begin
% course-table-style: 2026-10
% course-panel:T32-P01:begin
\par\begingroup
\fontsize{10.5}{13.2}\selectfont
\setlength{\LTleft}{0pt}\setlength{\LTright}{\fill}
\noindent
\begin{tabular}{@{}>{\raggedright\arraybackslash}p{39.459mm}>{\raggedright\arraybackslash}p{39.459mm}>{\raggedright\arraybackslash}p{39.459mm}@{}}
\toprule
\textbf{Person and number} & \textbf{Present} & \textbf{Preterite} \\
\midrule
1 sg. & fordilgie & fordilgode \\
2 sg. & fordilgast & fordilgodest \\
3 sg. & fordilgað & fordilgode \\
1–3 pl. & fordilgiað & fordilgodon \\
\bottomrule
\end{tabular}
\par\endgroup
% course-panel:T32-P01:end
% mdtable:T32:end
% teaching: NEW-U13-PAR-GEWUNIAN-IND
% teaching: NEW-U13-GRAM-U13-GEWUNIAN-INDICATIVE
% teaching: NEW-U13-PAR-ASCIAN-IND
% teaching: NEW-U13-GRAM-U13-ASCIAN-INDICATIVE
% teaching: NEW-U13-PAR-HALSIAN-IND
% teaching: NEW-U13-GRAM-U13-HALSIAN-INDICATIVE
% teaching: NEW-U13-PAR-WUNIAN-IND
% teaching: NEW-U13-GRAM-U13-WUNIAN-INDICATIVE
% teaching: NEW-U13-PAR-HATIAN-IND
% teaching: NEW-U13-GRAM-U13-HATIAN-INDICATIVE
% mdtable:T33:begin
\par\begingroup
\fontsize{10.5}{13.2}\selectfont
\setlength{\LTleft}{0pt}\setlength{\LTright}{\fill}
\noindent
\begin{tabular}{@{}>{\raggedright\arraybackslash}p{21.33mm}>{\raggedright\arraybackslash}p{21.33mm}>{\raggedright\arraybackslash}p{22.08mm}>{\raggedright\arraybackslash}p{25.17mm}>{\raggedright\arraybackslash}p{22.83mm}@{}}
\toprule
\textbf{Verb} & \textbf{Stem} & \textbf{Present 3 sg} & \textbf{Preterite 3 sg} & \textbf{Complete model} \\
\midrule
fordilgian & fordilg & fordilgað & fordilgode & Table 25 \\
gewunian & gewun & gewunað & gewunode & Table 25 \\
āscian & āsc & āscað & āscode & Table 25 \\
hālsian & hāls & hālsað & hālsode & Table 25 \\
wunian & wun & wunað & wunode & Table 25 \\
hatian & hat & hatað & hatode & Table 25 \\
\bottomrule
\end{tabular}
\par\endgroup

% mdtable:T33:end
% mdparagraph:P126:begin
For gewunian, āscian, hālsian, wunian and hatian, replace initial fordilg- in every cell of Table 25 by gewun-, āsc-, hāls-, wun- or hat-, respectively. Retain the rest of the cell exactly. This gives all twelve present and preterite indicative positions for each verb; no mood or nonfinite system is added by this substitution.
% mdparagraph:P126:end

% mdparagraph:P127:begin
Weak I verbs add a dental preterite, but their present stems and consonant outcomes must be learned separately. Āstręccean alternates geminate and ungeminated present consonants; its unmutated past is āstreahte, not a past formed mechanically from stręcc-. Nēalǣcan selects nēalǣcte. The tables retain the distinct present, past and consonant patterns of onlīesan, āsęttan, gescrȳdan and āsęndan.
% mdparagraph:P127:end

% teaching: NEW-U13-PAR-ASTRECCEAN-IND
% teaching: NEW-U13-GRAM-U13-ASTRECCEAN-INDICATIVE
\reserve{7\baselineskip}
% mdparagraph:P128:begin
\subsection*{Table 26. āstręccean: present and preterite indicative}
% mdparagraph:P128:end

% table: NEW-U13-PAR-ASTRECCEAN-IND
% mdtable:T34:begin
% course-table-style: 2026-10
% course-panel:T34-P01:begin
\par\begingroup
\fontsize{10.5}{13.2}\selectfont
\setlength{\LTleft}{0pt}\setlength{\LTright}{\fill}
\noindent
\begin{tabular}{@{}>{\raggedright\arraybackslash}p{39.459mm}>{\raggedright\arraybackslash}p{39.459mm}>{\raggedright\arraybackslash}p{39.459mm}@{}}
\toprule
\textbf{Person and number} & \textbf{Present} & \textbf{Preterite} \\
\midrule
1 sg. & āstręcce & āstreahte \\
2 sg. & āstręcst & āstreahtest \\
3 sg. & āstręcð & āstreahte \\
1–3 pl. & āstręcceað & āstreahton \\
\bottomrule
\end{tabular}
\par\endgroup
% course-panel:T34-P01:end
% mdtable:T34:end
% teaching: NEW-U13-PAR-ONLIESAN-IND
% teaching: NEW-U13-GRAM-U13-ONLIESAN-INDICATIVE
\reserve{7\baselineskip}
% mdparagraph:P129:begin
\subsection*{Table 27. onlīesan: present and preterite indicative}
% mdparagraph:P129:end

% table: NEW-U13-PAR-ONLIESAN-IND
% mdtable:T35:begin
% course-table-style: 2026-10
% course-panel:T35-P01:begin
\par\begingroup
\fontsize{10.5}{13.2}\selectfont
\setlength{\LTleft}{0pt}\setlength{\LTright}{\fill}
\noindent
\begin{tabular}{@{}>{\raggedright\arraybackslash}p{39.459mm}>{\raggedright\arraybackslash}p{39.459mm}>{\raggedright\arraybackslash}p{39.459mm}@{}}
\toprule
\textbf{Person and number} & \textbf{Present} & \textbf{Preterite} \\
\midrule
1 sg. & onlīese & onlīesde \\
2 sg. & onlīest & onlīesdest \\
3 sg. & onlīest & onlīesde \\
1–3 pl. & onlīesað & onlīesdon \\
\bottomrule
\end{tabular}
\par\endgroup
% course-panel:T35-P01:end
% mdtable:T35:end
% teaching: NEW-U13-PAR-NEALAECAN-IND
% teaching: NEW-U13-GRAM-U13-NEALAECAN-INDICATIVE
\reserve{7\baselineskip}
% mdparagraph:P130:begin
\subsection*{Table 28. nēalǣcan: present and preterite indicative}
% mdparagraph:P130:end

% table: NEW-U13-PAR-NEALAECAN-IND
% mdtable:T36:begin
% course-table-style: 2026-10
% course-panel:T36-P01:begin
\par\begingroup
\fontsize{10.5}{13.2}\selectfont
\setlength{\LTleft}{0pt}\setlength{\LTright}{\fill}
\noindent
\begin{tabular}{@{}>{\raggedright\arraybackslash}p{39.459mm}>{\raggedright\arraybackslash}p{39.459mm}>{\raggedright\arraybackslash}p{39.459mm}@{}}
\toprule
\textbf{Person and number} & \textbf{Present} & \textbf{Preterite} \\
\midrule
1 sg. & nēalǣce & nēalǣcte \\
2 sg. & nēalǣcst & nēalǣctest \\
3 sg. & nēalǣcð & nēalǣcte \\
1–3 pl. & nēalǣcað & nēalǣcton \\
\bottomrule
\end{tabular}
\par\endgroup
% course-panel:T36-P01:end
% mdtable:T36:end
% teaching: NEW-U13-PAR-ASETTAN-IND
% teaching: NEW-U13-GRAM-U13-ASETTAN-INDICATIVE
\reserve{7\baselineskip}
% mdparagraph:P131:begin
\subsection*{Table 29. āsęttan: present and preterite indicative}
% mdparagraph:P131:end

% table: NEW-U13-PAR-ASETTAN-IND
% mdtable:T37:begin
% course-table-style: 2026-10
% course-panel:T37-P01:begin
\par\begingroup
\fontsize{10.5}{13.2}\selectfont
\setlength{\LTleft}{0pt}\setlength{\LTright}{\fill}
\noindent
\begin{tabular}{@{}>{\raggedright\arraybackslash}p{39.459mm}>{\raggedright\arraybackslash}p{39.459mm}>{\raggedright\arraybackslash}p{39.459mm}@{}}
\toprule
\textbf{Person and number} & \textbf{Present} & \textbf{Preterite} \\
\midrule
1 sg. & āsętte & āsętte \\
2 sg. & āsętst & āsęttest \\
3 sg. & āsętt & āsętte \\
1–3 pl. & āsęttað & āsętton \\
\bottomrule
\end{tabular}
\par\endgroup
% course-panel:T37-P01:end
% mdtable:T37:end
% teaching: NEW-U13-PAR-GESCRYDAN-IND
% teaching: NEW-U13-GRAM-U13-GESCRYDAN-INDICATIVE
\reserve{7\baselineskip}
% mdparagraph:P132:begin
\subsection*{Table 30. gescrȳdan: present and preterite indicative}
% mdparagraph:P132:end

% table: NEW-U13-PAR-GESCRYDAN-IND
% mdtable:T38:begin
% course-table-style: 2026-10
% course-panel:T38-P01:begin
\par\begingroup
\fontsize{10.5}{13.2}\selectfont
\setlength{\LTleft}{0pt}\setlength{\LTright}{\fill}
\noindent
\begin{tabular}{@{}>{\raggedright\arraybackslash}p{39.459mm}>{\raggedright\arraybackslash}p{39.459mm}>{\raggedright\arraybackslash}p{39.459mm}@{}}
\toprule
\textbf{Person and number} & \textbf{Present} & \textbf{Preterite} \\
\midrule
1 sg. & gescrȳde & gescrȳdde \\
2 sg. & gescrȳtst & gescrȳddest \\
3 sg. & gescrȳt & gescrȳdde \\
1–3 pl. & gescrȳdað & gescrȳddon \\
\bottomrule
\end{tabular}
\par\endgroup
% course-panel:T38-P01:end
% mdtable:T38:end
% teaching: NEW-U13-PAR-ASENDAN-IND
% teaching: NEW-U13-GRAM-U13-ASENDAN-INDICATIVE
\reserve{7\baselineskip}
% mdparagraph:P133:begin
\subsection*{Table 31. āsęndan: present and preterite indicative}
% mdparagraph:P133:end

% table: NEW-U13-PAR-ASENDAN-IND
% mdtable:T39:begin
% course-table-style: 2026-10
% course-panel:T39-P01:begin
\par\begingroup
\fontsize{10.5}{13.2}\selectfont
\setlength{\LTleft}{0pt}\setlength{\LTright}{\fill}
\noindent
\begin{tabular}{@{}>{\raggedright\arraybackslash}p{39.459mm}>{\raggedright\arraybackslash}p{39.459mm}>{\raggedright\arraybackslash}p{39.459mm}@{}}
\toprule
\textbf{Person and number} & \textbf{Present} & \textbf{Preterite} \\
\midrule
1 sg. & āsęnde & āsęnde \\
2 sg. & āsęntst & āsęndest \\
3 sg. & āsęnt & āsęnde \\
1–3 pl. & āsęndað & āsęndon \\
\bottomrule
\end{tabular}
\par\endgroup
% course-panel:T39-P01:end
% mdtable:T39:end
\reserve{7\baselineskip}
% mdparagraph:P134:begin
\subsection*{7.2. Strong and irregular finite verbs}
% mdparagraph:P134:end

% mdparagraph:P135:begin
Strong verbs distinguish present, preterite singular, preterite plural and past-participle stems. The second singular preterite takes the same vowel grade as the plural. In forceorfan the present second and third singular use forcierf-; the past distinguishes forcearf from forcurf-. Oferwinnan and ęfthweorfan have their own related strong III patterns. Keep every change shown in the tables, including nn plus ð becoming nð and the selected consonant clusters before -st.
% mdparagraph:P135:end

% teaching: NEW-U13-PAR-FORCEORFAN-IND
% teaching: NEW-U13-GRAM-U13-FORCEORFAN-INDICATIVE
\reserve{7\baselineskip}
% mdparagraph:P136:begin
\subsection*{Table 32. forceorfan: present and preterite indicative}
% mdparagraph:P136:end

% table: NEW-U13-PAR-FORCEORFAN-IND
% mdtable:T40:begin
% course-table-style: 2026-10
% course-panel:T40-P01:begin
\par\begingroup
\fontsize{10.5}{13.2}\selectfont
\setlength{\LTleft}{0pt}\setlength{\LTright}{\fill}
\noindent
\begin{tabular}{@{}>{\raggedright\arraybackslash}p{39.459mm}>{\raggedright\arraybackslash}p{39.459mm}>{\raggedright\arraybackslash}p{39.459mm}@{}}
\toprule
\textbf{Person and number} & \textbf{Present} & \textbf{Preterite} \\
\midrule
1 sg. & forceorfe & forcearf \\
2 sg. & forcierfst & forcurfe \\
3 sg. & forcierfð & forcearf \\
1–3 pl. & forceorfað & forcurfon \\
\bottomrule
\end{tabular}
\par\endgroup
% course-panel:T40-P01:end
% mdtable:T40:end
% teaching: NEW-U13-PAR-OFERWINNAN-IND
% teaching: NEW-U13-GRAM-U13-OFERWINNAN-INDICATIVE
\reserve{7\baselineskip}
% mdparagraph:P137:begin
\subsection*{Table 33. oferwinnan: present and preterite indicative}
% mdparagraph:P137:end

% table: NEW-U13-PAR-OFERWINNAN-IND
% mdtable:T41:begin
% course-table-style: 2026-10
% course-panel:T41-P01:begin
\par\begingroup
\fontsize{10.5}{13.2}\selectfont
\setlength{\LTleft}{0pt}\setlength{\LTright}{\fill}
\noindent
\begin{tabular}{@{}>{\raggedright\arraybackslash}p{39.459mm}>{\raggedright\arraybackslash}p{39.459mm}>{\raggedright\arraybackslash}p{39.459mm}@{}}
\toprule
\textbf{Person and number} & \textbf{Present} & \textbf{Preterite} \\
\midrule
1 sg. & oferwinne & oferwann \\
2 sg. & oferwinst & oferwunne \\
3 sg. & oferwinð & oferwann \\
1–3 pl. & oferwinnað & oferwunnon \\
\bottomrule
\end{tabular}
\par\endgroup
% course-panel:T41-P01:end
% mdtable:T41:end
% teaching: NEW-U13-PAR-EFTHWEORFAN-IND
% teaching: NEW-U13-GRAM-U13-EFTHWEORFAN-INDICATIVE
\reserve{7\baselineskip}
% mdparagraph:P138:begin
\subsection*{Table 34. ęfthweorfan: present and preterite indicative}
% mdparagraph:P138:end

% table: NEW-U13-PAR-EFTHWEORFAN-IND
% mdtable:T42:begin
% course-table-style: 2026-10
% course-panel:T42-P01:begin
\par\begingroup
\fontsize{10.5}{13.2}\selectfont
\setlength{\LTleft}{0pt}\setlength{\LTright}{\fill}
\noindent
\begin{tabular}{@{}>{\raggedright\arraybackslash}p{39.459mm}>{\raggedright\arraybackslash}p{39.459mm}>{\raggedright\arraybackslash}p{39.459mm}@{}}
\toprule
\textbf{Person and number} & \textbf{Present} & \textbf{Preterite} \\
\midrule
1 sg. & ęfthweorfe & ęfthwearf \\
2 sg. & ęfthwierfst & ęfthwurfe \\
3 sg. & ęfthwierfð & ęfthwearf \\
1–3 pl. & ęfthweorfað & ęfthwurfon \\
\bottomrule
\end{tabular}
\par\endgroup
% course-panel:T42-P01:end
% mdtable:T42:end
% teaching: NEW-U13-PRINCIPAL-PARTS-FORCEORFAN
% support: NEW-U13-PRINCIPAL-PARTS-FORCEORFAN
% mdparagraph:P139:begin
\textbf{Principal-part support.} Citation support: forceorfan, forcearf, forcurfon, forcorfen. These four printed parts identify the present, preterite singular, preterite plural and past participle stems. Keep the row visible; only the separately assigned systems and uses are recall targets.
% mdparagraph:P139:end

% mdparagraph:P140:begin
Behealdan and ondrǣdan have the strong VII past stems behēold- and ondrēd-. Bebēodan has the strong II grades bēod-, bēad- and bud-, with bebīetst and bebīett in the present second and third singular. These lexical systems do not make every similarly prefixed verb predictable.
% mdparagraph:P140:end

% teaching: NEW-U13-PAR-BEHEALDAN-IND
% teaching: NEW-U13-GRAM-U13-BEHEALDAN-INDICATIVE
\reserve{7\baselineskip}
% mdparagraph:P141:begin
\subsection*{Table 35. behealdan: present and preterite indicative}
% mdparagraph:P141:end

% table: NEW-U13-PAR-BEHEALDAN-IND
% mdtable:T43:begin
% course-table-style: 2026-10
% course-panel:T43-P01:begin
\par\begingroup
\fontsize{10.5}{13.2}\selectfont
\setlength{\LTleft}{0pt}\setlength{\LTright}{\fill}
\noindent
\begin{tabular}{@{}>{\raggedright\arraybackslash}p{39.459mm}>{\raggedright\arraybackslash}p{39.459mm}>{\raggedright\arraybackslash}p{39.459mm}@{}}
\toprule
\textbf{Person and number} & \textbf{Present} & \textbf{Preterite} \\
\midrule
1 sg. & behealde & behēold \\
2 sg. & behieltst & behēolde \\
3 sg. & behielt & behēold \\
1–3 pl. & behealdað & behēoldon \\
\bottomrule
\end{tabular}
\par\endgroup
% course-panel:T43-P01:end
% mdtable:T43:end
% teaching: NEW-U13-PAR-ONDRAEDAN-IND
% teaching: NEW-U13-GRAM-U13-ONDRAEDAN-INDICATIVE
\reserve{7\baselineskip}
% mdparagraph:P142:begin
\subsection*{Table 36. ondrǣdan: present and preterite indicative}
% mdparagraph:P142:end

% table: NEW-U13-PAR-ONDRAEDAN-IND
% mdtable:T44:begin
% course-table-style: 2026-10
% course-panel:T44-P01:begin
\par\begingroup
\fontsize{10.5}{13.2}\selectfont
\setlength{\LTleft}{0pt}\setlength{\LTright}{\fill}
\noindent
\begin{tabular}{@{}>{\raggedright\arraybackslash}p{39.459mm}>{\raggedright\arraybackslash}p{39.459mm}>{\raggedright\arraybackslash}p{39.459mm}@{}}
\toprule
\textbf{Person and number} & \textbf{Present} & \textbf{Preterite} \\
\midrule
1 sg. & ondrǣde & ondrēd \\
2 sg. & ondrǣtst & ondrēde \\
3 sg. & ondrǣt & ondrēd \\
1–3 pl. & ondrǣdað & ondrēdon \\
\bottomrule
\end{tabular}
\par\endgroup
% course-panel:T44-P01:end
% mdtable:T44:end
% teaching: NEW-U13-PAR-BEBEODAN-IND
% teaching: NEW-U13-GRAM-U13-BEBEODAN-INDICATIVE
\reserve{7\baselineskip}
% mdparagraph:P143:begin
\subsection*{Table 37. bebēodan: present and preterite indicative}
% mdparagraph:P143:end

% table: NEW-U13-PAR-BEBEODAN-IND
% mdtable:T45:begin
% course-table-style: 2026-10
% course-panel:T45-P01:begin
\par\begingroup
\fontsize{10.5}{13.2}\selectfont
\setlength{\LTleft}{0pt}\setlength{\LTright}{\fill}
\noindent
\begin{tabular}{@{}>{\raggedright\arraybackslash}p{39.459mm}>{\raggedright\arraybackslash}p{39.459mm}>{\raggedright\arraybackslash}p{39.459mm}@{}}
\toprule
\textbf{Person and number} & \textbf{Present} & \textbf{Preterite} \\
\midrule
1 sg. & bebēode & bebēad \\
2 sg. & bebīetst & bebude \\
3 sg. & bebīett & bebēad \\
1–3 pl. & bebēodað & bebudon \\
\bottomrule
\end{tabular}
\par\endgroup
% course-panel:T45-P01:end
% mdtable:T45:end
% mdparagraph:P144:begin
Fordōn and ingān use the corresponding finite patterns of dōn and gān with their attested prefixes. The complete ingān subjunctive follows its indicative table. The new productive use is the present subjunctive in the worthy-content construction below; the other subjunctive cells complete the morphological system without adding new clause uses. The reading’s ingangende is separately supplied and is not generated as a regular participle from ingān.
% mdparagraph:P144:end

% teaching: NEW-U13-PAR-FORDON-IND
% teaching: NEW-U13-GRAM-U13-FORDON-INDICATIVE
\reserve{7\baselineskip}
% mdparagraph:P145:begin
\subsection*{Table 38. fordōn: present and preterite indicative}
% mdparagraph:P145:end

% table: NEW-U13-PAR-FORDON-IND
% mdtable:T46:begin
% course-table-style: 2026-10
% course-panel:T46-P01:begin
\par\begingroup
\fontsize{10.5}{13.2}\selectfont
\setlength{\LTleft}{0pt}\setlength{\LTright}{\fill}
\noindent
\begin{tabular}{@{}>{\raggedright\arraybackslash}p{39.459mm}>{\raggedright\arraybackslash}p{39.459mm}>{\raggedright\arraybackslash}p{39.459mm}@{}}
\toprule
\textbf{Person and number} & \textbf{Present} & \textbf{Preterite} \\
\midrule
1 sg. & fordō & fordyde \\
2 sg. & fordēst & fordydest \\
3 sg. & fordēð & fordyde \\
1–3 pl. & fordōð & fordydon \\
\bottomrule
\end{tabular}
\par\endgroup
% course-panel:T46-P01:end
% mdtable:T46:end
% teaching: NEW-U13-PAR-INGAN-IND
% teaching: NEW-U13-GRAM-U13-INGAN-INDICATIVE
\reserve{7\baselineskip}
% mdparagraph:P146:begin
\subsection*{Table 39. ingān: present and preterite indicative}
% mdparagraph:P146:end

% table: NEW-U13-PAR-INGAN-IND
% mdtable:T47:begin
% course-table-style: 2026-10
% course-panel:T47-P01:begin
\par\begingroup
\fontsize{10.5}{13.2}\selectfont
\setlength{\LTleft}{0pt}\setlength{\LTright}{\fill}
\noindent
\begin{tabular}{@{}>{\raggedright\arraybackslash}p{39.459mm}>{\raggedright\arraybackslash}p{39.459mm}>{\raggedright\arraybackslash}p{39.459mm}@{}}
\toprule
\textbf{Person and number} & \textbf{Present} & \textbf{Preterite} \\
\midrule
1 sg. & ingā & inēode \\
2 sg. & ingǣst & inēodest \\
3 sg. & ingǣð & inēode \\
1–3 pl. & ingāð & inēodon \\
\bottomrule
\end{tabular}
\par\endgroup
% course-panel:T47-P01:end
% mdtable:T47:end
% teaching: NEW-U13-PAR-INGAN-SUBJ
% teaching: NEW-U13-GRAM-U13-INGAN-SUBJUNCTIVE
\reserve{7\baselineskip}
% mdparagraph:P147:begin
\subsection*{Table 40. ingān: present and preterite subjunctive}
% mdparagraph:P147:end

% table: NEW-U13-PAR-INGAN-SUBJ
% mdtable:T48:begin
% course-table-style: 2026-10
% course-panel:T48-P01:begin
\par\begingroup
\fontsize{10.5}{13.2}\selectfont
\setlength{\LTleft}{0pt}\setlength{\LTright}{\fill}
\noindent
\begin{tabular}{@{}>{\raggedright\arraybackslash}p{39.459mm}>{\raggedright\arraybackslash}p{39.459mm}>{\raggedright\arraybackslash}p{39.459mm}@{}}
\toprule
\textbf{Person and number} & \textbf{Present} & \textbf{Preterite} \\
\midrule
1 sg. & ingā & inēode \\
2 sg. & ingā & inēode \\
3 sg. & ingā & inēode \\
1–3 pl. & ingān & inēoden \\
\bottomrule
\end{tabular}
\par\endgroup
% course-panel:T48-P01:end
% mdtable:T48:end
\reserve{7\baselineskip}
% mdparagraph:P148:begin
\subsection*{7.3. Commands, infinitives and participle bases}
% mdparagraph:P148:end

% teaching: NEW-U13-VOC-WINNAN
% teaching: NEW-U13-PAR-WINNAN
% mdparagraph:P149:begin
The imperative table gives singular and plural commands for all four verbs. Winnan means “toil”; behealdan “behold”; ondrǣdan “fear”; lǣdan “lead or bring a person.” Prohibitions use the already taught ne with the appropriate command. The earlier begān command pair remains begā, begāð, now with the explicit sense “pursue study.”
% mdparagraph:P149:end

% teaching: NEW-U13-PAR-WINNAN-IMP
% teaching: NEW-U13-PAR-BEHEALDAN-IMP
% teaching: NEW-U13-PAR-ONDRAEDAN-IMP
% teaching: NEW-U13-PAR-LAEDAN-IMP
% teaching: NEW-U13-GRAM-U13-WINNAN-IMPERATIVE
% teaching: NEW-U13-GRAM-U13-BEHEALDAN-IMPERATIVE
% teaching: NEW-U13-GRAM-U13-ONDRAEDAN-IMPERATIVE
% teaching: NEW-U13-GRAM-U13-LAEDAN-IMPERATIVE
\reserve{7\baselineskip}
% mdparagraph:P150:begin
\subsection*{Table 41. Complete singular and plural commands}
% mdparagraph:P150:end

% table: NEW-U13-DISPLAY-IMP
% mdtable:T49:begin
% course-table-style: 2026-10
% course-panel:T49-P01:begin
\par\begingroup
\fontsize{10.5}{13.2}\selectfont
\setlength{\LTleft}{0pt}\setlength{\LTright}{\fill}
\noindent
\begin{tabular}{@{}>{\raggedright\arraybackslash}p{24.232mm}>{\raggedright\arraybackslash}p{22.13mm}>{\raggedright\arraybackslash}p{22.13mm}>{\raggedright\arraybackslash}p{22.13mm}>{\raggedright\arraybackslash}p{22.13mm}@{}}
\toprule
\textbf{Person and number} & \textbf{winnan} & \textbf{behealdan} & \textbf{ondrǣdan} & \textbf{lǣdan} \\
\midrule
2 sg. & win & beheald & ondrǣd & lǣd \\
2 pl. & winnað & behealdað & ondrǣdað & lǣdað \\
\bottomrule
\end{tabular}
\par\endgroup
% course-panel:T49-P01:end
% mdtable:T49:end
% mdparagraph:P151:begin
The plain infinitive follows taught modal and willan constructions. The inflected infinitive adds -ne after the plain infinitive’s final n, producing endings including -anne, -enne, -ianne and, in fordōnne, -ōnne. Use tō with the inflected form in previously taught functions and in the explicitly supported adjective expression below. The table does not introduce every possible infinitive construction.
% mdparagraph:P151:end

% teaching: NEW-U13-PAR-ASCIAN-INF
% teaching: NEW-U13-PAR-FORDON-INF
% teaching: NEW-U13-PAR-LAEDAN-INF
% teaching: NEW-U13-PAR-SPRECAN-INF
% teaching: NEW-U13-GRAM-U13-ASCIAN-INFINITIVE-PAIR
% teaching: NEW-U13-GRAM-U13-FORDON-INFINITIVE-PAIR
% teaching: NEW-U13-GRAM-U13-LAEDAN-INFINITIVE-PAIR
% teaching: NEW-U13-GRAM-U13-SPRECAN-INFINITIVE-PAIR
\reserve{7\baselineskip}
% mdparagraph:P152:begin
\subsection*{Table 42. Complete plain and inflected infinitives}
% mdparagraph:P152:end

% table: NEW-U13-DISPLAY-INF
% mdtable:T50:begin
\par\begingroup
\fontsize{10.5}{13.2}\selectfont
\setlength{\LTleft}{0pt}\setlength{\LTright}{\fill}
\noindent
\begin{tabular}{@{}>{\raggedright\arraybackslash}p{31.82mm}>{\raggedright\arraybackslash}p{40.41mm}>{\raggedright\arraybackslash}p{46.14mm}@{}}
\toprule
\textbf{Verb} & \textbf{Plain infinitive} & \textbf{Inflected infinitive} \\
\midrule
āscian & āscian & āscianne \\
fordōn & fordōn & fordōnne \\
lǣdan & lǣdan & lǣdanne \\
sprecan & sprecan & sprecanne \\
\bottomrule
\end{tabular}
\par\endgroup

% mdtable:T50:end
% mdparagraph:P153:begin
The following table gives complete present and past participle base pairs. Base formation is a production target. In the passive constructions taught here, only the past base is used, uninflected, with wesan; the patient becomes the subject. The table does not introduce full attributive participle declensions or unrestricted uses of every present participle. Gescrȳd has one final d; distinguish the finite present gescrȳt from this past base.
% mdparagraph:P153:end

% teaching: NEW-U13-PAR-FORDILGIAN-PART
% teaching: NEW-U13-PAR-ONLIESAN-PART
% teaching: NEW-U13-PAR-ASETTAN-PART
% teaching: NEW-U13-PAR-GESCRYDAN-PART
% teaching: NEW-U13-PAR-ASENDAN-PART
% teaching: NEW-U13-GRAM-U13-FORDILGIAN-PARTICIPLE-BASES
% teaching: NEW-U13-GRAM-U13-ONLIESAN-PARTICIPLE-BASES
% teaching: NEW-U13-GRAM-U13-ASETTAN-PARTICIPLE-BASES
% teaching: NEW-U13-GRAM-U13-GESCRYDAN-PARTICIPLE-BASES
% teaching: NEW-U13-GRAM-U13-ASENDAN-PARTICIPLE-BASES
\reserve{7\baselineskip}
% mdparagraph:P154:begin
\subsection*{Table 43. Complete present and past participle bases}
% mdparagraph:P154:end

% table: NEW-U13-DISPLAY-PART
% mdtable:T51:begin
\par\begingroup
\fontsize{10.5}{13.2}\selectfont
\setlength{\LTleft}{0pt}\setlength{\LTright}{\fill}
\noindent
\begin{tabular}{@{}>{\raggedright\arraybackslash}p{32.04mm}>{\raggedright\arraybackslash}p{47.09mm}>{\raggedright\arraybackslash}p{39.25mm}@{}}
\toprule
\textbf{Verb} & \textbf{Present participle} & \textbf{Past participle} \\
\midrule
fordilgian & fordilgiende & fordilgod \\
onlīesan & onlīesende & onlīesed \\
āsęttan & āsęttende & āsętt \\
gescrȳdan & gescrȳdende & gescrȳd \\
āsęndan & āsęndende & āsęnd \\
\bottomrule
\end{tabular}
\par\endgroup

% mdtable:T51:end
\reserve{7\baselineskip}
% mdparagraph:P155:begin
\section*{8. Ordinary constructions and smaller words}
% mdparagraph:P155:end

% mdparagraph:P156:begin
Learn a verb with the case or clause it governs. English prepositions are not a reliable guide to the Old English construction. The following frames state the uses required by the exercises; an accusative object remains accusative when it is fronted.
% mdparagraph:P156:end

% teaching: NEW-U13-GRAM-U13-FORDILGIAN-ACC
% teaching: NEW-U13-GRAM-U13-GEWUNIAN-INF
% teaching: NEW-U13-GRAM-U13-ASCIAN-PERSON-THING
% teaching: NEW-U13-GRAM-U13-ASTRECCEAN-REFLEXIVE
% teaching: NEW-U13-GRAM-U13-ONLIESAN-ACC-PASSIVE
% teaching: NEW-U13-GRAM-U13-NEALAECAN-TO-DAT
% teaching: NEW-U13-GRAM-U13-HALSIAN-REQUEST
% teaching: NEW-U13-GRAM-U13-FORCEORFAN-ACC
% teaching: NEW-U13-GRAM-U13-ASETTAN-ACC-PASSIVE
% teaching: NEW-U13-GRAM-U13-WUNIAN-LOCATION
% teaching: NEW-U13-GRAM-U13-OFERWINNAN-ACC
% teaching: NEW-U13-GRAM-U13-GESCRYDAN-CLOTHING
% teaching: NEW-U13-GRAM-U13-ONDRAEDAN-ACC
% teaching: NEW-U13-GRAM-U13-FORDON-ACC
% teaching: NEW-U13-GRAM-U13-INGAN-UNDER
% teaching: NEW-U13-GRAM-U13-ASENDAN-PERSON
% teaching: NEW-U13-GRAM-U13-BEHEALDAN-ACC
% teaching: NEW-U13-GRAM-U13-HATIAN-ACC
% teaching: NEW-U13-GRAM-U13-BEBEODAN-CONTENT
% teaching: NEW-U13-GRAM-U13-EFTHWEORFAN-RETURN
% teaching: NEW-U13-GRAM-U13-LAEDAN-PERSON-DIRECTION
% teaching: NEW-U13-GRAM-U13-SPRECAN-CONVERSATION-MID
% teaching: NEW-U13-GRAM-U13-BEGAN-STUDY
% teaching: NEW-U13-GRAM-U13-ASTIGAN-DESCENT
% teaching: NEW-U13-GRAM-U13-SELLAN-SALE
% teaching: NEW-U13-GRAM-U13-WYRCEAN-MAKING
% teaching: NEW-U13-GRAM-U13-INSTRUMENTAL-ROUTE
% teaching: NEW-U13-GRAM-U13-PLEASANT-INFINITIVE
% teaching: NEW-U13-GRAM-U13-GE-PAIRED
% teaching: NEW-U13-GRAM-U13-BEFORAN-POSITIONING-ACC
% teaching: NEW-U13-GRAM-U13-WITH-PROXIMITY-ACC
% teaching: NEW-U13-GRAM-U13-UNDER-DIRECTION-ACC
% teaching: NEW-U13-GRAM-U13-WORTHY-CONTENT
% teaching: NEW-U13-GRAM-U13-HWILC-INDEFINITE
% teaching: NEW-U13-GRAM-U13-DON-WELL-DATIVE
% teaching: NEW-U13-GRAM-U13-SINGAN-SONG
\reserve{7\baselineskip}
% mdparagraph:P157:begin
\subsection*{8.1. New finite verbs}
% mdparagraph:P157:end

% mdtable:T52:begin
\par\begingroup
\fontsize{10.5}{13.2}\selectfont
\setlength{\LTleft}{0pt}\setlength{\LTright}{\fill}
\begin{longtable}{@{}>{\raggedright\arraybackslash}p{21.67mm}>{\raggedright\arraybackslash}p{99.52mm}@{}}
\toprule
\textbf{Word or expression} & \textbf{Construction} \\
\midrule
\endfirsthead
\toprule
\textbf{Word or expression} & \textbf{Construction} \\
\midrule
\endhead
\midrule
\multicolumn{2}{r}{\footnotesize Continued on the following page.}\\
\endfoot
\bottomrule
\endlastfoot
fordilgian & Destroy an accusative thing. The corresponding passive uses the uninflected past participle with wesan. \\
gewunian & Be accustomed to a plain infinitive. The subject of the finite verb also performs the infinitive’s action. \\
āscian & Ask an accusative person about a genitive thing. A supported indirect question can give the content; do not use two accusatives for person and thing. \\
āstręccean & Prostrate oneself: the reflexive pronoun is accusative and refers to the subject. \\
onlīesan & Release an accusative person or thing. In the passive, use the uninflected past base onlīesed with wesan. \\
nēalǣcan & Approach a goal expressed by tō with the dative. \\
hālsian & Entreat an accusative person. A requested action may be expressed by ðæt with the subjunctive. \\
forceorfan & Cut down an accusative tree or other object. \\
āsęttan & Set or place an accusative object. The passive uses the uninflected past base with wesan; the placement phrase has its own case requirements. \\
wunian & Dwell or remain, either without a complement or with an ordinary static location. \\
oferwinnan & Overcome an accusative person or force. \\
gescrȳdan; mid & Clothe an accusative person; mid with the dative names the garment. The passive retains mid with the dative. English “clothed in” does not require Old English on. \\
ondrǣdan & Fear an accusative person or thing. A dative pronoun referring to the subject is possible but need not be added to every sentence. \\
fordōn & Destroy accusative persons or things. Its plain infinitive can follow a taught modal. \\
ingān & Enter or go in. In the supported roof-entry construction, under governs an accusative noun phrase. \\
āsęndan & Send an accusative person or thing. The corresponding passive uses the uninflected past base with wesan. \\
behealdan & Behold or look at an accusative object. \\
hatian & Hate an accusative person. \\
bebēodan & Command a dative addressee. A ðæt clause may express the command with the subjunctive; a relative can identify what is commanded. \\
ęfthweorfan & Return or turn back, without a direct object. A route expression can specify the way. \\
\end{longtable}
\par\endgroup

% mdtable:T52:end
\reserve{7\baselineskip}
% mdparagraph:P158:begin
\subsection*{8.2. Additional senses of familiar verbs}
% mdparagraph:P158:end

% mdtable:T53:begin
\par\begingroup
\fontsize{10.5}{13.2}\selectfont
\setlength{\LTleft}{0pt}\setlength{\LTright}{\fill}
\noindent
\begin{tabular}{@{}>{\raggedright\arraybackslash}p{21.53mm}>{\raggedright\arraybackslash}p{99.66mm}@{}}
\toprule
\textbf{Word or expression} & \textbf{Construction} \\
\midrule
lǣdan & Lead, bring or take an accusative person. Mid with the dative can name a companion, and hider can give the direction “here.” \\
sprecan; mid & Speak with someone: mid with the dative names the interlocutor. \\
begān & Pursue or practise an accusative activity, including leornung “study.” Use the finite and command forms already taught. \\
āstīgan; of & Descend; of with the dative marks the source. This is an additional meaning of the previously encountered verb. \\
sęllan & Sell an accusative possession. This sense is distinct from the earlier meaning “give.” \\
wyrcean & Make an accusative thing. Its earlier “work; perform” sense also permits an accusative deed, including abstract wrongdoing. \\
\bottomrule
\end{tabular}
\par\endgroup

% mdtable:T53:end
\reserve{7\baselineskip}
% mdparagraph:P159:begin
\subsection*{8.3. Other required constructions}
% mdparagraph:P159:end

% mdtable:T54:begin
\par\begingroup
\fontsize{10.5}{13.2}\selectfont
\setlength{\LTleft}{0pt}\setlength{\LTright}{\fill}
\noindent
\begin{tabular}{@{}>{\raggedright\arraybackslash}p{20.97mm}>{\raggedright\arraybackslash}p{100.22mm}@{}}
\toprule
\textbf{Word or expression} & \textbf{Construction} \\
\midrule
ðȳ ilcan wege & “By the same way”: ðȳ is instrumental; ilcan and wege have the already familiar dative shapes in this route expression. Learn the whole expression; it does not add a new ilca declension. \\
wynsum tō hīeranne & “Pleasant to hear”: tō takes the inflected infinitive. The adjective agrees with what is pleasant. This expression supplies the evaluative use needed here. \\
ge … ge & “Both … and”: the paired conjunction joins parallel constituents. Both occurrences have short e. \\
beforan & Before a person, with the accusative in a movement or positioning construction. A static dative is also a valid construal; follow the requested exercise target where one is given. \\
wið & Beside, with the accusative, including static standing beside someone. This accusative does not itself express motion. \\
under & Under, with the accusative for entering beneath a roof. Keep this distinct from the earlier static under with the dative. \\
wierðe & Worthy, with a ðæt clause and present subjunctive for the proposed entry. This bounded construction uses ingān; it does not introduce a general rule for every adjective. \\
hwilc & Any or some, as an indefinite adjective. This use is distinct from the interrogative “which?”; agreement follows the noun qualified. \\
dōn wel & Do well to a person in the dative, without an expressed object naming a deed. \\
singan & Sing an accusative song: sang can be the direct object although it names the action’s own product. \\
\bottomrule
\end{tabular}
\par\endgroup

% mdtable:T54:end
% mdparagraph:P160:begin
In familiar verbs, the earlier inflection remains available, but the meanings just specified must be learned separately. Thus knowing sęllan “give” did not by itself establish “sell,” nor did knowing mid for accompaniment establish its conversational or clothing use.
% mdparagraph:P160:end

% teaching: NEW-U13-VOC-THAER
% teaching: NEW-U13-PAR-THAER
% teaching: NEW-U13-VOC-EFT
% teaching: NEW-U13-PAR-EFT
% teaching: NEW-U13-VOC-OF
% teaching: NEW-U13-PAR-OF
% teaching: NEW-U13-VOC-BUTAN
% teaching: NEW-U13-PAR-BUTAN
% teaching: NEW-U13-VOC-OTH-PREPOSITION
% teaching: NEW-U13-PAR-OTH-PREPOSITION
% teaching: NEW-U13-VOC-THIDER
% teaching: NEW-U13-PAR-THIDER
% teaching: NEW-U13-VOC-AER-ADVERB
% teaching: NEW-U13-PAR-AER-ADVERB
% teaching: NEW-U13-VOC-STYCCEMAELUM
% teaching: NEW-U13-PAR-STYCCEMAELUM
% teaching: NEW-U13-VOC-NEAH-ADVERB
% teaching: NEW-U13-PAR-NEAH-ADVERB
% teaching: NEW-U13-VOC-OTHOTHE
% teaching: NEW-U13-PAR-OTHOTHE
% teaching: NEW-U13-VOC-GE-CONJUNCTION
% teaching: NEW-U13-PAR-GE-CONJUNCTION
% teaching: NEW-U13-VOC-BEFORAN
% teaching: NEW-U13-VOC-WITH
% teaching: NEW-U13-VOC-UNDER
% teaching: NEW-U13-VOC-MID
% teaching: NEW-U13-PAR-MID
% teaching: NEW-U13-VOC-HWILC
% mdtable:T55:begin
\par\begingroup
\fontsize{10.5}{13.2}\selectfont
\setlength{\LTleft}{0pt}\setlength{\LTright}{\fill}
\noindent
\begin{tabular}{@{}>{\raggedright\arraybackslash}p{41.92mm}>{\raggedright\arraybackslash}p{79.27mm}@{}}
\toprule
\textbf{Word} & \textbf{Meaning and use} \\
\midrule
ðǣr & there \\
ęft & again; back \\
of & out of; from, governing dative \\
būtan & without, governing dative \\
oð & up to; until, governing accusative \\
ðider & thither; to that place \\
ǣr & before; earlier \\
ǣrest & first; at first \\
styccemǣlum & little by little \\
nēah & near; at hand \\
oððe & or \\
ge & both ... and (paired ge ... ge) \\
\bottomrule
\end{tabular}
\par\endgroup

% mdtable:T55:end
% mdparagraph:P161:begin
Distinguish ðǣr “there,” a location, from ðider “to that place,” a direction. Here ęft has the explicitly taught sense “again; back”; the reading also supplies its contextual translation “afterwards.” Nēah in the ordinary expression “be near” or “be at hand” is an adverb. Oð with the accusative marks a limit; the reading separately supplies the conjunction oð-ðæt “until.”
% mdparagraph:P161:end

\reserve{7\baselineskip}
% mdparagraph:P162:begin
\section*{9. New practice}
% mdparagraph:P162:end

% mdparagraph:P163:begin
For analysis tasks, preserve all possible references when the context does not decide between them. For composition, meet the stated target while keeping the intended participants, agreement and meaning. A valid alternative order and compliance with a particular requested order are separate questions.
% mdparagraph:P163:end

\reserve{7\baselineskip}
% mdparagraph:P164:begin
\subsection*{Forms}
% mdparagraph:P164:end

\par\begin{minipage}{\linewidth}\raggedright
% mdparagraph:P165:begin
\subsection*{F01}
% mdparagraph:P165:end

% item: NEW-U13-F01
% mdparagraph:P166:begin
Give the complete nominative, genitive, dative and accusative singular and plural of mōdor, dohtor and brōðor. Explain the selected oblique and plural stems.
% mdparagraph:P166:end

\end{minipage}\par\addvspace{4pt}

\par\begin{minipage}{\linewidth}\raggedright
% mdparagraph:P167:begin
\subsection*{F02}
% mdparagraph:P167:end

% item: NEW-U13-F02
% mdparagraph:P168:begin
Give the complete nominative, genitive, dative and accusative singular and plural of æx and ēagðȳrel. Explain the retained stem in æx and the syncopated stem before endings in ēagðȳrel.
% mdparagraph:P168:end

\end{minipage}\par\addvspace{4pt}

\par\begin{minipage}{\linewidth}\raggedright
% mdparagraph:P169:begin
\subsection*{F03}
% mdparagraph:P169:end

% item: NEW-U13-F03
% mdparagraph:P170:begin
Give the genitive and dative singular and the nominative and genitive plural of cirice, mynster and wyrttruma. State the model used for each.
% mdparagraph:P170:end

\end{minipage}\par\addvspace{4pt}

\par\begin{minipage}{\linewidth}\raggedright
% mdparagraph:P171:begin
\subsection*{F04}
% mdparagraph:P171:end

% item: NEW-U13-F04
% mdparagraph:P172:begin
For swēte, give all strong nominative, genitive, dative and accusative singular and plural forms in all three genders. Then give the weak masculine nominative singular of its comparative and both strong and weak masculine nominative singular of its superlative. Explain the stems and endings.
% mdparagraph:P172:end

\end{minipage}\par\addvspace{4pt}

\par\begin{minipage}{\linewidth}\raggedright
% mdparagraph:P173:begin
\subsection*{F05}
% mdparagraph:P173:end

% item: NEW-U13-F05
% mdparagraph:P174:begin
Give the complete present and preterite indicative of forceorfan: all three persons in singular and plural. Using the printed principal parts, explain the present singular changes and the distribution of preterite stems.
% mdparagraph:P174:end

% support: NEW-U13-PRINCIPAL-PARTS-FORCEORFAN
% mdparagraph:P175:begin
\textbf{Supplied help.} Citation support: forceorfan, forcearf, forcurfon, forcorfen. These four printed parts identify the present, preterite singular, preterite plural and past participle stems. Keep the row visible; only the separately assigned systems and uses are recall targets.
% mdparagraph:P175:end

\end{minipage}\par\addvspace{4pt}

\par\begin{minipage}{\linewidth}\raggedright
% mdparagraph:P176:begin
\subsection*{F06}
% mdparagraph:P176:end

% item: NEW-U13-F06
% mdparagraph:P177:begin
Give the singular and plural commands of winnan, begān and behealdan. For fordilgian and onlīesan, give the base present and past participles.
% mdparagraph:P177:end

\end{minipage}\par\addvspace{4pt}

\reserve{7\baselineskip}
% mdparagraph:P178:begin
\subsection*{Clause analysis and reference}
% mdparagraph:P178:end

\par\begin{minipage}{\linewidth}\raggedright
% mdparagraph:P179:begin
\subsection*{A01}
% mdparagraph:P179:end

% item: NEW-U13-A01
% mdparagraph:P180:begin
Translate both clauses. Identify the subject, direct object and finite verb in each; explain which forms settle the participants.
% mdparagraph:P180:end

% mdparagraph:P181:begin
\textbf{Se bisceop grētte ðone munuc.}
% mdparagraph:P181:end

% mdparagraph:P182:begin
\textbf{Ðone munuc grētte se bisceop.}
% mdparagraph:P182:end

\end{minipage}\par\addvspace{4pt}

\par\begin{minipage}{\linewidth}\raggedright
% mdparagraph:P183:begin
\subsection*{A02}
% mdparagraph:P183:end

% item: NEW-U13-A02
% mdparagraph:P184:begin
Divide this sentence into clauses; identify each finite verb and its subject. Label every ðā, and translate the whole sentence.
% mdparagraph:P184:end

% mdparagraph:P185:begin
\textbf{Ðā se munuc cōm, ðā grētte se bisceop ðā brōðor.}
% mdparagraph:P185:end

\end{minipage}\par\addvspace{4pt}

\par\begin{minipage}{\linewidth}\raggedright
% mdparagraph:P186:begin
\subsection*{A03}
% mdparagraph:P186:end

% item: NEW-U13-A03
% mdparagraph:P187:begin
Translate. Identify the main finite verb and the relative-clause finite verb; give the relative antecedent and its role in each clause.
% mdparagraph:P187:end

% mdparagraph:P188:begin
\textbf{Se bisceop, ðe se munuc grētte, sæt on ðǣm hūse.}
% mdparagraph:P188:end

\end{minipage}\par\addvspace{4pt}

\par\begin{minipage}{\linewidth}\raggedright
% mdparagraph:P189:begin
\subsection*{A04}
% mdparagraph:P189:end

% item: NEW-U13-A04
% mdparagraph:P190:begin
Translate. Identify the person greeted, the person who gave the food, and the recipient. Explain the understood subject and the pronoun cases.
% mdparagraph:P190:end

% mdparagraph:P191:begin
\textbf{Se munuc grētte ðone bisceop and sealde him ðone męte.}
% mdparagraph:P191:end

\end{minipage}\par\addvspace{4pt}

\par\begin{minipage}{\linewidth}\raggedright
% mdparagraph:P192:begin
\subsection*{A05}
% mdparagraph:P192:end

% item: NEW-U13-A05
% mdparagraph:P193:begin
Translate the second sentence after each context. State the referent of hē in each; explain what morphology establishes and what the context establishes.
% mdparagraph:P193:end

% mdparagraph:P194:begin
\textbf{Se sang wæs swēte. Hē wæs wynsum.}
% mdparagraph:P194:end

% mdparagraph:P195:begin
\textbf{Se bisceop cōm. Hē wæs on ðǣm hūse.}
% mdparagraph:P195:end

\end{minipage}\par\addvspace{4pt}

\par\begin{minipage}{\linewidth}\raggedright
% mdparagraph:P196:begin
\subsection*{A06}
% mdparagraph:P196:end

% item: NEW-U13-A06
% mdparagraph:P197:begin
Give both possible referents of hē in the second sentence. Is either one selected by gender and number alone? Rewrite that sentence twice, naming each possible subject explicitly.
% mdparagraph:P197:end

% mdparagraph:P198:begin
\textbf{Se bisceop geseah ðone munuc. Hē sæt on ðǣm hūse.}
% mdparagraph:P198:end

\end{minipage}\par\addvspace{4pt}

\par\begin{minipage}{\linewidth}\raggedright
% mdparagraph:P199:begin
\subsection*{A07}
% mdparagraph:P199:end

% item: NEW-U13-A07
% mdparagraph:P200:begin
The monk addresses the bishop. Translate, and identify each speaker/\allowbreak{}addressee reference in the quotation.
% mdparagraph:P200:end

% mdparagraph:P201:begin
\textbf{Se munuc cwæð: “Ic hīerde ðone sang. Hīerdest ðū hine?”}
% mdparagraph:P201:end

\end{minipage}\par\addvspace{4pt}

\par\begin{minipage}{\linewidth}\raggedright
% mdparagraph:P202:begin
\subsection*{A08}
% mdparagraph:P202:end

% item: NEW-U13-A08
% mdparagraph:P203:begin
Translate. Identify the content-clause boundary, its finite verb, and the antecedent of hē. The monk is reporting his own presence.
% mdparagraph:P203:end

% mdparagraph:P204:begin
\textbf{Se munuc sægde ðæt hē on ðǣm hūse wǣre.}
% mdparagraph:P204:end

\end{minipage}\par\addvspace{4pt}

\par\begin{minipage}{\linewidth}\raggedright
% mdparagraph:P205:begin
\subsection*{A09}
% mdparagraph:P205:end

% item: NEW-U13-A09
% mdparagraph:P206:begin
Translate, restoring the omitted repeated predicate in English. Then write the second clause in Old English with that predicate expressed and agreeing with its own subject.
% mdparagraph:P206:end

% mdparagraph:P207:begin
\textbf{Se bisceop wæs on ðǣre cirican, and ðā brōðor on ðǣm hūse.}
% mdparagraph:P207:end

\end{minipage}\par\addvspace{4pt}

\par\begin{minipage}{\linewidth}\raggedright
% mdparagraph:P208:begin
\subsection*{A10}
% mdparagraph:P208:end

% item: NEW-U13-A10
% mdparagraph:P209:begin
Translate both sentences. Explain which clause each initial phrase belongs to and identify the finite verbs. Are initial words themselves a reliable way to decide who performs an action?
% mdparagraph:P209:end

% mdparagraph:P210:begin
\textbf{On ðǣm hūse sæt se munuc.}
% mdparagraph:P210:end

% mdparagraph:P211:begin
\textbf{Ðā se munuc on ðǣm hūse sæt, ðā cōm se bisceop.}
% mdparagraph:P211:end

\end{minipage}\par\addvspace{4pt}

\par\begin{minipage}{\linewidth}\raggedright
% mdparagraph:P212:begin
\subsection*{A11}
% mdparagraph:P212:end

% item: NEW-U13-A11
% mdparagraph:P213:begin
Translate and draw the clause structure, using brackets or indentation. Identify the subject of every finite verb and the referent/\allowbreak{}function of ðe and ðā.
% mdparagraph:P213:end

% mdparagraph:P214:begin
\textbf{Ðā se munuc, ðe se bisceop grētte, cōm, ðā sungon ðā brōðor.}
% mdparagraph:P214:end

\end{minipage}\par\addvspace{4pt}

\par\begin{minipage}{\linewidth}\raggedright
% mdparagraph:P215:begin
\subsection*{A12}
% mdparagraph:P215:end

% item: NEW-U13-A12
% mdparagraph:P216:begin
A learner says: “Every main clause has the finite verb as its second word, and every dependent clause puts it last.” Explain two problems with this statement, referring to the chapter’s checked examples. Then explain why Cook 25.7 can nevertheless require a final verb as a particular exercise target.
% mdparagraph:P216:end

% mdparagraph:P217:begin
The two source fragments and their complete supplied order help are in §2. Keep that help visible for this analysis.
% mdparagraph:P217:end

\end{minipage}\par\addvspace{4pt}

\reserve{7\baselineskip}
% mdparagraph:P218:begin
\subsection*{Composition}
% mdparagraph:P218:end

\par\begin{minipage}{\linewidth}\raggedright
% mdparagraph:P219:begin
\subsection*{C01}
% mdparagraph:P219:end

% item: NEW-U13-C01
% mdparagraph:P220:begin
Translate “The brothers heard the song” twice: first with the subject phrase first, then with the object phrase first. Keep the same participants.
% mdparagraph:P220:end

\end{minipage}\par\addvspace{4pt}

\par\begin{minipage}{\linewidth}\raggedright
% mdparagraph:P221:begin
\subsection*{C02}
% mdparagraph:P221:end

% item: NEW-U13-C02
% mdparagraph:P222:begin
Translate: Then the bishop greeted the brothers. Begin with ðā and place the finite verb before the nominal subject.
% mdparagraph:P222:end

\end{minipage}\par\addvspace{4pt}

\par\begin{minipage}{\linewidth}\raggedright
% mdparagraph:P223:begin
\subsection*{C03}
% mdparagraph:P223:end

% item: NEW-U13-C03
% mdparagraph:P224:begin
Translate: When the brothers came, the bishop sang. Use a ðā temporal clause and a main-clause correlative ðā; use a late verb in the dependent clause.
% mdparagraph:P224:end

\end{minipage}\par\addvspace{4pt}

\par\begin{minipage}{\linewidth}\raggedright
% mdparagraph:P225:begin
\subsection*{C04}
% mdparagraph:P225:end

% item: NEW-U13-C04
% mdparagraph:P226:begin
Translate: The brothers whom the bishop greeted went to the church. Use invariant ðe and keep the relative beside its antecedent.
% mdparagraph:P226:end

\end{minipage}\par\addvspace{4pt}

\par\begin{minipage}{\linewidth}\raggedright
% mdparagraph:P227:begin
\subsection*{C05}
% mdparagraph:P227:end

% item: NEW-U13-C05
% mdparagraph:P228:begin
Translate as two sentences: The bishop greeted the monk. The monk then went to the church. Repeat the second subject as a noun phrase so that the reference is explicit.
% mdparagraph:P228:end

\end{minipage}\par\addvspace{4pt}

\par\begin{minipage}{\linewidth}\raggedright
% mdparagraph:P229:begin
\subsection*{C06}
% mdparagraph:P229:end

% item: NEW-U13-C06
% mdparagraph:P230:begin
The monk addresses several brothers. Translate: The monk said, “I want to speak with you.” Preserve the person and number appropriate to direct speech.
% mdparagraph:P230:end

\end{minipage}\par\addvspace{4pt}

\par\begin{minipage}{\linewidth}\raggedright
% mdparagraph:P231:begin
\subsection*{C07}
% mdparagraph:P231:end

% item: NEW-U13-C07
% mdparagraph:P232:begin
Translate: The monk saw the bishop and greeted him. Express the subject once; both actions have that subject.
% mdparagraph:P232:end

\end{minipage}\par\addvspace{4pt}

\par\begin{minipage}{\linewidth}\raggedright
% mdparagraph:P233:begin
\subsection*{C08}
% mdparagraph:P233:end

% item: NEW-U13-C08
% mdparagraph:P234:begin
Translate: When the bishop whom the monk greeted came, the brothers sang. Use ðe, temporal ðā and a correlative main-clause ðā. Mark the three clauses after writing the sentence.
% mdparagraph:P234:end

\end{minipage}\par\addvspace{4pt}

\reserve{7\baselineskip}
% mdparagraph:P235:begin
\section*{10. Cook VII.22–25}
% mdparagraph:P235:end

% mdparagraph:P236:begin
Translate each English sentence into Old English. The wording, item numbers and original grammatical controls are retained. Parenthesized English and Old English belong to Cook’s prompt. Keep the local helps visible; supplied forms and constructions are not independent recall targets. Original section references are explained in the editorial notes.
% mdparagraph:P236:end

% teaching: NEW-U13-COOK-RELATIVE-CONTROLS
% mdparagraph:P237:begin
Cook’s 87.b requests an inflected demonstrative followed by ðe. Its gender and number follow the antecedent, while its case follows its role inside the relative clause; use this strategy in 24.9, 25.1 and 25.6. Cook’s 87.d requests invariant ðe alone; use it in 25.7. These are the relative strategies taught in Unit 8. A different grammatical strategy can preserve the meaning while missing the stated exercise target.
% mdparagraph:P237:end

\reserve{7\baselineskip}
% mdparagraph:P238:begin
\subsection*{Cook VII.22}
% mdparagraph:P238:end

% mdparagraph:P239:begin
Items 1–15. Use the ordinary senses and verb constructions in §8, including the explicit new senses of familiar words.
% mdparagraph:P239:end

\par\begin{minipage}{\linewidth}\raggedright
% mdparagraph:P240:begin
\subsection*{Cook VII.22.1}
% mdparagraph:P240:end

% item: EX-VII-22-01
% mdparagraph:P241:begin
Idleness is the mother of sin.
% mdparagraph:P241:end

\end{minipage}\par\addvspace{4pt}

\par\begin{minipage}{\linewidth}\raggedright
% mdparagraph:P242:begin
\subsection*{Cook VII.22.2}
% mdparagraph:P242:end

% item: EX-VII-22-02
% mdparagraph:P243:begin
Etheldred was (om. the) daughter of the king of the East Anglians.
% mdparagraph:P243:end

\end{minipage}\par\addvspace{4pt}

% teaching: U13-COOK-HELP-NAMES
% support: U13-COOK-HELP-NAMES
% mdparagraph:P244:begin
\textbf{Supplied help.} Supply Æðelðrȳð as the nominative personal name and Eastęngla as the genitive plural “of the East Anglians.” No productive proper-name system is assigned.
% mdparagraph:P244:end

% teaching: NEW-U13-VOC-AETHELTHRYTH
% teaching: NEW-U13-VOC-EASTENGLE
% teaching: NEW-U13-LOCAL-HELPS
% teaching: NEW-U13-GRAM-U13-COOK-LOCAL-NAMES
\par\begin{minipage}{\linewidth}\raggedright
% mdparagraph:P245:begin
\subsection*{Cook VII.22.3}
% mdparagraph:P245:end

% item: EX-VII-22-03
% mdparagraph:P246:begin
The bishop’s beautiful monastery was destroyed (fordilgian).
% mdparagraph:P246:end

\end{minipage}\par\addvspace{4pt}

\par\begin{minipage}{\linewidth}\raggedright
% mdparagraph:P247:begin
\subsection*{Cook VII.22.4}
% mdparagraph:P247:end

% item: EX-VII-22-04
% mdparagraph:P248:begin
Labor [toil, sing.] with your hands.
% mdparagraph:P248:end

\end{minipage}\par\addvspace{4pt}

\par\begin{minipage}{\linewidth}\raggedright
% mdparagraph:P249:begin
\subsection*{Cook VII.22.5}
% mdparagraph:P249:end

% item: EX-VII-22-05
% mdparagraph:P250:begin
Pursue (sing.) your study.
% mdparagraph:P250:end

\end{minipage}\par\addvspace{4pt}

\par\begin{minipage}{\linewidth}\raggedright
% mdparagraph:P251:begin
\subsection*{Cook VII.22.6}
% mdparagraph:P251:end

% item: EX-VII-22-06
% mdparagraph:P252:begin
He goes to church.
% mdparagraph:P252:end

\end{minipage}\par\addvspace{4pt}

\par\begin{minipage}{\linewidth}\raggedright
% mdparagraph:P253:begin
\subsection*{Cook VII.22.7}
% mdparagraph:P253:end

% item: EX-VII-22-07
% mdparagraph:P254:begin
Depart from me, ye that work unrighteousness.
% mdparagraph:P254:end

\end{minipage}\par\addvspace{4pt}

\par\begin{minipage}{\linewidth}\raggedright
% mdparagraph:P255:begin
\subsection*{Cook VII.22.8}
% mdparagraph:P255:end

% item: EX-VII-22-08
% mdparagraph:P256:begin
What is your name?
% mdparagraph:P256:end

\end{minipage}\par\addvspace{4pt}

\par\begin{minipage}{\linewidth}\raggedright
% mdparagraph:P257:begin
\subsection*{Cook VII.22.9}
% mdparagraph:P257:end

% item: EX-VII-22-09
% mdparagraph:P258:begin
Tell (hātan) our three friends to go thither to him.
% mdparagraph:P258:end

\end{minipage}\par\addvspace{4pt}

\par\begin{minipage}{\linewidth}\raggedright
% mdparagraph:P259:begin
\subsection*{Cook VII.22.10}
% mdparagraph:P259:end

% item: EX-VII-22-10
% mdparagraph:P260:begin
The song descends from heaven (plur.) to earth.
% mdparagraph:P260:end

\end{minipage}\par\addvspace{4pt}

\par\begin{minipage}{\linewidth}\raggedright
% mdparagraph:P261:begin
\subsection*{Cook VII.22.11}
% mdparagraph:P261:end

% item: EX-VII-22-11
% mdparagraph:P262:begin
I am wont to visit the bishop.
% mdparagraph:P262:end

\end{minipage}\par\addvspace{4pt}

\par\begin{minipage}{\linewidth}\raggedright
% mdparagraph:P263:begin
\subsection*{Cook VII.22.12}
% mdparagraph:P263:end

% item: EX-VII-22-12
% mdparagraph:P264:begin
I wish (willan) first (om. to) ask you (sing.) one thing (156).
% mdparagraph:P264:end

\end{minipage}\par\addvspace{4pt}

\par\begin{minipage}{\linewidth}\raggedright
% mdparagraph:P265:begin
\subsection*{Cook VII.22.13}
% mdparagraph:P265:end

% item: EX-VII-22-13
% mdparagraph:P266:begin
I wish to take (lǣdan) him with me.
% mdparagraph:P266:end

\end{minipage}\par\addvspace{4pt}

\par\begin{minipage}{\linewidth}\raggedright
% mdparagraph:P267:begin
\subsection*{Cook VII.22.14}
% mdparagraph:P267:end

% item: EX-VII-22-14
% mdparagraph:P268:begin
The priest prostrated himself before (beforan, 172) the bishop.
% mdparagraph:P268:end

\end{minipage}\par\addvspace{4pt}

\par\begin{minipage}{\linewidth}\raggedright
% mdparagraph:P269:begin
\subsection*{Cook VII.22.15}
% mdparagraph:P269:end

% item: EX-VII-22-15
% mdparagraph:P270:begin
On the third day the holy soul was released from the body.
% mdparagraph:P270:end

\end{minipage}\par\addvspace{4pt}

\reserve{7\baselineskip}
% mdparagraph:P271:begin
\subsection*{Cook VII.23}
% mdparagraph:P271:end

% mdparagraph:P272:begin
Items 1–13. The supplied placement, taking and wonder constructions are local to their items.
% mdparagraph:P272:end

\par\begin{minipage}{\linewidth}\raggedright
% mdparagraph:P273:begin
\subsection*{Cook VII.23.1}
% mdparagraph:P273:end

% item: EX-VII-23-01
% mdparagraph:P274:begin
The largest window of the church is two feet (154. c) broad, and four feet high.
% mdparagraph:P274:end

\end{minipage}\par\addvspace{4pt}

\par\begin{minipage}{\linewidth}\raggedright
% mdparagraph:P275:begin
\subsection*{Cook VII.23.2}
% mdparagraph:P275:end

% item: EX-VII-23-02
% mdparagraph:P276:begin
Little by little the sweet song drew nigh.
% mdparagraph:P276:end

\end{minipage}\par\addvspace{4pt}

\par\begin{minipage}{\linewidth}\raggedright
% mdparagraph:P277:begin
\subsection*{Cook VII.23.3}
% mdparagraph:P277:end

% item: EX-VII-23-03
% mdparagraph:P278:begin
I entreat you to tell [that you tell] me what is (194. b) the angel’s name.
% mdparagraph:P278:end

\end{minipage}\par\addvspace{4pt}

\par\begin{minipage}{\linewidth}\raggedright
% mdparagraph:P279:begin
\subsection*{Cook VII.23.4}
% mdparagraph:P279:end

% item: EX-VII-23-04
% mdparagraph:P280:begin
The soldiers took his raiment, and made (wyrcean) four parts, (om. to) every soldier a part; thus did the soldiers.
% mdparagraph:P280:end

\end{minipage}\par\addvspace{4pt}

\par\begin{minipage}{\linewidth}\raggedright
% mdparagraph:P281:begin
\subsection*{Cook VII.23.5}
% mdparagraph:P281:end

% item: EX-VII-23-05
% mdparagraph:P282:begin
The mother of Jesus stood by (wið, 172) the cross.
% mdparagraph:P282:end

\end{minipage}\par\addvspace{4pt}

\par\begin{minipage}{\linewidth}\raggedright
% mdparagraph:P283:begin
\subsection*{Cook VII.23.6}
% mdparagraph:P283:end

% item: EX-VII-23-06
% mdparagraph:P284:begin
Many read this writing.
% mdparagraph:P284:end

\end{minipage}\par\addvspace{4pt}

\par\begin{minipage}{\linewidth}\raggedright
% mdparagraph:P285:begin
\subsection*{Cook VII.23.7}
% mdparagraph:P285:end

% item: EX-VII-23-07
% mdparagraph:P286:begin
With his little hatchet he cut down (forceorfan, 134) the tree.
% mdparagraph:P286:end

\end{minipage}\par\addvspace{4pt}

\par\begin{minipage}{\linewidth}\raggedright
% mdparagraph:P287:begin
\subsection*{Cook VII.23.8}
% mdparagraph:P287:end

% item: EX-VII-23-08
% mdparagraph:P288:begin
The ax is placed (āsęttan) into (tō) the root (wyrttruma, 53) of the trees.
% mdparagraph:P288:end

\end{minipage}\par\addvspace{4pt}

% teaching: U13-COOK-HELP-GENITIVE-VARIANT
% support: U13-COOK-HELP-GENITIVE-VARIANT
% mdparagraph:P289:begin
\textbf{Supplied help.} Ðǣra is the supplied genitive plural variant corresponding to ordinary ðāra. It agrees with the plural trees.
% mdparagraph:P289:end

% teaching: NEW-U13-GRAM-U13-COOK-LOCAL-GENITIVE-VARIANT
% teaching: U13-COOK-HELP-ROOT
% support: U13-COOK-HELP-ROOT
% mdparagraph:P290:begin
\textbf{Supplied help.} For the local placement phrase, use tō with dative wyrttruman and a genitive plural trees phrase. The phrase denotes placement at/\allowbreak{}to the root; it does not create a general static tō rule.
% mdparagraph:P290:end

% teaching: NEW-U13-VOC-TO
% teaching: NEW-U13-PAR-TO
% teaching: NEW-U13-GRAM-u13-to-at-root-support
% teaching: NEW-U13-GRAM-U13-COOK-LOCAL-ROOT
\par\begin{minipage}{\linewidth}\raggedright
% mdparagraph:P291:begin
\subsection*{Cook VII.23.9}
% mdparagraph:P291:end

% item: EX-VII-23-09
% mdparagraph:P292:begin
He took an ax in his hand.
% mdparagraph:P292:end

\end{minipage}\par\addvspace{4pt}

% teaching: U13-COOK-HELP-TAKE-IN-HAND
% support: U13-COOK-HELP-TAKE-IN-HAND
% mdparagraph:P293:begin
\textbf{Supplied help.} For this taking action, the bounded phrase on his hand means “into his hand,” with accusative hand. This is a declared composition applying Cook’s qualified directional use of on and the ordinary sense of niman; it is not an attested Old English quotation or a general productive niman-plus-goal rule. A static dative “in his hand” is a defensible alternative construal.
% mdparagraph:P293:end

% teaching: NEW-U13-GRAM-U13-COOK-LOCAL-TAKE-IN-HAND
\par\begin{minipage}{\linewidth}\raggedright
% mdparagraph:P294:begin
\subsection*{Cook VII.23.10}
% mdparagraph:P294:end

% item: EX-VII-23-10
% mdparagraph:P295:begin
Sell (sęllan) me five hatchets.
% mdparagraph:P295:end

\end{minipage}\par\addvspace{4pt}

\par\begin{minipage}{\linewidth}\raggedright
% mdparagraph:P296:begin
\subsection*{Cook VII.23.11}
% mdparagraph:P296:end

% item: EX-VII-23-11
% mdparagraph:P297:begin
His bodily infirmity was very grievous.
% mdparagraph:P297:end

\end{minipage}\par\addvspace{4pt}

\par\begin{minipage}{\linewidth}\raggedright
% mdparagraph:P298:begin
\subsection*{Cook VII.23.12}
% mdparagraph:P298:end

% item: EX-VII-23-12
% mdparagraph:P299:begin
Your companionship is very pleasant.
% mdparagraph:P299:end

\end{minipage}\par\addvspace{4pt}

\par\begin{minipage}{\linewidth}\raggedright
% mdparagraph:P300:begin
\subsection*{Cook VII.23.13}
% mdparagraph:P300:end

% item: EX-VII-23-13
% mdparagraph:P301:begin
Is it any wonder that [though] the faith of God grew (196. e)?
% mdparagraph:P301:end

\end{minipage}\par\addvspace{4pt}

% teaching: U13-COOK-HELP-WONDER
% support: U13-COOK-HELP-WONDER
% mdparagraph:P302:begin
\textbf{Supplied help.} In the local wonder question, hit anticipates the following circumstance. Wundor is a nominal predicate; ðēah introduces the conceded circumstance with a preterite subjunctive. An admitted fact is compatible with this subjunctive.
% mdparagraph:P302:end

% teaching: NEW-U13-GRAM-U13-COOK-LOCAL-WONDER
\reserve{7\baselineskip}
% mdparagraph:P303:begin
\subsection*{Cook VII.24}
% mdparagraph:P303:end

% mdparagraph:P304:begin
Items 1–12. The relative controls and the supports beside the longer sentences remain part of the task.
% mdparagraph:P304:end

\par\begin{minipage}{\linewidth}\raggedright
% mdparagraph:P305:begin
\subsection*{Cook VII.24.1}
% mdparagraph:P305:end

% item: EX-VII-24-01
% mdparagraph:P306:begin
He that (sē ðe) sitteth on the throne shall dwell (wunian) in the midst of them.
% mdparagraph:P306:end

\end{minipage}\par\addvspace{4pt}

% teaching: U13-COOK-HELP-RELATIVE-LONG
% support: U13-COOK-HELP-RELATIVE-LONG
% mdparagraph:P307:begin
\textbf{Supplied help.} Sē is the retained long source spelling of the nominative masculine singular demonstrative/\allowbreak{}relative. The ordinary short se remains a valid taught alternative; the added quantity is not an independent recall target.
% mdparagraph:P307:end

% teaching: NEW-U13-GRAM-U13-COOK-LOCAL-RELATIVE-LONG
\par\begin{minipage}{\linewidth}\raggedright
% mdparagraph:P308:begin
\subsection*{Cook VII.24.2}
% mdparagraph:P308:end

% item: EX-VII-24-02
% mdparagraph:P309:begin
Angels are spirits without body.
% mdparagraph:P309:end

\end{minipage}\par\addvspace{4pt}

\par\begin{minipage}{\linewidth}\raggedright
% mdparagraph:P310:begin
\subsection*{Cook VII.24.3}
% mdparagraph:P310:end

% item: EX-VII-24-03
% mdparagraph:P311:begin
I saw a star (169) fall from heaven unto [until] the earth.
% mdparagraph:P311:end

\end{minipage}\par\addvspace{4pt}

% teaching: U13-COOK-HELP-FALL
% support: U13-COOK-HELP-FALL
% mdparagraph:P312:begin
\textbf{Supplied help.} Feallan is the supplied plain infinitive “fall.” Use the already taught gesēon + accusative actor + plain infinitive construction.
% mdparagraph:P312:end

% teaching: NEW-U13-GRAM-U13-COOK-LOCAL-FALL
\par\begin{minipage}{\linewidth}\raggedright
% mdparagraph:P313:begin
\subsection*{Cook VII.24.4}
% mdparagraph:P313:end

% item: EX-VII-24-04
% mdparagraph:P314:begin
He that overcometh, the same shall be clothed (gescrȳdan) in white (hwīt) raiment.
% mdparagraph:P314:end

% support: U13-COOK-HELP-RELATIVE-LONG
% mdparagraph:P315:begin
\textbf{Supplied help.} Sē is the retained long source spelling of the nominative masculine singular demonstrative/\allowbreak{}relative. The ordinary short se remains a valid taught alternative; the added quantity is not an independent recall target.
% mdparagraph:P315:end

\end{minipage}\par\addvspace{4pt}

\par\begin{minipage}{\linewidth}\raggedright
% mdparagraph:P316:begin
\subsection*{Cook VII.24.5}
% mdparagraph:P316:end

% item: EX-VII-24-05
% mdparagraph:P317:begin
These things saith he that is holy.
% mdparagraph:P317:end

% support: U13-COOK-HELP-RELATIVE-LONG
% mdparagraph:P318:begin
\textbf{Supplied help.} Sē is the retained long source spelling of the nominative masculine singular demonstrative/\allowbreak{}relative. The ordinary short se remains a valid taught alternative; the added quantity is not an independent recall target.
% mdparagraph:P318:end

\end{minipage}\par\addvspace{4pt}

\par\begin{minipage}{\linewidth}\raggedright
% mdparagraph:P319:begin
\subsection*{Cook VII.24.6}
% mdparagraph:P319:end

% item: EX-VII-24-06
% mdparagraph:P320:begin
Fear him (ðone) that is able to destroy (fordōn) both soul and body in (on, 166) hell.
% mdparagraph:P320:end

\end{minipage}\par\addvspace{4pt}

\par\begin{minipage}{\linewidth}\raggedright
% mdparagraph:P321:begin
\subsection*{Cook VII.24.7}
% mdparagraph:P321:end

% item: EX-VII-24-07
% mdparagraph:P322:begin
I am not worthy (wierðe) that thou enter under (172) my roof.
% mdparagraph:P322:end

\end{minipage}\par\addvspace{4pt}

% teaching: U13-COOK-HELP-WORTHY
% support: U13-COOK-HELP-WORTHY
% mdparagraph:P323:begin
\textbf{Supplied help.} Uncontracted ne eom is available for “I am not.” The bounded worthy-content construction uses a present subjunctive; under in the entry phrase takes an accusative roof. No negative-copula contraction is required.
% mdparagraph:P323:end

% teaching: NEW-U13-GRAM-U13-COOK-LOCAL-WORTHY
\par\begin{minipage}{\linewidth}\raggedright
% mdparagraph:P324:begin
\subsection*{Cook VII.24.8}
% mdparagraph:P324:end

% item: EX-VII-24-08
% mdparagraph:P325:begin
Great is your reward in heaven (plur.).
% mdparagraph:P325:end

\end{minipage}\par\addvspace{4pt}

\par\begin{minipage}{\linewidth}\raggedright
% mdparagraph:P326:begin
\subsection*{Cook VII.24.9}
% mdparagraph:P326:end

% item: EX-VII-24-09
% mdparagraph:P327:begin
If God so clothes (om. the) grass of the field, which (87. b) to-day is, and is (bið) to-morrow sent (āsęndan) into (on, 172) the oven, much (178) more he (om. shall) clothe you, O ye of little faith.
% mdparagraph:P327:end

\end{minipage}\par\addvspace{4pt}

% teaching: U13-COOK-HELP-MUCH-MORE
% support: U13-COOK-HELP-MUCH-MORE
% mdparagraph:P328:begin
\textbf{Supplied help.} Supply micle, instrumental neuter singular used adverbially before mā, for “much more.” Recognition of that cell elsewhere does not establish independent productive recall here.
% mdparagraph:P328:end

% teaching: NEW-U13-GRAM-U13-COOK-LOCAL-MUCH-MORE
% teaching: U13-COOK-HELP-TOMORROW
% support: U13-COOK-HELP-TOMORROW
% mdparagraph:P329:begin
\textbf{Supplied help.} Tō-morgen means “tomorrow” and is supplied as one invariant temporal expression.
% mdparagraph:P329:end

% teaching: NEW-U13-VOC-TO-MORGEN
% teaching: NEW-U13-PAR-TO-MORGEN
% teaching: NEW-U13-GRAM-U13-COOK-LOCAL-TOMORROW
% teaching: U13-COOK-HELP-LITTLE-FAITH
% support: U13-COOK-HELP-LITTLE-FAITH
% mdparagraph:P330:begin
\textbf{Supplied help.} The address “you of little faith” is locally supplied as gē lȳtles gelēafan. Lȳtles gelēafan is the bounded genitive quality expression; no general productive rule is inferred.
% mdparagraph:P330:end

% teaching: NEW-U13-GRAM-U13-COOK-LOCAL-LITTLE-FAITH
\par\begin{minipage}{\linewidth}\raggedright
% mdparagraph:P331:begin
\subsection*{Cook VII.24.10}
% mdparagraph:P331:end

% item: EX-VII-24-10
% mdparagraph:P332:begin
Better is the soul than food (męte), and the body than its raiment.
% mdparagraph:P332:end

\end{minipage}\par\addvspace{4pt}

\par\begin{minipage}{\linewidth}\raggedright
% mdparagraph:P333:begin
\subsection*{Cook VII.24.11}
% mdparagraph:P333:end

% item: EX-VII-24-11
% mdparagraph:P334:begin
Behold (behealdan) the birds of heaven; your heavenly Father feedeth (fēdan) them.
% mdparagraph:P334:end

\end{minipage}\par\addvspace{4pt}

\par\begin{minipage}{\linewidth}\raggedright
% mdparagraph:P335:begin
\subsection*{Cook VII.24.12}
% mdparagraph:P335:end

% item: EX-VII-24-12
% mdparagraph:P336:begin
What shall (sculon) we eat, or what drink?
% mdparagraph:P336:end

\end{minipage}\par\addvspace{4pt}

\reserve{7\baselineskip}
% mdparagraph:P337:begin
\subsection*{Cook VII.25}
% mdparagraph:P337:end

% mdparagraph:P338:begin
Items 1–12. In 25.7, “verb last” applies to the relative translating “that which I command you”: place its verb bebēode last in that relative. The surrounding conditional clause need not be forced into a universal dependent-clause pattern.
% mdparagraph:P338:end

\par\begin{minipage}{\linewidth}\raggedright
% mdparagraph:P339:begin
\subsection*{Cook VII.25.1}
% mdparagraph:P339:end

% item: EX-VII-25-01
% mdparagraph:P340:begin
He that (sē ðe) loveth not his brother, whom (87. b) he seeth, how can (mugan) he love God, whom he seeth not?
% mdparagraph:P340:end

% support: U13-COOK-HELP-RELATIVE-LONG
% mdparagraph:P341:begin
\textbf{Supplied help.} Sē is the retained long source spelling of the nominative masculine singular demonstrative/\allowbreak{}relative. The ordinary short se remains a valid taught alternative; the added quantity is not an independent recall target.
% mdparagraph:P341:end

\end{minipage}\par\addvspace{4pt}

\par\begin{minipage}{\linewidth}\raggedright
% mdparagraph:P342:begin
\subsection*{Cook VII.25.2}
% mdparagraph:P342:end

% item: EX-VII-25-02
% mdparagraph:P343:begin
The time of my departure is nigh at hand.
% mdparagraph:P343:end

\end{minipage}\par\addvspace{4pt}

\par\begin{minipage}{\linewidth}\raggedright
% mdparagraph:P344:begin
\subsection*{Cook VII.25.3}
% mdparagraph:P344:end

% item: EX-VII-25-03
% mdparagraph:P345:begin
He did something (hwæthwega) of this sort, while his companions were going to town (burg).
% mdparagraph:P345:end

\end{minipage}\par\addvspace{4pt}

% teaching: U13-COOK-HELP-INDEFINITE
% support: U13-COOK-HELP-INDEFINITE
% mdparagraph:P346:begin
\textbf{Supplied help.} Hwæthwega means “something” and takes the supplied partitive expression swilces, “of this kind.” The reading’s different exact hwæthwugu spelling remains separately recorded.
% mdparagraph:P346:end

% teaching: NEW-U13-VOC-HWAETHWEGA
% teaching: NEW-U13-PAR-HWAETHWEGA
% teaching: NEW-U13-GRAM-U13-COOK-LOCAL-INDEFINITE
\par\begin{minipage}{\linewidth}\raggedright
% mdparagraph:P347:begin
\subsection*{Cook VII.25.4}
% mdparagraph:P347:end

% item: EX-VII-25-04
% mdparagraph:P348:begin
Be with the brethren.
% mdparagraph:P348:end

\end{minipage}\par\addvspace{4pt}

\par\begin{minipage}{\linewidth}\raggedright
% mdparagraph:P349:begin
\subsection*{Cook VII.25.5}
% mdparagraph:P349:end

% item: EX-VII-25-05
% mdparagraph:P350:begin
Open the windows of the house.
% mdparagraph:P350:end

\end{minipage}\par\addvspace{4pt}

\par\begin{minipage}{\linewidth}\raggedright
% mdparagraph:P351:begin
\subsection*{Cook VII.25.6}
% mdparagraph:P351:end

% item: EX-VII-25-06
% mdparagraph:P352:begin
Love (plur.) your enemies, do (om. to) them well that hate you, that ye (om. may) be (196. f) children of your Father, which (87. b) is in heaven (plur.).
% mdparagraph:P352:end

% support: U13-COOK-HELP-RELATIVE-LONG
% mdparagraph:P353:begin
\textbf{Supplied help.} Sē is the retained long source spelling of the nominative masculine singular demonstrative/\allowbreak{}relative. The ordinary short se remains a valid taught alternative; the added quantity is not an independent recall target.
% mdparagraph:P353:end

\end{minipage}\par\addvspace{4pt}

\par\begin{minipage}{\linewidth}\raggedright
% mdparagraph:P354:begin
\subsection*{Cook VII.25.7}
% mdparagraph:P354:end

% item: EX-VII-25-07
% mdparagraph:P355:begin
Ye are my friends, if ye do (ind.) the things which (87. d) I command (bebēodan) you (verb last).
% mdparagraph:P355:end

\end{minipage}\par\addvspace{4pt}

\par\begin{minipage}{\linewidth}\raggedright
% mdparagraph:P356:begin
\subsection*{Cook VII.25.8}
% mdparagraph:P356:end

% item: EX-VII-25-08
% mdparagraph:P357:begin
The birds sang the sweetest song, winsome (acc.) to hear.
% mdparagraph:P357:end

\end{minipage}\par\addvspace{4pt}

\par\begin{minipage}{\linewidth}\raggedright
% mdparagraph:P358:begin
\subsection*{Cook VII.25.9}
% mdparagraph:P358:end

% item: EX-VII-25-09
% mdparagraph:P359:begin
The roof of the church is very high.
% mdparagraph:P359:end

\end{minipage}\par\addvspace{4pt}

\par\begin{minipage}{\linewidth}\raggedright
% mdparagraph:P360:begin
\subsection*{Cook VII.25.10}
% mdparagraph:P360:end

% item: EX-VII-25-10
% mdparagraph:P361:begin
There (ðā) was suddenly with the angel (om. a) multitude [number] of the heavenly host, praising (gen. plur.; cf. 189. 2) God, and thus saying.
% mdparagraph:P361:end

\end{minipage}\par\addvspace{4pt}

% teaching: U13-COOK-HELP-HEAVENLY-HOST
% support: U13-COOK-HELP-HEAVENLY-HOST
% mdparagraph:P362:begin
\textbf{Supplied help.} Unordered local material: cweðendra, “of people saying”; getæl, “number/\allowbreak{}company”; hęriendra, “of people praising”; heofonlices weorodes, “of a heavenly host”; sęmninga, “suddenly.” The genitive plural participles are substantival and supported in this collective-number formula. Identify their functions with this help; do not independently derive a participle paradigm.
% mdparagraph:P362:end

% teaching: NEW-U13-VOC-GETAEL
% teaching: NEW-U13-VOC-WEOROD
% teaching: NEW-U13-VOC-SEMNINGA
% teaching: NEW-U13-VOC-HERIAN
% teaching: NEW-U13-VOC-CWETHAN
% teaching: NEW-U13-PAR-SEMNINGA
% teaching: NEW-U13-PAR-HERIAN
% teaching: NEW-U13-PAR-CWETHAN
% teaching: NEW-U13-GRAM-u13-bounded-formula-support
% teaching: NEW-U13-GRAM-U13-COOK-LOCAL-HEAVENLY-HOST
\par\begin{minipage}{\linewidth}\raggedright
% mdparagraph:P363:begin
\subsection*{Cook VII.25.11}
% mdparagraph:P363:end

% item: EX-VII-25-11
% mdparagraph:P364:begin
To God be (193) glory (wuldor) in (om. the) heights, and on earth peace (om. to) men of good will.
% mdparagraph:P364:end

\end{minipage}\par\addvspace{4pt}

% teaching: U13-COOK-HELP-DOXOLOGY
% support: U13-COOK-HELP-DOXOLOGY
% mdparagraph:P365:begin
\textbf{Supplied help.} Unordered local material: Gode, dative recipient; gōdes willan, the bounded quality expression “of good will”; hēanessum, dative plural “heights”; sibb, nominative “peace”; sīe, optative “be”; wuldor, nominative “glory.” Supply the complete doxological construction: dative recipient with nominative glory, and an understood parallel wish of peace for the men described. This formula does not establish unrestricted genitive-quality syntax.
% mdparagraph:P365:end

% teaching: NEW-U13-VOC-HEANES
% teaching: NEW-U13-VOC-WULDOR
% teaching: NEW-U13-VOC-SIBB
% teaching: NEW-U13-VOC-WILLA
% teaching: NEW-U13-GRAM-U13-COOK-LOCAL-DOXOLOGY
\par\begin{minipage}{\linewidth}\raggedright
% mdparagraph:P366:begin
\subsection*{Cook VII.25.12}
% mdparagraph:P366:end

% item: EX-VII-25-12
% mdparagraph:P367:begin
He returned (ęfthweorfan) (om. by) the same way that he had come (see ǣr, adv.).
% mdparagraph:P367:end

\end{minipage}\par\addvspace{4pt}

\reserve{7\baselineskip}
% mdparagraph:P368:begin
\section*{11. The bishop and the monk}
% mdparagraph:P368:end

% mdparagraph:P369:begin
This connected account is new teaching material. It prepares the clause patterns and references of the source reading without claiming to retell the historical episode.
% mdparagraph:P369:end

% passage: NEW-U13-READING
% mdparagraph:P370:begin
\par\begin{minipage}{\linewidth}\raggedright
\readingtext{Se bisceop sæt on ðǣm hūse, and se munuc wæs ðǣr. Ðā se munuc ðone sang hīerde, ðā ēode hē tō ðǣm bisceope. Se munuc cwæð: “Ic hīerde ðone sang, ðe swīðe wynsum wæs. Hīerdest ðū hine?” Se bisceop cwæð: “Gā nū, and lǣd ðā brōðor hider.” Ðā ēode se munuc tō ðǣre cirican, and cōm ęft mid ðǣm brōðrum. Se bisceop geseah hīe and cwæð: “Gē sindon mīne frīend. Ic wille mid ēow sprecan.” Se munuc sæt mid ðǣm brōðrum, and hīe hīerdon ðæs bisceopes word.}
\end{minipage}\par
% mdparagraph:P370:end

\par\begin{minipage}{\linewidth}\raggedright
% mdparagraph:P371:begin
\subsection*{R01}
% mdparagraph:P371:end

% item: NEW-U13-R01
% mdparagraph:P372:begin
Translate the complete new account, The bishop and the monk. Preserve the inquiry, command and final speech.
% mdparagraph:P372:end

\end{minipage}\par\addvspace{4pt}

\par\begin{minipage}{\linewidth}\raggedright
% mdparagraph:P373:begin
\subsection*{R02}
% mdparagraph:P373:end

% item: NEW-U13-R02
% mdparagraph:P374:begin
Identify the referents of hē after ēode, hine in the monk’s question, hīe after geseah, and Ic and ēow in the bishop’s final speech. Explain why hine need not refer to a man.
% mdparagraph:P374:end

\end{minipage}\par\addvspace{4pt}

\par\begin{minipage}{\linewidth}\raggedright
% mdparagraph:P375:begin
\subsection*{R03}
% mdparagraph:P375:end

% item: NEW-U13-R03
% mdparagraph:P376:begin
Find the temporal clause and its main clause, the relative clause, and a coordinated finite verb with an unexpressed repeated subject. Give each finite verb and recover each subject.
% mdparagraph:P376:end

\end{minipage}\par\addvspace{4pt}

\par\begin{minipage}{\linewidth}\raggedright
% mdparagraph:P377:begin
\subsection*{R04}
% mdparagraph:P377:end

% item: NEW-U13-R04
% mdparagraph:P378:begin
Give the narrative order of hearing the song, the monk’s inquiry, the bishop’s command, the monk’s return and the bishop’s final speech. Who is present at the end, and which words establish this?
% mdparagraph:P378:end

\end{minipage}\par\addvspace{4pt}

\reserve{7\baselineskip}
% mdparagraph:P379:begin
\section*{12. The Passing of Chad: the music and summons}
% mdparagraph:P379:end

% teaching: NEW-U13-SOURCE-READING
\reserve{7\baselineskip}
% mdparagraph:P380:begin
\subsection*{English context}
% mdparagraph:P380:end

% mdparagraph:P381:begin
Owini had left a court household to live and work at Chad’s monastery. While the other brothers studied inside, he did needed work outside. Chad was alone in the local oratory; the companions had gone to the community church.
% mdparagraph:P381:end

% mdparagraph:P382:begin
This orientation is new teaching material based on the preceding passage. The two buildings must be distinguished; see original note 152.5. The extract ends with Chad’s summons to the brothers and Owini.
% mdparagraph:P382:end

\reserve{7\baselineskip}
% mdparagraph:P383:begin
\subsection*{Old English text}
% mdparagraph:P383:end

% mdparagraph:P384:begin
Cook, printed pages 152–153, editorial lines 152.3–153.9, ending partway through the last line. The complete original notes and their attachment points follow in §13. Numbers in the local helps below identify explanations, not extra words in the source.
% mdparagraph:P384:end

\par\begin{minipage}{\linewidth}\raggedright
% passage: FB-VII-CHAD-MUSIC-SUMMONS
% mdparagraph:P385:begin
\readingtext{Ðā hē ðā sume dæge hwæthwugu swilces ūte dyde, and his gefēran tō byrig tō ciricean ēodon, swā hīe gelōmlice dydon, and se bisceop, āna in ðǣre ciricean, oððe in bēcrǣdinge oððe in gebedum geornfull wæs, ðā gehīerde hē sęmninga, swā-swā hē ęft æfter ðon sægde, ðā swētestan stefne and ðā fægerestan, singendra and blissiendra, of heofonum oð eorðan āstīgan. Ðā stefne and ðone sang hē cwæð ðæt hē ǣrest gehīerde fram ēastsūðdǣle heofones, ðæt is fram hēanesse ðǣre winterlican sunnan ūpganges; and ðanan tō him styccemǣlum nēalǣcton, oð-ðæt hē becōm tō ðæcean ðǣre ciricean ðe se bisceop in wæs; and, ingangende, eall gefylde, and in ymbhwyrfte ymbsealde. And hē ðā geornlice his mōd āðęnede in ðā ðing ðe hē gehīerde. Ðā gehīerde hē ęft swā-swā healfre tīde fæce, of hrōfe ðǣre ilcan ciricean ūpāstīgan ðone ilcan blissesang, and, ðȳ ilcan wege ðe hē ǣr cōm, ūp oð heofonas mid unasęcgendre swētnesse ęfthweorfan. Ðā wunode hē ðǣr sum fæc tīde, wundriende and wafiende; and mid behygdigum mōde ðōhte and smēade hwæt ðā ðing bēon sceolden. Ðā ontȳnde se bisceop ðæt ēagðȳrel ðǣre ciricean, and mid his handa slōg tācen, swā-swā his gewuna wæs gif hwilc mann ūte wǣre, ðæt hē in tō him ēode. Ðā ēode hē sōna in tō him. Cwæð hē, se bisceop, him tō: “Gang hraðe tō ciricean, and hāt ūre seofon brōðor hider tō mē cuman; and ðū ēac swilce mid wes.”}
% mdparagraph:P385:end

\end{minipage}\par
\reserve{7\baselineskip}
% mdparagraph:P386:begin
\subsection*{Local reading help}
% mdparagraph:P386:end

% teaching: NEW-U13-GRAM-u13-source-local-idioms
% mdparagraph:P387:begin
All eighteen helps remain available for translation and analysis. They supply exact source forms and bounded constructions; they do not add full paradigms for every word in the extract.
% mdparagraph:P387:end

% teaching: U13-SRC-HELP-01
% support: U13-SRC-HELP-01
% mdparagraph:P388:begin
\textbf{Reading help 1 — sume dæge; hwæthwugu swilces; ūte.} Sume dæge is a time expression, “one day”; sume is the instrumental form. Hwæthwugu means “something”; swilces is neuter genitive, “of this kind.” Ūte means “outside.”
% mdparagraph:P388:end

% teaching: NEW-U13-VOC-HWAETHWUGU
% teaching: NEW-U13-VOC-UTE
% teaching: NEW-U13-PAR-HWAETHWUGU
% teaching: NEW-U13-PAR-UTE
% teaching: U13-SRC-HELP-02
% support: U13-SRC-HELP-02
% mdparagraph:P389:begin
\textbf{Reading help 2 — ciricean.} Ciricean is the source genitive or dative spelling corresponding to productive cirican: genitive when dependent on roof or window, dative after tō or locative in. The same noun can name either the community church or Chad’s oratory; context and Cook’s note distinguish the buildings.
% mdparagraph:P389:end

% teaching: U13-SRC-HELP-03
% support: U13-SRC-HELP-03
% mdparagraph:P390:begin
\textbf{Reading help 3 — gelōmlice; āna; in bēcrǣdinge /\allowbreak{} in gebedum geornfull.} Gelōmlice means “often”; āna, “alone.” Bēcrǣdinge means “reading books”; gebedum is dative plural, “prayers”; geornfull here means “occupied” or “intent.” In governs the dative in these expressions.
% mdparagraph:P390:end

% teaching: NEW-U13-VOC-BECRAEDING
% teaching: NEW-U13-VOC-GEBED
% teaching: NEW-U13-VOC-GEORNFULL
% teaching: NEW-U13-VOC-GELOMLICE
% teaching: NEW-U13-VOC-IN
% teaching: NEW-U13-VOC-AN
% teaching: NEW-U13-PAR-GELOMLICE
% teaching: NEW-U13-PAR-IN
% teaching: NEW-U13-PAR-AN
% teaching: U13-SRC-HELP-04
% support: U13-SRC-HELP-04
% mdparagraph:P391:begin
\textbf{Reading help 4 — gehīerde; sęmninga; swā-swā; ęft æfter ðon.} Gehīerde is prefixed “heard.” Sęmninga means “suddenly”; swā-swā, “just as”; ęft, “again” or “afterwards” as the context requires. Æfter ðon means “after that”; ðon is the source dative demonstrative variant.
% mdparagraph:P391:end

% teaching: NEW-U13-VOC-SWA-SWA
% teaching: NEW-U13-PAR-SWA-SWA
% teaching: U13-SRC-HELP-05
% support: U13-SRC-HELP-05
% mdparagraph:P392:begin
\textbf{Reading help 5 — swētestan /\allowbreak{} fægerestan stefne; singendra and blissiendra; oð; āstīgan.} Stefne is accusative feminine singular, “voice.” Fægerestan means “most beautiful,” agreeing with stefne. Singendra and blissiendra are genitive plural participles used substantivally, “of people singing and rejoicing.” Oð means “as far as.” Āstīgan here means “descend.”
% mdparagraph:P392:end

% teaching: NEW-U13-VOC-STEFN
% teaching: NEW-U13-VOC-FAEGER-SUPERL
% teaching: NEW-U13-VOC-SINGAN
% teaching: NEW-U13-VOC-BLISSIAN
% teaching: NEW-U13-PAR-SINGAN
% teaching: NEW-U13-PAR-BLISSIAN
% teaching: U13-SRC-HELP-06
% support: U13-SRC-HELP-06
% mdparagraph:P393:begin
\textbf{Reading help 6 — Ðā stefne and ðone sang … gehīerde.} The initial accusative objects belong with the later gehīerde inside the reported clause. Cwæð introduces the report. Ǣrest means “first.” Do not turn the fronted accusatives into subjects.
% mdparagraph:P393:end

% teaching: U13-SRC-HELP-07
% support: U13-SRC-HELP-07
% mdparagraph:P394:begin
\textbf{Reading help 7 — ēastsūðdǣle; hēanesse; winterlican sunnan ūpganges; ðanan.} Ēastsūðdǣle means “southeastern part”; hēanesse, “height”; winterlican, “winter,” agrees with genitive sunnan; ūpganges means “rising.” Ðanan means “from there.”
% mdparagraph:P394:end

% teaching: NEW-U13-VOC-WINTERLIC
% teaching: NEW-U13-VOC-UPGANG
% teaching: NEW-U13-VOC-THANAN
% teaching: NEW-U13-PAR-THANAN
% teaching: U13-SRC-HELP-08
% support: U13-SRC-HELP-08
% mdparagraph:P395:begin
\textbf{Reading help 8 — styccemǣlum; nēalǣcton; oð-ðæt hē becōm.} Styccemǣlum means “little by little.” Nēalǣcton is third plural: the voice and song approached. Oð-ðæt means “until.” Becōm means “arrived”; the following singular reference is explained by Cook’s note 9a.
% mdparagraph:P395:end

% teaching: U13-SRC-HELP-09
% support: U13-SRC-HELP-09
% mdparagraph:P396:begin
\textbf{Reading help 9 — ðæcean … ðe … in.} Ðæcean is dative singular, “roof.” The relative ðe is governed by the later in, “in which”; the preposition follows its relative.
% mdparagraph:P396:end

% teaching: NEW-U13-VOC-THAECE
% teaching: U13-SRC-HELP-10
% support: U13-SRC-HELP-10
% mdparagraph:P397:begin
\textbf{Reading help 10 — ingangende; eall; in ymbhwyrfte ymbsealde.} Ingangende is a present participle, “entering.” Eall is accusative neuter singular, “everything.” Gefylde means “filled.” In ymbhwyrfte means “all around”; ymbsealde means “surrounded.” These continue the description of the song.
% mdparagraph:P397:end

% teaching: NEW-U13-VOC-YMBHWYRFT
% teaching: NEW-U13-VOC-INGANGAN
% teaching: NEW-U13-VOC-YMBSELLAN
% teaching: NEW-U13-PAR-INGANGAN
% teaching: NEW-U13-PAR-YMBSELLAN
% teaching: U13-SRC-HELP-11
% support: U13-SRC-HELP-11
% mdparagraph:P398:begin
\textbf{Reading help 11 — geornlice; his mōd āðęnede in ….} Geornlice means “attentively.” Mōd is “mind”; āðęnede is “directed/\allowbreak{}stretched.” The whole attention expression means to direct one’s mind or attention to something; in has the accusative here.
% mdparagraph:P398:end

% teaching: NEW-U13-VOC-MOD
% teaching: NEW-U13-VOC-GEORNLICE
% teaching: NEW-U13-VOC-ATHENIAN
% teaching: NEW-U13-PAR-GEORNLICE
% teaching: NEW-U13-PAR-ATHENIAN
% teaching: U13-SRC-HELP-12
% support: U13-SRC-HELP-12
% mdparagraph:P399:begin
\textbf{Reading help 12 — ęft swā-swā healfre tīde fæce.} Here swā-swā means “about,” not “just as.” Healfre tīde is genitive, “of half an hour”; fæce is dative of the interval after which the event occurs. The hour sense of tīde is supplied separately.
% mdparagraph:P399:end

% teaching: NEW-U13-VOC-HEALF
% teaching: NEW-U13-VOC-FAEC
% teaching: NEW-U13-VOC-TID
% teaching: U13-SRC-HELP-13
% support: U13-SRC-HELP-13
% mdparagraph:P400:begin
\textbf{Reading help 13 — ūpāstīgan; blissesang; ðȳ ilcan wege; unasęcgendre swētnesse; ęfthweorfan.} Ūpāstīgan means “rise up”; blissesang, “song of joy.” Ðȳ ilcan wege is the supported route expression, “by the same way”; ǣr means “previously.” Unasęcgendre swētnesse means “with indescribable sweetness” after mid. Ęfthweorfan means “return.” The two infinitives depend on heard.
% mdparagraph:P400:end

% teaching: NEW-U13-VOC-BLISSESANG
% teaching: NEW-U13-VOC-UNASECGENDE
% teaching: NEW-U13-VOC-SWETNES
% teaching: NEW-U13-VOC-UPASTIGAN
% teaching: NEW-U13-PAR-UPASTIGAN
% teaching: U13-SRC-HELP-14
% support: U13-SRC-HELP-14
% mdparagraph:P401:begin
\textbf{Reading help 14 — wunode; sum fæc tīde; wundriende and wafiende.} Wunode means “remained”; ðǣr, “there.” Sum fæc tīde is an accusative duration expression, “for a while.” Wundriende and wafiende are uninflected participles describing the observer as wondering and amazed.
% mdparagraph:P401:end

% teaching: NEW-U13-VOC-WUNDRIAN
% teaching: NEW-U13-VOC-WAFIAN
% teaching: NEW-U13-PAR-WUNDRIAN
% teaching: NEW-U13-PAR-WAFIAN
% teaching: U13-SRC-HELP-15
% support: U13-SRC-HELP-15
% mdparagraph:P402:begin
\textbf{Reading help 15 — mid behygdigum mōde; ðōhte and smēade; hwæt … bēon sceolden.} Behygdigum mōde means “with an attentive mind.” Ðōhte and smēade mean “thought” and “pondered.” The embedded question asks what the things might be or mean; sceolden is preterite subjunctive plural, with infinitive bēon.
% mdparagraph:P402:end

% teaching: NEW-U13-VOC-BEHYGDIG
% teaching: NEW-U13-VOC-THENCAN
% teaching: NEW-U13-VOC-SMEAGAN
% teaching: NEW-U13-VOC-SCULAN
% teaching: NEW-U13-PAR-THENCAN
% teaching: NEW-U13-PAR-SMEAGAN
% teaching: NEW-U13-PAR-SCULAN
% teaching: U13-SRC-HELP-16
% support: U13-SRC-HELP-16
% mdparagraph:P403:begin
\textbf{Reading help 16 — ontȳnde; ēagðȳrel; handa; slōg tācen; swā-swā; gewuna.} Ontȳnde means “opened.” Handa is the source dative form of hand. Slōg tācen means “made a signal,” literally “struck a sign.” Swā-swā means “just as.” Gewuna means “custom” or “habit.”
% mdparagraph:P403:end

% teaching: NEW-U13-VOC-GEWUNA
% teaching: NEW-U13-VOC-ONTYNAN
% teaching: NEW-U13-VOC-SLEAN
% teaching: NEW-U13-PAR-ONTYNAN
% teaching: NEW-U13-VOC-HAND
% teaching: NEW-U13-PAR-SLEAN
% teaching: U13-SRC-HELP-17
% support: U13-SRC-HELP-17
% mdparagraph:P404:begin
\textbf{Reading help 17 — gif hwilc mann … wǣre; ðæt … ēode; in; sōna.} Hwilc mann here means “anyone.” Wǣre is preterite subjunctive in the habitual condition. The ðæt clause states the content or purpose of the signal. In is the adverb “in”; sōna means “immediately.”
% mdparagraph:P404:end

% teaching: U13-SRC-HELP-18
% support: U13-SRC-HELP-18
% mdparagraph:P405:begin
\textbf{Reading help 18 — Cwæð … him tō; Gang; hāt; ēac swilce mid wes.} The separated him … tō identifies the person addressed. Gang means “go” (singular imperative); hraðe, “quickly.” Hāt is singular “command/\allowbreak{}tell,” with an accusative subject and infinitive. Hider means “hither/\allowbreak{}here.” Ēac swilce means “also”; mid is adverbial “with them”; wes is singular “be.”
% mdparagraph:P405:end

% teaching: NEW-U13-VOC-GAN
% teaching: NEW-U13-VOC-HRATHE
% teaching: NEW-U13-PAR-GAN
% teaching: NEW-U13-PAR-HRATHE
\reserve{7\baselineskip}
% mdparagraph:P406:begin
\subsection*{Source-reading tasks}
% mdparagraph:P406:end

\par\begin{minipage}{\linewidth}\raggedright
% mdparagraph:P407:begin
\subsection*{S01}
% mdparagraph:P407:end

% item: NEW-U13-S01
% mdparagraph:P408:begin
Translate the complete music-and-summons extract from The Passing of Chad. Use all printed local helps and original notes.
% mdparagraph:P408:end

\end{minipage}\par\addvspace{4pt}

\par\begin{minipage}{\linewidth}\raggedright
% mdparagraph:P409:begin
\subsection*{S02}
% mdparagraph:P409:end

% item: NEW-U13-S02
% mdparagraph:P410:begin
Divide the opening sentence through of heofonum oð eorðan āstīgan into its major clauses. Identify the subjects of dyde, ēodon, dydon, wæs, gehīerde and the parenthetical sægde; distinguish the first two occurrences of ðā and the main-clause ðā before gehīerde.
% mdparagraph:P410:end

\end{minipage}\par\addvspace{4pt}

\par\begin{minipage}{\linewidth}\raggedright
% mdparagraph:P411:begin
\subsection*{S03}
% mdparagraph:P411:end

% item: NEW-U13-S03
% mdparagraph:P412:begin
Follow the sound through the account. Identify what hē refers to in hē becōm and in ðe hē ǣr cōm; describe its path using the printed vocabulary help.
% mdparagraph:P412:end

\end{minipage}\par\addvspace{4pt}

\par\begin{minipage}{\linewidth}\raggedright
% mdparagraph:P413:begin
\subsection*{S04}
% mdparagraph:P413:end

% item: NEW-U13-S04
% mdparagraph:P414:begin
Identify Owini’s actions and Chad’s actions from the arrival of the sound until the command to fetch the brothers. Where does se bisceop make a reference explicit?
% mdparagraph:P414:end

\end{minipage}\par\addvspace{4pt}

\par\begin{minipage}{\linewidth}\raggedright
% mdparagraph:P415:begin
\subsection*{S05}
% mdparagraph:P415:end

% item: NEW-U13-S05
% mdparagraph:P416:begin
In ðe se bisceop in wæs, give the antecedent and explain the role of in using the local help. Translate the complete noun phrase and relative clause.
% mdparagraph:P416:end

\end{minipage}\par\addvspace{4pt}

\par\begin{minipage}{\linewidth}\raggedright
% mdparagraph:P417:begin
\subsection*{S06}
% mdparagraph:P417:end

% item: NEW-U13-S06
% mdparagraph:P418:begin
In Ðā stefne and ðone sang hē cwæð ðæt hē ǣrest gehīerde, identify the two fronted objects and the verb that governs them. Use Cook’s note 8 and the local help. Show the reporting clause and its content clause without assuming that the fronted objects belong to cwæð.
% mdparagraph:P418:end

\end{minipage}\par\addvspace{4pt}

\par\begin{minipage}{\linewidth}\raggedright
% mdparagraph:P419:begin
\subsection*{S07}
% mdparagraph:P419:end

% item: NEW-U13-S07
% mdparagraph:P420:begin
In Chad’s final speech, identify the speaker, the addressee, and the referents of ūre, mē and ðū. Which brothers are to come, and is Owini included? Cite the wording and the original note distinguishing the churches.
% mdparagraph:P420:end

\end{minipage}\par\addvspace{4pt}

\par\begin{minipage}{\linewidth}\raggedright
% mdparagraph:P421:begin
\subsection*{S08}
% mdparagraph:P421:end

% item: NEW-U13-S08
% mdparagraph:P422:begin
Choose one main clause and one dependent clause in the extract. Identify their finite verbs, subjects and first constituents. Explain how their actual order relates to the chapter’s tendencies; do not rewrite the source into a compulsory pattern.
% mdparagraph:P422:end

\end{minipage}\par\addvspace{4pt}

\reserve{8\baselineskip}
% mdparagraph:P423:begin
\section*{13. Cook’s original notes and attachments}
% mdparagraph:P423:end

\begingroup\fontsize{10}{12.8}\selectfont
\setlength{\parskip}{1.5pt}
% mdparagraph:P424:begin
The following notes preserve Cook’s wording and original numbering. References such as 176 are to sections of Cook’s grammar. His questions about later English words remain source apparatus; they are not additional assessed tasks in this unit. The local helps in §12 supply what is needed to read the extract.
% mdparagraph:P424:end

% mdtable:T56:begin
\par\begingroup
\fontsize{10.5}{13.2}\selectfont
\setlength{\LTleft}{0pt}\setlength{\LTright}{\fill}
\begin{longtable}{@{}>{\raggedright\arraybackslash}p{35.74mm}>{\raggedright\arraybackslash}p{43.41mm}>{\raggedright\arraybackslash}p{39.23mm}@{}}
\toprule
\textbf{Original note} & \textbf{Attached word} & \textbf{Editorial line} \\
\midrule
\endfirsthead
\toprule
\textbf{Original note} & \textbf{Attached word} & \textbf{Editorial line} \\
\midrule
\endhead
\midrule
\multicolumn{3}{r}{\footnotesize Continued on the following page.}\\
\endfoot
\bottomrule
\endlastfoot
152.3 & sume & 152.3 \\
152.3 & dæge & 152.3 \\
152.4 & swilces & 152.3 \\
152.5 & ciricean & 152.5 \\
152.6 & stefne & 152.8 \\
152.7 & blissiendra & 152.8 \\
152.8 & stefne & 152.9 \\
152.8 & sang & 152.10 \\
152.9 & styccemǣlum & 152.12 \\
152.9a & hē & 152.13 \\
152.10 & ðæcean & 152.13 \\
152.11 & ðe & 152.13 \\
152.12 & eall & 152.14 \\
152.13 & mōd & 152.15 \\
152.13 & āðęnede & 152.16 \\
152.14 & swā-swā & 152.17 \\
152.15 & fæce & 152.17 \\
152.16 & blissesang & 152.18 \\
152.17 & wege & 152.19 \\
152.18 & unasęcgendre & 152.19 \\
152.19 & ęfthweorfan & 152.20 \\
153.1 & sum & 153.1 \\
153.1 & fæc & 153.1 \\
153.1 & tīde & 153.1 \\
153.2 & ēagðȳrel & 153.4 \\
153.3 & tō & 153.7 \\
153.4 & ciricean & 153.7 \\
153.5 & hāt & 153.8 \\
153.6 & wes & 153.9 \\
\end{longtable}
\par\endgroup

% mdtable:T56:end
\reserve{7\baselineskip}
% mdparagraph:P425:begin
\subsection*{152.3}
% mdparagraph:P425:end

% source-note: FB152-N3
% mdparagraph:P426:begin
See 176.
% mdparagraph:P426:end

\reserve{7\baselineskip}
% mdparagraph:P427:begin
\subsection*{152.4}
% mdparagraph:P427:end

% source-note: FB152-N4
% mdparagraph:P428:begin
See 154. b.
% mdparagraph:P428:end

\reserve{7\baselineskip}
% mdparagraph:P429:begin
\subsection*{152.5}
% mdparagraph:P429:end

% source-note: FB152-N5
% mdparagraph:P430:begin
Lat. oratorio loci. The monks had gone to the church. Cf. below, p. 153, l. 7.
% mdparagraph:P430:end

\reserve{7\baselineskip}
% mdparagraph:P431:begin
\subsection*{152.6}
% mdparagraph:P431:end

% source-note: FB152-N6
% mdparagraph:P432:begin
See 169. For this word see Chaucer, Knight’s Tale 1704.
% mdparagraph:P432:end

\reserve{7\baselineskip}
% mdparagraph:P433:begin
\subsection*{152.7}
% mdparagraph:P433:end

% source-note: FB152-N7
% mdparagraph:P434:begin
Lat. vocem suavissimam cantantium atque lætantium. What adjective is concealed in bliss- (see 34)?
% mdparagraph:P434:end

\reserve{7\baselineskip}
% mdparagraph:P435:begin
\subsection*{152.8}
% mdparagraph:P435:end

% source-note: FB152-N8
% mdparagraph:P436:begin
Acc. after gehīerde.
% mdparagraph:P436:end

\reserve{7\baselineskip}
% mdparagraph:P437:begin
\subsection*{152.9}
% mdparagraph:P437:end

% source-note: FB152-N9
% mdparagraph:P438:begin
See 72.
% mdparagraph:P438:end

\reserve{7\baselineskip}
% mdparagraph:P439:begin
\subsection*{152.9a}
% mdparagraph:P439:end

% source-note: FB152-N9a
% mdparagraph:P440:begin
Se sang.
% mdparagraph:P440:end

\reserve{7\baselineskip}
% mdparagraph:P441:begin
\subsection*{152.10}
% mdparagraph:P441:end

% source-note: FB152-N10
% mdparagraph:P442:begin
Lat. tectum, for which l. 17 has hrōf.
% mdparagraph:P442:end

\reserve{7\baselineskip}
% mdparagraph:P443:begin
\subsection*{152.11}
% mdparagraph:P443:end

% source-note: FB152-N11
% mdparagraph:P444:begin
Governed by in.
% mdparagraph:P444:end

\reserve{7\baselineskip}
% mdparagraph:P445:begin
\subsection*{152.12}
% mdparagraph:P445:end

% source-note: FB152-N12
% mdparagraph:P446:begin
Acc. sing.
% mdparagraph:P446:end

\reserve{7\baselineskip}
% mdparagraph:P447:begin
\subsection*{152.13}
% mdparagraph:P447:end

% source-note: FB152-N13
% mdparagraph:P448:begin
Lat. animum intenderet.
% mdparagraph:P448:end

\reserve{7\baselineskip}
% mdparagraph:P449:begin
\subsection*{152.14}
% mdparagraph:P449:end

% source-note: FB152-N14
% mdparagraph:P450:begin
Lat. quasi.
% mdparagraph:P450:end

\reserve{7\baselineskip}
% mdparagraph:P451:begin
\subsection*{152.15}
% mdparagraph:P451:end

% source-note: FB152-N15
% mdparagraph:P452:begin
See 176.
% mdparagraph:P452:end

\reserve{7\baselineskip}
% mdparagraph:P453:begin
\subsection*{152.16}
% mdparagraph:P453:end

% source-note: FB152-N16
% mdparagraph:P454:begin
Lat. lætitiæ canticum. See 169.
% mdparagraph:P454:end

\reserve{7\baselineskip}
% mdparagraph:P455:begin
\subsection*{152.17}
% mdparagraph:P455:end

% source-note: FB152-N17
% mdparagraph:P456:begin
See 176. 2.
% mdparagraph:P456:end

\reserve{7\baselineskip}
% mdparagraph:P457:begin
\subsection*{152.18}
% mdparagraph:P457:end

% source-note: FB152-N18
% mdparagraph:P458:begin
Lat. ineffabili.
% mdparagraph:P458:end

\reserve{7\baselineskip}
% mdparagraph:P459:begin
\subsection*{152.19}
% mdparagraph:P459:end

% source-note: FB152-N19
% mdparagraph:P460:begin
Lat. reverti.
% mdparagraph:P460:end

\reserve{7\baselineskip}
% mdparagraph:P461:begin
\subsection*{153.1}
% mdparagraph:P461:end

% source-note: FB153-N1
% mdparagraph:P462:begin
Lat. aliquantulum horæ (170).
% mdparagraph:P462:end

\reserve{7\baselineskip}
% mdparagraph:P463:begin
\subsection*{153.2}
% mdparagraph:P463:end

% source-note: FB153-N2
% mdparagraph:P464:begin
In what Mod. Eng. word is a disguised form of -ðȳrel to be found? What is the etymology of window?
% mdparagraph:P464:end

\reserve{7\baselineskip}
% mdparagraph:P465:begin
\subsection*{153.3}
% mdparagraph:P465:end

% source-note: FB153-N3
% mdparagraph:P466:begin
See 201. 1.
% mdparagraph:P466:end

\reserve{7\baselineskip}
% mdparagraph:P467:begin
\subsection*{153.4}
% mdparagraph:P467:end

% source-note: FB153-N4
% mdparagraph:P468:begin
What is the etymology of church? From what language is it originally derived?
% mdparagraph:P468:end

\reserve{7\baselineskip}
% mdparagraph:P469:begin
\subsection*{153.5}
% mdparagraph:P469:end

% source-note: FB153-N5
% mdparagraph:P470:begin
See hātan.
% mdparagraph:P470:end

\reserve{7\baselineskip}
% mdparagraph:P471:begin
\subsection*{153.6}
% mdparagraph:P471:end

% source-note: FB153-N6
% mdparagraph:P472:begin
Imper. sing. When followed by the adj. hāl, what Mod. Eng. word does it give rise to?
% mdparagraph:P472:end

\reserve{7\baselineskip}
% mdparagraph:P473:begin
\section*{14. Editorial and source notes}
% mdparagraph:P473:end

\fontsize{10}{12}\selectfont
% mdparagraph:P474:begin
\textbf{Editions and locators.} The source grammar and reading are Albert S. Cook, A First Book in Old English, second edition (1895). The exercise prompts are Cook’s Exercises in Old English, Group VII, Exercises 22–25, printed pages 11–13 (PDF pages 19–21). The Chad extract is on printed pages 152–153 of the First Book (PDF pages 168–169). Editorial line counts start at the first text line on each printed page; Cook’s marginal five-line numerals remain distinct from this counting convention.
% mdparagraph:P474:end

% mdparagraph:P475:begin
\textbf{The two order fragments.} Æfter ðon hē underðēodde is exact source wording from printed page 153, editorial line 16, outside the continuous extract; Cook’s note 13 glosses the last verb with Latin subjunxit. The dependent-clause fragment is already part of the assigned extract on page 152. Its English brackets supply the objects identified by Cook’s original note 8. The general treatment of order follows the qualified tendencies described by R. D. Fulk, An Introductory Grammar of Old English, §78. Neither source supports an exceptionless finite-second or dependent-final rule.
% mdparagraph:P475:end

% mdparagraph:P476:begin
\textbf{Forms and derivation.} The noun and adjective models follow Cook §§43–63; verbal systems follow the strong, weak and anomalous paradigms on printed pages 54–80 and the lexical entries for the individual verbs. The full kinship systems are also checked against Joseph Wright, Old English Grammar, §415. Wright §259.3 supplies the consonant simplification in swētra; the selected rule-derived genitive plural is not presented as an independently attested quotation. Forceorfan’s selected present forcierf- is a declared Cook-style normalization of dictionary cirf/\allowbreak{}cyrf variants. The movement, proximity and clothing constructions are checked against the corresponding Bosworth–Toller entries, including beforan, wið, under and gescrȳdan.
% mdparagraph:P476:end

% mdparagraph:P477:begin
\textbf{Source typography.} The continuous reading joins source line and page breaks, removes purely typographic word breaks, and represents original note cues through the attachment table. No word has been added to complete its sense. Original quantities, e/\allowbreak{}ę, ð and source variants such as ciricean, hwæthwugu and Gang are retained. Dictionary separators in joined learner headwords, including the compound nouns and adjectives, do not change their letters or vowel quantities. Wright’s thorn is represented by Cook-style eth in the derived kinship paradigms. These are declared presentation choices, not silent changes to the source reading.
% mdparagraph:P477:end

% mdparagraph:P478:begin
\textbf{C03: Cook 23.7.} The original prompt’s reference “forceorfan, 134” is retained. Section 134 of the selected First Book does not supply the forceorfan paradigm required here. Use the complete indicative table and the printed principal-part support in this unit; the original reference is not a prerequisite. No cause is assigned to the discrepancy.
% mdparagraph:P478:end

% mdparagraph:P479:begin
\textbf{Exercise controls and variants.} Cook’s directions choose targets within the language. The long sē and genitive plural ðǣra are supplied where needed; ordinary se and ðāra remain the familiar forms. The English “in white raiment” is rendered through the taught mid-with-dative clothing construction. Static dative interpretations of “before” or “in his hand” can be defensible even where a task selects a different construction. Such alternatives must be judged separately from compliance with the selected target.
% mdparagraph:P479:end

\endgroup
\end{document}
